%%%%%%%%%%%%%%%%%%%%%%%%%%%%%%%%%%%%%%%%%%%%%%%%%%%%%%%%%%%%%%%%%%%%%%%%%%%%%%%
% PREAMBOLO
%%%%%%%%%%%%%%%%%%%%%%%%%%%%%%%%%%%%%%%%%%%%%%%%%%%%%%%%%%%%%%%%%%%%%%%%%%%%%%%%
\documentclass[12pt, a4paper]{book}

% --- Lingua e Codifica ---
\usepackage[utf8]{inputenc}
\usepackage[T1]{fontenc}
\usepackage[italian]{babel}

% --- Colori (Caricati prima di TikZ per evitare conflitti) ---
\usepackage{xcolor}
\definecolor{colorA}{HTML}{003f5c} % 1 - Blu Profondo
\definecolor{colorB}{HTML}{2f4b7c} % 2 - Indaco
\definecolor{colorC}{HTML}{665191} % 3 - Viola Bluastro
\definecolor{colorD}{HTML}{a05195} % 4 - Magenta Scuro
\definecolor{colorE}{HTML}{d45087} % 5 - Rosso Porpora
\definecolor{colorF}{HTML}{f95d6a} % 6 - Arancio-Rosso
\definecolor{colorG}{HTML}{ff7c43} % 7 - Arancio Brillante

% --- Matematica e Grafica ---
\usepackage{amsfonts}
\usepackage{amsmath, amssymb, amsthm, mathtools}

\usepackage{tikz}
\usepackage{pgfplots}
\pgfplotsset{compat=1.18}
\usetikzlibrary{arrows.meta, positioning, calc, shapes.geometric, decorations.pathreplacing}

% --- Layout, Box e Font ---
\usepackage{graphicx}
\usepackage[a4paper, top=2.5cm, bottom=2.5cm, left=2.5cm, right=2.5cm]{geometry}
\usepackage{fancyhdr} 
\setlength{\headheight}{14.5pt}

\usepackage{wrapfig} 
\usepackage{caption} 
\captionsetup{labelfont=bf}

\usepackage{fontawesome5} 
\usepackage[most]{tcolorbox}
\usepackage{draftwatermark}
\usepackage{lmodern} 

% --- Watermark ---
\SetWatermarkText{\sffamily\bfseries\textcolor{gray!10}{github.com/elisabettagnello}}
\SetWatermarkScale{0.3}
\SetWatermarkAngle{45}

% --- Configurazione Stile Pagina ---
\pagestyle{fancy}
\fancyhf{} 
\fancyhead[LE,RO]{\nouppercase{\leftmark}} 
\fancyfoot[C]{\thepage} 
\renewcommand{\headrulewidth}{0.4pt} 
\renewcommand{\footrulewidth}{0.4pt} 
\renewcommand{\chaptermark}[1]{\markboth{#1}{}}

\usepackage{titlesec}
\titleformat{\chapter}[display]
  {\normalfont\huge\bfseries} 
  {}                
  {0pt}                      
  {\Huge}                     

\usepackage{hyperref} 
\hypersetup{
    colorlinks=true,
    linkcolor=colorA,
    filecolor=colorE,      
    urlcolor=colorB,
}

%%%%%%%%%%%%%%%%%%%%%%%%%%%%%%%%%%%%%%%%%%%%%%%%%%%%%%%%%%%%%%%%%%%%%%%%%%%%%%%%
% INIZIO DEL DOCUMENTO
%%%%%%%%%%%%%%%%%%%%%%%%%%%%%%%%%%%%%%%%%%%%%%%%%%%%%%%%%%%%%%%%%%%%%%%%%%%%%%%%
\begin{document}

% --- Copertina Bio-Fisica ---
\begin{titlepage}
    \thispagestyle{empty}
    
    \begin{tikzpicture}[remember picture, overlay]
        
        % 1. Sfondo con gradiente
        \shade[top color=colorA, bottom color=colorC] (current page.north west) rectangle (current page.south east);
        
        % 2. Fascia laterale a contrasto
        \fill[colorE] (current page.south west) rectangle ([xshift=0.4cm]current page.north west);
        
        % 3. Rete di Hopfield
        \begin{scope}[shift={([xshift=-1.5cm, yshift=-8cm]current page.east)}]
            \def\numnodes{24} 
            \def\radius{9} 
            
            \foreach \i in {1,...,\numnodes} {
                \foreach \j in {\i,...,\numnodes} {
                    \draw[white, opacity=0.15, line width=0.3pt] ({\i*360/\numnodes}:\radius) -- ({\j*360/\numnodes}:\radius);
                }
            }
            
            \foreach \i in {1,...,\numnodes} {
                \node[circle, fill=white, opacity=0.15, inner sep=0pt, minimum size=14pt] at ({\i*360/\numnodes}:\radius) {};
                \node[circle, fill=white, inner sep=0pt, minimum size=4pt] at ({\i*360/\numnodes}:\radius) {};
            }
        \end{scope}
    \end{tikzpicture}

    % --- Area Testi: Intestazione ---
    \begin{flushleft}
        \vspace*{1.5cm} 
        \textcolor{colorE}{\rule{2cm}{2pt}}\\[0.4cm] 
        \textcolor{white}{\sffamily\Large\textbf{SAPIENZA UNIVERSITÀ DI ROMA}}\\[0.2cm]
        \textcolor{white!70}{\sffamily\large Laurea Magistrale in Fisica}\\
        
        \vspace{4.5cm}
        
        \textcolor{colorE}{\sffamily\fontsize{50}{70}\selectfont\textbf{Physics of}}\\[0.1cm]
        \textcolor{colorE}{\sffamily\fontsize{50}{70}\selectfont\textbf{Complex Systems}}\\[0.6cm]
        \textcolor{white!80}{\sffamily\Large Appunti del corso}\\[0.5cm]
        \textcolor{white!30}{\rule{6cm}{1pt}} 
    \end{flushleft}
    
    \vfill 
    
    % --- Area Testi: Autore ---
    \begin{flushleft}
        \textcolor{white!60}{\sffamily\large A cura di}\\[0.2cm]
        \textcolor{white}{\sffamily\huge\textbf{Elisabetta Agnello}}\\[1.2cm]
        
        \textcolor{white!60}{\sffamily\large Anno Accademico}\\[0.2cm]
        \textcolor{white}{\sffamily\Large\textbf{2025/2026}}
    \end{flushleft}
    
    \vspace*{1cm} 
\end{titlepage}
% --- Fine Copertina ---

% --- Indice ---
\frontmatter 

% --- Disclaimer ---
\chapter*{Disclaimer}
\thispagestyle{empty} 

Il presente documento costituisce una raccolta di appunti personali presi durante le lezioni del corso di \textit{Physics of Complex Systems}. 

Si prega di notare che questo materiale \textbf{non è stato sottoposto a una revisione sistematica} né è stato verificato dal docente. Pertanto, è possibile che siano presenti:
\begin{itemize}
    \item Errori concettuali o imprecisioni tecniche;
    \item Refusi o errori di battitura;
    \item Parti incomplete o sezioni mancanti.
\end{itemize}

Questi appunti sono destinati esclusivamente come supporto allo studio e non sostituiscono in alcun modo il materiale didattico ufficiale. L'autore non si assume alcuna responsabilità per l'uso delle informazioni qui contenute.

% =============================================================================
\vspace{3cm} 

\begin{tcolorbox}[
    colback=colorA!5,      
    colframe=colorA,       
    title={Contatti e Segnalazioni},
    fonttitle=\bfseries\sffamily,
    arc=3mm,               
    boxrule=1.2pt,         
    halign title=flush left
]
Per consultare l'archivio aggiornato, segnalare refusi o proporre integrazioni agli appunti, puoi utilizzare i seguenti canali:
\begin{description}
    \item[\textsf{\faGlobe\ Archivio Appunti:}] \hfill \\ \href{https://elisabettagnello.github.io/elisabetta-physics-notes/}{\texttt{elisabettagnello.github.io/elisabetta-physics-notes/}}
    \item[\textsf{\faGithub\ GitHub:}] \href{https://github.com/elisabettagnello}{\texttt{github.com/elisabettagnello}}
    \item[\textsf{\faLinkedin\ LinkedIn:}] \href{https://www.linkedin.com/in/elisabetta-agnello}{\texttt{linkedin.com/in/elisabetta-agnello}}
    \item[\textsf{\faEnvelope\ Email:}] \href{mailto:agnelloe24@gmail.com}{\texttt{agnelloe24@gmail.com}}
\end{description}
\end{tcolorbox}
% =============================================================================

\clearpage 
\tableofcontents

\mainmatter

% --- INCLUSIONE DEI CAPITOLI ---

% File: lezione_01.tex
\chapter{Lezione 1}
\label{chap:lezione_01} % Etichetta per riferirsi al capitolo

\begin{flushright}
\textit{Data: 01/10/2025}
\end{flushright}


\section{Invarianza di Scala}

Uno dei concetti fondamentali nella fisica dei sistemi complessi è l'invarianza di scala. Questa proprietà descrive sistemi le cui caratteristiche statistiche non cambiano sotto una trasformazione di scala, ovvero un ingrandimento o una riduzione (zoom).

Consideriamo una trasformazione di scala sulla variabile $\epsilon$:
\begin{equation}
\epsilon \rightarrow \epsilon' = r\epsilon
\end{equation}
dove $r$ è il fattore di scala. Un sistema presenta invarianza di scala se una sua proprietà $S(\epsilon)$ si trasforma come segue:
\begin{equation}
S(\epsilon') = S(r\epsilon) = A(r)S(\epsilon)
\end{equation}
dove $A(r)$ è una funzione che dipende solo dal fattore di scala $r$ e non dalla scala originale $\epsilon$.

Si può dimostrare matematicamente che l'unica funzione che soddisfa questa relazione è la legge di potenza (power-law):
\begin{equation}
S(\epsilon) = a\epsilon^k
\end{equation}
dove $a$ e $k$ sono costanti.

Un esempio classico è il problema della misurazione della lunghezza della costa della Bretagna. A seconda della scala di misura (la lunghezza del "righello" utilizzato), la lunghezza totale misurata cambia. Questo perché a scale più piccole si possono cogliere dettagli e anfratti che vengono ignorati a scale più grandi. Se si osserva una foto di una roccia senza alcun riferimento dimensionale, è impossibile determinarne la grandezza reale. Zoomando sull'immagine, si continuano a osservare strutture statisticamente simili a quelle visibili a una scala più grande. L'invarianza di scala è il segnale matematico di questa proprietà.
\chapter{Lezione 2}
\label{chap:lezione_02} 

\begin{flushright}
\textit{Data: 02/10/2025}
\end{flushright}


\section{Funzione di Densità di Probabilità (PDF)}

Consideriamo una variabile casuale (random variable, R.V.) continua $X$ e la sua funzione di densità di probabilità (Probability Density Function, PDF) $\rho(x)$. La probabilità che la variabile $X$ assuma un valore compreso nell'intervallo infinitesimo $[x, x+dx]$ è data da:
\begin{equation}
P(x \le X < x+dx) = \rho(x)dx
\end{equation}
La PDF deve soddisfare la condizione di normalizzazione, ovvero l'integrale su tutto il dominio $D$ della variabile deve essere uguale a 1:
\begin{equation}
\int_D \rho(x)dx = 1
\end{equation}
La probabilità che $X$ cada in un intervallo finito $[a,b]$ si calcola integrando la PDF su quell'intervallo:
\begin{equation}
P(a < X < b) = \int_a^b \rho(x)dx
\end{equation}
Il valore atteso (o valore di aspettazione) di un'osservabile $O(x)$ si calcola come:
\begin{equation}
\langle O \rangle = \int_D \rho(x)O(x)dx
\end{equation}
In pratica, spesso non conosciamo la vera PDF $\rho(x)$. Quello che facciamo è acquisire dati tramite esperimenti o altre fonti. A partire da un set di $N$ misure $\{o_i\}$, possiamo calcolare uno stimatore del valore atteso, ovvero la media campionaria:
\begin{equation}
\langle O \rangle_{EST} = \frac{1}{N}\sum_{i=1}^N o_i
\end{equation}
Per la legge dei grandi numeri, nel limite in cui il numero di misure $N$ tende a infinito, lo stimatore converge al vero valore atteso:
\begin{equation}
\lim_{N \to +\infty} \langle O \rangle_{EST} \to \langle O \rangle
\end{equation}



Consideriamo ora una variabile casuale discreta $Z$, che può assumere un insieme discreto di valori. L'equivalente del valore atteso si calcola tramite una sommatoria. Se abbiamo $N$ osservazioni $o_i$ della nostra osservabile, il limite della media campionaria può essere espresso come:
\begin{equation}
\lim_{N \to \infty} \frac{1}{N} \sum_{i=1}^N o_i = \lim_{N \to \infty} \sum_{z=z_{min}}^{z_{max}} O_z f_z
\end{equation}
dove $O_z$ è il valore dell'osservabile quando la variabile casuale è $z$, e $f_z = N_z/N$ è la frequenza con cui si osserva il valore $z$ (con $N_z$ numero di occorrenze di $z$). Nel limite $N \to \infty$, la frequenza $f_z$ converge alla probabilità $p_z$ di osservare il valore $z$. Pertanto, il valore atteso diventa:
\begin{equation}
\langle O \rangle = \sum_{z=z_{min}}^{z_{max}} O_z p_z
\end{equation}

\subsubsection{Esempio di una distribuzione con media non definita:}

Non tutte le distribuzioni di probabilità hanno momenti finiti. Consideriamo la seguente distribuzione a legge di potenza per $x \ge 1$:
\begin{equation}
\rho(x) = \frac{1}{2}x^{-3/2}
\end{equation}
Calcoliamo il valore atteso di $x$:
\begin{equation}
\langle x \rangle = \int_1^\infty x \rho(x) dx = \int_1^\infty x \left(\frac{1}{2}x^{-3/2}\right) dx = \frac{1}{2} \int_1^\infty x^{-1/2} dx = \left[x^{1/2}\right]_1^\infty \to \infty
\end{equation}
Il valore atteso diverge. Questo ha una conseguenza pratica importante. Se campioniamo $N$ valori da questa distribuzione e calcoliamo la media, questo valore non convergerà. Anzi, all'aumentare di $N$, la media campionaria tenderà a crescere, poiché aumenta la probabilità di campionare valori molto grandi dalla "coda" pesante della distribuzione. Questo comportamento della media campionaria può essere un segnale empirico che la distribuzione sottostante non ha un valore atteso finito.

\subsection{Funzioni di Variabili Casuali}

Spesso siamo interessati alla distribuzione di una variabile casuale $Y$ che è funzione di un'altra variabile casuale $X$ di cui conosciamo la distribuzione, $y=f(x)$. Se la funzione $f$ è invertibile, $x=f^{-1}(y)$, la probabilità che la variabile $X$ si trovi nell'intervallo $[x, x+dx]$ deve essere uguale alla probabilità che la variabile $Y$ si trovi nell'intervallo corrispondente $[y, y+dy]$. Questo principio di conservazione della probabilità si scrive come:
\begin{equation}
\tilde{\rho}(y)|dy| = \rho(x)|dx|
\end{equation}
Da cui si ottiene la PDF per $Y$:
\begin{equation}
\tilde{\rho}(y) = \rho(x) \left|\frac{dx}{dy}\right|  
\end{equation}
Il valore assoluto è necessario perché le probabilità devono essere positive.

\subsection{Funzione Cumulativa Inferiore}

Un caso particolarmente importante e utile di funzione di variabile casuale è la funzione di ripartizione (o cumulativa inferiore, CDF):
\begin{equation}
y = P_<(x) = \int_a^x \rho(x')dx'
\end{equation}
dove la variabile $X$ è definita nell'intervallo $[a,b]$. Per costruzione, la variabile $Y$ è definita nell'intervallo $[0,1]$, poiché $P_<(a)=0$ e $P_<(b)=1$.

Per trovare la distribuzione di $Y$, calcoliamo la derivata di $y$ rispetto a $x$:
\begin{equation}
\frac{dy}{dx} = \rho(x)
\end{equation}
Usando la formula di trasformazione, otteniamo la PDF di $Y$:
\begin{equation}
\tilde{\rho}(y) = \rho(x) \left|\frac{dx}{dy}\right| = \rho(x) \frac{1}{\rho(x)} = 1
\end{equation}
Questo significa che la variabile $Y$ è distribuita uniformemente nell'intervallo $[0,1]$.

Questo risultato è estremamente potente e fornisce un metodo generale per generare numeri casuali secondo una qualsiasi distribuzione $\rho(x)$ desiderata:
\begin{enumerate}
    \item Generare un numero casuale $y$ da una distribuzione uniforme in $[0,1]$.
    \item Calcolare la funzione di ripartizione $P_<(x)$ per la distribuzione $\rho(x)$.
    \item Invertire la relazione $y = P_<(x)$ per ottenere $x = P_<^{-1}(y)$.
    \item Il valore di $x$ così ottenuto sarà un numero casuale distribuito secondo $\rho(x)$.
\end{enumerate}

\textbf{Esempio:} Generare numeri casuali da $\rho(x) = cx^{-2}$ per $x \in [a,b]$.
Prima normalizziamo la distribuzione:
\begin{equation}
\int_a^b c x^{-2} dx = c \left[-\frac{1}{x}\right]_a^b = c\left(\frac{1}{a} - \frac{1}{b}\right) = 1 \implies c = \frac{1}{\frac{1}{a}-\frac{1}{b}}
\end{equation}
Calcoliamo la CDF:
\begin{equation}
y = P_<(x) = \int_a^x c(x')^{-2} dx' = c\left(-\frac{1}{x} + \frac{1}{a}\right)
\end{equation}
Infine, invertiamo la relazione per trovare $x(y)$:
\begin{equation}
x = \frac{c}{c/a - y}
\end{equation}
Generando valori di $y$ uniformi in $[0,1]$ e inserendoli in questa formula, si ottengono valori di $x$ distribuiti come $x^{-2}$.

\section{Somma di Variabili Casuali e Funzione Caratteristica}

Consideriamo due variabili casuali indipendenti $X_1$ e $X_2$, con PDF $\rho_1(x_1)$ e $\rho_2(x_2)$. Vogliamo trovare la distribuzione della loro somma $X = X_1 + X_2$. La PDF di $X$, $\rho(x)$, è data dalla convoluzione delle due PDF:
\begin{equation}
\rho(x) = \int_{-\infty}^{+\infty} dx_1 \rho_1(x_1) \rho_2(x-x_1)
\end{equation}
Questo integrale può essere complicato da calcolare direttamente. Un metodo più agevole passa attraverso l'uso della \textbf{funzione caratteristica}, che è definita come la trasformata di Fourier della PDF:
\begin{equation}
\hat{\rho}(z) = \int_D e^{izx} \rho(x) dx
\end{equation}
La PDF può essere recuperata tramite l'antitrasformata di Fourier:
\begin{equation}
\rho(x) = \frac{1}{2\pi} \int_{-\infty}^{+\infty} e^{-izx} \hat{\rho}(z) dz
\end{equation}
La proprietà fondamentale è che la funzione caratteristica della somma di due variabili casuali indipendenti è il prodotto delle loro funzioni caratteristiche individuali:
\begin{equation}
\hat{\rho}_X(z) = \hat{\rho}_{X_1}(z) \cdot \hat{\rho}_{X_2}(z)
\end{equation}

\textbf{\textit{Dimostrazione:}}
\begin{align*}
\hat{\rho}(z) &= \int e^{izx} \rho(x) dx = \int e^{izx} \left( \int \rho_1(x_1) \rho_2(x-x_1) dx_1 \right) dx \\
&= \iint e^{izx} \rho_1(x_1) \rho_2(x-x_1) dx_1 dx
\end{align*}
Per scambiare l'ordine di integrazione (usando il teorema di Fubini), si verifica che l'integrale del valore assoluto converga, il che è vero poiché le PDF sono normalizzate. 

Con il cambio di variabile $y=x-x_1$ (e quindi $x=y+x_1$), l'integrale diventa:
\begin{align*}
\hat{\rho}(z) &= \int \rho_1(x_1) \left( \int \rho_2(y) e^{iz(y+x_1)} dy \right) dx_1 \\
&= \int \rho_1(x_1) e^{izx_1} dx_1 \int \rho_2(y) e^{izy} dy \\
&= \hat{\rho}_1(z) \hat{\rho}_2(z)
\end{align*}

\subsection{Cumulanti e Momenti}

La funzione caratteristica è uno strumento potente per calcolare i momenti e i cumulanti di una distribuzione. I \textbf{momenti} di ordine $n$ sono definiti come $\mu_n = \langle x^n \rangle$. Si possono ottenere derivando la funzione caratteristica:
\begin{equation}
\mu_n = \langle x^n \rangle = \left[\frac{1}{i}\frac{d}{dz}\right]^n \hat{\rho}(z) \bigg|_{z=0}
\end{equation}
Ad esempio, $\mu_1 = \langle x \rangle$ e $\mu_2 = \langle x^2 \rangle$.

I \textbf{cumulanti} $C_n$ sono definiti in modo simile, ma a partire dal logaritmo della funzione caratteristica:
\begin{equation}
C_n = \left[\frac{1}{i}\frac{d}{dz}\right]^n \log\hat{\rho}(z) \bigg|_{z=0}
\end{equation}
I primi cumulanti sono legati ai momenti in modo semplice:
\begin{itemize}
    \item $C_1 = \langle x \rangle = \mu$ (la media)
    \item $C_2 = \langle x^2 \rangle - \langle x \rangle^2 = \sigma^2$ (la varianza)
    \item $C_3 = \langle (x-\mu)^3 \rangle$ (legato alla skewness, $\gamma_1 = C_3/\sigma^3$)
    \item $C_4 = \langle (x-\mu)^4 \rangle - 3\sigma^4$ (legato alla curtosi, $\gamma_2 = C_4/\sigma^4$)
\end{itemize}
La proprietà più importante dei cumulanti deriva dalla relazione del prodotto per la funzione caratteristica della somma di variabili indipendenti. Prendendo il logaritmo, si ha $\log \hat{\rho}_X(z) = \log \hat{\rho}_{X_1}(z) + \log \hat{\rho}_{X_2}(z)$. Poiché la derivata è un operatore lineare, ne consegue che \textbf{i cumulanti della somma sono la somma dei cumulanti}:
\begin{equation}
C_{n,X} = C_{n,X_1} + C_{n,X_2}
\end{equation}

\subsection{Il Teorema del Limite Centrale}

Consideriamo la somma $X_N = \sum_{n=1}^N x_n$ di $N$ variabili casuali indipendenti e identicamente distribuite (i.i.d.), ciascuna con media $\mu_0$ e varianza $\sigma_0^2$. Per la proprietà di additività dei cumulanti:
\begin{itemize}
    \item $C_{1,X_N} = \sum C_{1,x_n} = N\mu_0$
    \item $C_{2,X_N} = \sum C_{2,x_n} = N\sigma_0^2$
    \item $C_{n,X_N} = N C_{n,1}$
\end{itemize}
Ora consideriamo la media campionaria $Y = X_N/N$. Usando le proprietà di trasformazione, si può dimostrare che la funzione caratteristica di $Y$ è $\hat{\rho}_Y(z) = \hat{\rho}_{X_N}(z/N)$. Questo implica che i cumulanti scalano come $C_{n,Y} = N^{-n}C_{n,X_N} = N^{1-n}C_{n,1}$.
Nel limite $N \to \infty$:
\begin{itemize}
    \item $C_{1,Y} = N^0 C_{1,1} = \mu_0 \to \mu_0$ (Legge dei Grandi Numeri)
    \item $C_{2,Y} = N^{-1} C_{2,1} = \frac{\sigma_0^2}{N} \to 0$
    \item $C_{n,Y} = N^{1-n} C_{n,1} \to 0$ per $n \ge 2$.
\end{itemize}
Una distribuzione con varianza e tutti i cumulanti superiori nulli è una funzione delta di Dirac centrata sulla media.

Per ottenere una distribuzione non degenere, dobbiamo standardizzare la variabile in modo diverso. Definiamo la variabile standardizzata $V$:
\begin{equation}
V = \frac{X_N - N\mu_0}{\sqrt{N}\sigma_0} = \frac{Y - \mu_0}{\sigma_0/\sqrt{N}}
\end{equation}
Si può dimostrare che i cumulanti di $V$ sono:
\begin{itemize}
    \item $C_{1,V} = 0$ (la media è zero)
    \item $C_{2,V} = 1$ (la varianza è uno)
    \item $C_{n,V} \propto N^{1-n/2} \to 0$ per $n \ge 3$
\end{itemize}
L'unica distribuzione con i primi due cumulanti finiti (0 e 1) e tutti gli altri nulli è la \textbf{distribuzione Gaussiana standard}. Questo dimostra il Teorema del Limite Centrale: la somma (o media) di un gran numero di variabili i.i.d. (con media e varianza finite) tende a una distribuzione Gaussiana, indipendentemente dalla distribuzione delle singole variabili.

\section{Leggi di Potenza (Power-Laws)}

Una variabile casuale segue una legge di potenza se la sua PDF ha la forma:
\begin{equation}
\rho(x) = \begin{cases} cx^{-\alpha} & \text{per } x \ge m \\ 0 & \text{per } x < m \end{cases}
\end{equation}
dove $\alpha > 0$ è l'esponente, $m > 0$ è il valore minimo, e $c$ è la costante di normalizzazione.
La normalizzazione richiede $\int_m^\infty cx^{-\alpha} dx = 1$. L'integrale converge solo se $\alpha > 1$. In tal caso:
\begin{equation}
c = (\alpha-1)m^{\alpha-1}
\end{equation}
Una caratteristica distintiva delle leggi di potenza è che appaiono come una linea retta su un grafico in scala log-log:
\begin{equation}
\log(\rho(x)) = \log(c) - \alpha \log(x)
\end{equation}
Questa è l'equazione di una retta con pendenza $-\alpha$.

\subsection{Stima dei Parametri: Metodo della Massima Verosimiglianza}

Supponiamo di avere un insieme di dati $\{x_i\}_{i=1...N}$ e di ipotizzare che provengano da una distribuzione a legge di potenza con un dato $m$ ma $\alpha$ sconosciuto. Per stimare $\alpha$, usiamo il \textbf{metodo della massima verosimiglianza} (Maximum Likelihood Estimation, MLE). Questo metodo cerca il valore del parametro $\alpha$ che massimizza la probabilità (verosimiglianza) di osservare i dati dati.

La funzione di verosimiglianza $\mathcal{L}(\alpha)$ è il prodotto delle probabilità di ogni singola osservazione:
\begin{equation}
\mathcal{L}(\alpha) = \prod_{i=1}^N \rho(x_i|\alpha) = \prod_{i=1}^N \frac{\alpha-1}{m} \left(\frac{x_i}{m}\right)^{-\alpha} = \left(\frac{\alpha-1}{m}\right)^N \prod_{i=1}^N \left(\frac{x_i}{m}\right)^{-\alpha}
\end{equation}
È più comodo massimizzare il logaritmo della verosimiglianza (log-likelihood), $L(\alpha) = \log\mathcal{L}(\alpha)$:
\begin{equation}
L(\alpha) = N\log(\alpha-1) - N\log(m) - \alpha \sum_{i=1}^N \log\left(\frac{x_i}{m}\right)
\end{equation}
Per trovare il massimo, deriviamo rispetto ad $\alpha$ e poniamo la derivata uguale a zero:
\begin{equation}
\frac{dL(\alpha)}{d\alpha} = \frac{N}{\alpha-1} - \sum_{i=1}^N \log\left(\frac{x_i}{m}\right) = 0
\end{equation}
Risolvendo per $\alpha$, otteniamo lo stimatore di massima verosimiglianza $\hat{\alpha}$:
\begin{equation}
\hat{\alpha} = 1 + N \left[ \sum_{i=1}^N \log\left(\frac{x_i}{m}\right) \right]^{-1}
\end{equation}
Questo può essere scritto in modo più compatto come:
\begin{equation}
\hat{\alpha} = 1 + \frac{1}{\langle \ln(x/m) \rangle}
\end{equation}


\chapter{Lezione 3}
\label{chap:lezione_03} 

\begin{flushright}
\textit{Data: 08/10/2025}
\end{flushright}

\chapter{Lezione 4}
\label{chap:lezione_04} 

\begin{flushright}
\textit{Data: 09/10/2025}
\end{flushright}


\section{Estensione del Teorema Centrale del Limite per code pesanti }

L'analisi dei momenti di una distribuzione a legge di potenza ($\rho(x) \sim x^{-\alpha}$) è fondamentale per comprendere il comportamento del Teorema Centrale del Limite.

Ricordiamo le condizioni di finitudine per i momenti:
\begin{itemize}
    \item La media $\langle x \rangle$ è finita per $\alpha > 2$.
    \item La varianza $\langle x^2 \rangle$ è finita per $\alpha > 3$.
\end{itemize}

L'obiettivo è estendere l'equazione del TLC a diversi casi. Iniziamo considerando il caso in cui $\alpha > 3$, ovvero la varianza è finita.

\subsection{Calcolo del Terzo Cumulante}

Per il caso in cui $\langle x^2 \rangle$ è finito, consideriamo momenti di ordine superiore (cumulanti di ordine $m \ge 3$) e consideriamo $\alpha$ appena superiore a 3, definendo:

\begin{equation}
\alpha = 3 + \delta
\end{equation}

dove $\delta > 0$.

Per una variabile casuale gaussiana, tutti i cumulanti di ordine $m \ge 3$ sono nulli, e la distribuzione della somma di variabili IID (Indipendenti e Identicamente Distribuite) converge a una Gaussiana. Vogliamo verificare se questo accade anche in questo caso, calcolando il terzo cumulante $C_{3,1}$ per una singola variabile.

Il cumulante di ordine tre per una singola variabile è dato dall'integrale, dove $X_N$ è il valore massimo delle $N$ estrazioni:

\begin{equation}
C_{3,1}\sim\int_{m}^{x_{N}}dx(\alpha-1)m^{\alpha-1}x^{-\alpha}(x-\mu)^{3}
\end{equation}

Espandendo l'integrale (utilizzando i termini dominanti della PDF come $\rho(x) \sim x^{-\alpha}$ e $\rho(x) (x-\mu)^3 \sim x^{3-\alpha}$):

\begin{equation}
=\int_{m}^{x_{N}}dx \; (x^{3-\alpha}-x^{-\alpha}\mu^{3}-3x^{2-\alpha}\mu-3x^{1-\alpha}\mu^{2}) 
\end{equation}

Sostituendo $\alpha = 3 + \delta$:

\begin{equation}
=\int_{m}^{x_{N}}dx(x^{-\delta}-x^{-3-\delta}\mu^{3}-3x^{-1-\delta}\mu-3x^{-2-\delta}\mu^{2}) 
\end{equation}

Il termine dominante nell'integrale è $x^{-\delta}$. L'integrazione fornisce:

\begin{equation}
C_{3,1}\sim\int_{m}^{x_{N}}x^{-\delta}dx\sim x_{N}^{1-\delta}-m^{1-\delta}
\end{equation}

Il termine dominante dipende dal valore di $\delta$:
\begin{itemize}
    \item Se $\delta > 1$ (ovvero $\alpha > 4$), l'esponente $1-\delta$ è negativo. Per $X_N \to \infty$, il termine $X_N^{1-\delta}$ tende a zero, e $C_{3,1}$ è una costante. Questo caso non è molto interessante.
    \item Se $0 < \delta < 1$ (ovvero $3 < \alpha < 4$), l'esponente $1-\delta$ è positivo e il termine dominante è $X_N^{1-\delta}$. Poiché $X_N \sim N^{1/(\alpha-1)} = N^{1/(2+\delta)}$, si ha:
    \begin{equation}
    C_{3,1}\sim N^{\frac{1-\delta}{2+\delta}}
    \end{equation}
\end{itemize}

\subsection{Cumulanti della Variabile Scalata }

Definiamo la variabile scalata $V$ (standardizzata) per la somma $Y = \frac{1}{N}\sum_i x_i$:

\begin{equation}
V:v=(Y-\mu_{0})\frac{\sqrt{N}}{\sigma_{0}}
\end{equation}

Per questa variabile standardizzata, i cumulanti di una Gaussiana sono $C_{1,v}=0$ e $C_{2,v}=1$. Vogliamo verificare la tendenza a zero dei cumulanti di ordine superiore, come $C_{3,v}$.

Il cumulante $C_{m}$ della media di $N$ variabili $\bar{X}$ è legato al cumulante della singola variabile $C_{m,1}$ dalla relazione:

\begin{equation}
C_{m, \bar{X}} = \frac{C_{m,N}}{N^{m}} = \frac{1}{N^{m}} C_{m,1} N = C_{m,1} N^{1-m}
\end{equation}

Il cumulante $C_{m,v}$ per la variabile scalata $V$ (normalizzata per deviazione standard) è dato da:

\begin{equation}
C_{m, v}=\frac{C_{m, 1}}{\sigma_{0}^{m}}N^{1-m/2}
\end{equation}

Per $m=3$, il terzo cumulante scalato è:

\begin{equation}
C_{3,v}=\frac{C_{3,1}}{\sqrt{N}{\sigma_{0}}^{3}}
\end{equation}

Sostituendo la dipendenza da $N$ di $C_{3,1}$ per $3 < \alpha < 4$:

\begin{equation}
C_{3,\nu}\sim N^{\frac{1-\delta}{2+\delta}}N^{-1/2}\sim N^{-\frac{3\delta}{2(2+\delta)}}
\end{equation}

Poiché $\delta > 0$, l'esponente è negativo. Quindi, per $N \rightarrow \infty$, si ha $C_{3,v} \rightarrow 0$.

In generale, il cumulante di ordine $m$ per la singola variabile scala come:
\begin{equation}
C_{m,1}\sim N^{\frac{m-2-\delta}{2+\delta}}
\end{equation}

E il cumulante di ordine $m$ per la variabile scalata $V$ scala come:
\begin{equation}
C_{m, v}\sim N^{\frac{-m\delta}{2(2+\delta)}} \rightarrow 0 \quad \text{per } N\rightarrow\infty
\end{equation}

Nel caso in cui $\alpha > 3$ (varianza finita), tutti i cumulanti di ordine $m \ge 3$ della variabile scalata $V$ tendono a zero. La distribuzione limite è dunque la \textbf{Distribuzione Gaussiana}. Per $\alpha > 3$ , dunque, una power law segue il Teorema Centrale del Limite.

\section{Distribuzioni Stabili di Lévy per code molto pesanti }

Consideriamo ora il caso in cui $1 < \alpha < 3$, dove la varianza è divergente. Il TLC non si applica, ma la distribuzione limite della somma di variabili converge a una \textbf{Distribuzione Stabile di Lévy}.

\subsection{La Funzione Caratteristica $\hat{\rho}(z)$}

La distribuzione di probabilità $\rho(x)$ per una variabile simmetrica con legge di potenza nelle code è:
\begin{equation}
\rho(x)=\begin{cases}\frac{1}{2}(\alpha-1)m^{\alpha-1}|x|^{-\alpha} & |x|\ge m \\ 0 & |x|<m\end{cases}
\end{equation}

\begin{figure}[h]
    \centering
\begin{tikzpicture}[scale=0.8]
    \begin{axis}[
        axis lines = middle,
        xmin = -5.5, xmax = 5.5,
        ymin = 0, ymax = 5.5,
        xtick = {-2, 2},
        xticklabels = {$-m$, $m$},
        ytick = \empty,
        extra x ticks = {0},
        extra x tick labels = {0},
        thick,
        xlabel = {}, ylabel = {},
        tick label style = {font=\large, color=purple},
        axis line style = {very thick}
    ]

    % Parametro m (distanza dall'asse y)
    \def\m{2}

    % Parte sinistra del grafico (fino a -m)
    \addplot [
        domain=-5:-2, 
        samples=100, 
        color=cyan!70!blue, 
        ultra thick
    ] {exp(x+2) * 4};

    % Parte destra del grafico (da m in poi)
    \addplot [
        domain=2:5, 
        samples=100, 
        color=cyan!70!blue, 
        ultra thick
    ] {exp(-x+2) * 4};

    % Linee tratteggiate verticali rosa
    \draw [dashed, ultra thick, red!50!purple!60] (axis cs:-2,0) -- (axis cs:-2,5.5);
    \draw [dashed, ultra thick, red!50!purple!60] (axis cs:2,0) -- (axis cs:2,5.5);

    \end{axis}
\end{tikzpicture}
\caption{Distribuzione di probabilità $\rho(x)$.}
\end{figure}

La funzione caratteristica $\hat{\rho}(z)$ è calcolata come:
\begin{equation}
1-\hat{\rho}(z)=\int_{-\infty}^{-m}dx \; (1-e^{izx})\rho(x)+\int_{m}^{\infty}dx~\rho(x)(1-e^{izx})
\end{equation}
\begin{equation}
=\left [\frac{(\alpha-1)m^{\alpha-1}}{2}\right ] \left [\int_{-\infty}^{-m}(-x)^{-\alpha}(1-e^{izx})dx+\int_{m}^{\infty}dx~x^{-\alpha}(1-e^{izx})\right ] 
\end{equation}
\begin{equation}
=\left [\frac{(\alpha-1)m^{\alpha-1}}{2}\right ] \left [\int_{m}^{\infty}dx~x^{-\alpha}(1-e^{-izx})+\int_{m}^{\infty}dx~x^{-\alpha}(1-e^{izx})\right ] 
\end{equation}

Dopo un cambio di variabile $u=|z|x$, si ottiene l'approssimazione:
\begin{equation}
1-\hat{\rho}(z)=1-(\alpha-1)(m|z|)^{\alpha-1}\int_{m|z|}^{+\infty}du \; \frac{1-\cos u}{u^{\alpha}}
\end{equation}

Per $u<<1$ vale: $1-\cos u \approx \frac{u^2}{2}$ 

Segue:
\begin{equation}
\int_{m|z|}^{+\infty}du \; u^{2-\alpha} \approx |z|^{3-\alpha}
\end{equation}

per $\alpha >3$ diverge, ma stiamo considerando il caso $1 < \alpha < 3$

Per $|z|<<1$ si ha:

\begin{equation}
\hat{\rho}(z)=1-K|z|^{\alpha-1}
\end{equation}
dove $K$ è una costante definita come:
\begin{equation}
K=(\alpha-1)m^{\alpha-1}C
\end{equation}
con la costante $C$ data da:
\begin{equation}
C=\int_{0}^{\infty}du\frac{1-cosu}{u^{\alpha}}= \alpha \Gamma(\alpha)\sin(\frac{\pi \alpha}{2})
\end{equation}

Questa forma è valida per $1 < \alpha < 3$.

\subsection{Distribuzione Limite}

Per trovare la distribuzione limite della somma di $N$ variabili, si usa un fattore di riscalamento diverso rispetto al TLC, basato sul valore massimo $\max\{x_i\}$ invece che sulla deviazione standard $\sigma_0$. Definiamo la variabile riscalata $V$ come:
\begin{equation}
V:v=\frac{S}{X_{N}}=\frac{\sum_{i}x_{i}}{\max\{x_{i}\}}
\end{equation}
dove $\max\{x_i\} \sim N^{1/(\alpha-1)}$.

La funzione caratteristica della somma riscalata è:
\begin{equation}
\hat{\rho}_{v}(z)=\hat{\rho}(|z|N^{-\frac{1}{\alpha-1}})^{N}\sim(1-K|z|^{\alpha-1}N^{-\frac{\alpha-1}{\alpha-1}})^{N}
\end{equation}

\begin{equation}
\sim(1-\frac{K|z|^{\alpha-1}}{N})^{N}\xrightarrow{N\rightarrow\infty} e^{-K|z|^{\alpha-1}} \quad \quad  \quad  \text{(\textbf{Distribuzione di Lévy})}
\end{equation}

Per $\alpha=2$ (dove la media è finita ma la varianza è infinita), la distribuzione limite è la \textbf{Distribuzione di Cauchy}:
\begin{equation}
\hat{\rho}_{v}\sim e^{-K|z|}\rightarrow\rho_{v}(x)=\frac{1}{\pi}\frac{K}{K^{2}+x^{2}}
\end{equation}

\newpage
\section{Legge di Benford e Processi Moltiplicativi}

La \textbf{Legge di Benford} descrive la probabilità $P(n)$ che la prima cifra significativa di un insieme di numeri sia $n$ (con $n=1, \dots, 9$). A differenza di una distribuzione uniforme ($P(n) \approx 1/9$), la Legge di Benford prevede che le cifre più piccole (soprattutto $n=1$) siano più frequenti.

\begin{figure}[h!]
    \centering
    \begin{tikzpicture}[scale=0.7]
        \begin{axis}[
            width=12cm,
            height=8cm,
            title={First 100.000 Fibonacci Numbers},
            xlabel={First Digit},
            ylabel={Probability},
            ymin=0, ymax=0.32,
            xmin=0.5, xmax=9.5,
            xtick={1,2,3,4,5,6,7,8,9},
            ytick={0, 0.05, 0.10, 0.15, 0.20, 0.25, 0.30},
            ymajorgrids=false,
            axis line style={gray!50},
            tick label style={font=\small, /pgf/number format/fixed, /pgf/number format/precision=2},
            legend style={at={(0.95,0.8)}, anchor=north east, draw=none, fill=none, font=\small},
            bar width=1.0, % Per rendere le barre adiacenti come nell'originale
        ]

        % Istogramma (Distribuzione delle prime cifre dei numeri di Fibonacci)
        % Valori approssimati basati sulla legge di Benford (a cui i Fibonacci convergono)
        \addplot[
            ybar,
            fill=colorD!50,
            draw=white,
            area legend
        ] coordinates {
            (1, 0.301) (2, 0.176) (3, 0.125) (4, 0.097) (5, 0.079) 
            (6, 0.067) (7, 0.058) (8, 0.051) (9, 0.046)
        };
        \addlegendentry{1st digit distribution}

        % Curva di Benford (Tratteggiata)
        \addplot[
            black,
            dashed,
            very thick,
            domain=1:9,
            samples=100,
        ] {log10(1 + 1/x)};
        \addlegendentry{Benford's law}

        \end{axis}
    \end{tikzpicture}
    \caption{Distribuzione della prima cifra per i primi 100.000 numeri della successione di Fibonacci confrontata con la Legge di Benford.}
\end{figure}

\subsubsection{Derivazione della Legge di Benford}

La probabilità $P(n)$ è definita come la probabilità che un numero $M$ rientri nell'intervallo $\left [n \cdot 10^k, (n+1) \cdot 10^k\right ] $ per qualsiasi scala $k$:

\begin{equation}
P(n) = \sum_{k} \int_{n \cdot 10^{k}}^{(n+1) \cdot 10^{k}} \rho(M) dM
\end{equation}

Questa formulazione approssimativa somma tutte le possibili scale $k$.

Il secondo punto fondamentale è l' \textbf{Invarianza di Scala}. Se si cambia l'unità di misura (la scala) da $M$ a $bM$ (ad esempio, lunghezza di un fiume in km o cm), la distribuzione di probabilità $\rho(M)$ non deve cambiare.
L'unica distribuzione di probabilità che soddisfa tale relazione di scaling è la legge di potenza $\rho(M) \propto 1/M$.

Sostituendo $\rho(M) \sim 1/M$ nell'integrale, si ottiene:
\begin{equation}
P(n) = \sum_{k} \int_{n \cdot 10^{k}}^{(n+1) \cdot 10^{k}} dM \frac{1}{M} = \sum_{k} \{\log\left [ (n+1)10^{k}\right ]  - \log\left [n10^{k}\right ] \}
\end{equation}
\begin{equation}
P(n) = \sum_{k} \log \frac{n+1}{n}
\end{equation}

Dopo la normalizzazione, la Legge di Benford è data da:
\begin{equation}
P(n) = \log_{10}\left(\frac{n+1}{n}\right)
\end{equation}


\subsubsection{Processo Additivo}

Un processo additivo è definito come una passeggiata aleatoria in cui si aggiunge un termine casuale $\xi$ ad ogni passo temporale $t$:
\begin{equation}
N_{t+1} = N_t + \xi
\end{equation}

Per $\xi$ distribuito simmetricamente con varianza finita, $N_t$ è la somma di $t$ variabili casuali.
La varianza $\sigma^2$ scala linearmente con $t$ e la deviazione standard scala come $\sigma \sim \sqrt{t}$.
Per $t \to \infty$, la distribuzione di $N_t$ diventa sempre più larga e tende a una distribuzione uniforme (che non è una distribuzione di scala).

\subsubsection{Processo Moltiplicativo}

Un processo moltiplicativo (tipico di crescita di popolazioni o variazioni percentuali del mercato azionario) è definito da:
\begin{equation}
N_{t+1} = N_t \cdot \xi
\end{equation}

Prendendo il logaritmo, il processo moltiplicativo si trasforma in un processo additivo nello spazio logaritmico:
\begin{equation}
\log(N_{t+1}) = \log(N_t) + \log(\xi)
\end{equation}

Per $t \to \infty$, il Teorema Centrale del Limite si applica a $\log(N_t)$, e la sua distribuzione converge a una Gaussiana. La distribuzione di $\log(N_t)$ diventa uniforme nello spazio logaritmico.

Una distribuzione uniforme su $\log(N)$ implica $\int p(\log N) d \log N \sim \int \frac{dN}{N}$, ovvero la distribuzione nello spazio lineare è $\rho(N) \propto 1/N$.

La dipendenza da $\frac{1}{N}$ è una caratteristica fondamentale delle distribuzioni che seguono la legge di Benford.

\subsection{La Distribuzione Log-Normale}

Quando una variabile casuale $X$ è il risultato del prodotto di un gran numero $N$ di variabili casuali indipendenti $y_i$:
$$X = \prod_{i=1}^{N} y_i$$

il suo logaritmo sarà la somma di molte variabili casuali:
$$\log X = \sum_{i=1}^{N} \log(y_i)$$

Se le variabili $\log(y_i)$ hanno media $\mu_l$ e varianza $\sigma_l^2$ finite, allora la distribuzione di $\log X$ tenderà a una distribuzione normale (Gaussiana) con media $N\mu_l$ e varianza $N\sigma_l^2$.
$$\tilde{\rho}(\log X) = C \exp\left\{-\frac{(\log X - N\mu_l)^2}{2\sigma_l^2 N}\right\}$$

Di conseguenza, la variabile $X$ segue una \textbf{distribuzione log-normale}. La sua funzione di densità di probabilità $\rho(X)$ si ottiene dalla densità del suo logaritmo $\tilde{\rho}(\log X)$ con un cambio di variabile:
$$\rho(X)dX = \tilde{\rho}(\log X) d\log X = \tilde{\rho}(\log X) \frac{dX}{X}$$

Questo porta alla forma esplicita della distribuzione log-normale:
$$\rho(X) = \frac{C}{X} \exp\left\{-\frac{(\log X - N\mu_l)^2}{2\sigma_l^2 N}\right\}$$

Vogliamo stimare quantitativamente l'intervallo di valori di $x$ per cui il termine esponenziale è vicino a 1, rendendo la distribuzione log-normale di fatto una legge di potenza del tipo $\frac{C}{X}$.

Consideriamo dei valori tipici per i parametri, come $\sigma_{l} = 1$ e $N=100$. La deviazione standard del termine logaritmico è $\sigma_{l}\sqrt{N} = 10$.

Imponiamo che il termine nell'argomento dell'esponenziale sia piccolo, ad esempio compreso tra 0 e 1:
$$ 0 \le \frac{(\ln X - N\mu_{l})^2}{2\sigma_{l}^2 N} \le 1 $$
Sostituendo i valori, otteniamo:
$$ (\ln X - N\mu_{l})^2 \le 200 $$
Estraendo la radice quadrata:
$$ |\ln X - N\mu_{l}| \le \sqrt{200} \approx 14 $$
Questo significa che il logaritmo di $X$ può discostarsi dalla media $N\mu_{l}$ di circa 14. In scala lineare, questo equivale a un fattore moltiplicativo:
$$ e^{-14} \le \frac{X}{e^{N\mu_{l}}} \le e^{14} $$
Poiché $e^{14} \approx 1.2 \times 10^6$, vediamo che l'approssimazione a legge di potenza vale su un intervallo che si estende per circa 6 ordini di grandezza. Su questa vastissima scala, le correzioni introdotte dal termine esponenziale sono trascurabili e il comportamento dominante della distribuzione è effettivamente:
$$ \rho(x) \approx \frac{C}{X} $$


\begin{figure}[h!]
    \centering
    \begin{tikzpicture}[scale=0.75]
        \begin{axis}[
            % Assi posizionati a sinistra e in basso (esterni)
            axis lines = left,
            % Posizionamento etichetta X: alla fine dell'asse (cs:1), sotto (anchor=north)
            xlabel={$\log x$},
            xlabel style={at={(ticklabel* cs:1)}, anchor=north west},
            % Posizionamento etichetta Y: in cima all'asse (cs:1), a sinistra (anchor=south east)
            ylabel={$\log \rho(x)$},
            ylabel style={at={(ticklabel* cs:1)}, anchor=south east},
            % Stile font
            label style={font=\Large},
            % Griglia
            grid=both,
            minor tick num=1,
            grid style={line width=.1pt, draw=gray!30},
            major grid style={line width=.2pt, draw=gray!50},
            % Rimuove i numeri
            xtick=\empty,
            ytick=\empty,
            % Range
            xmin=0, xmax=10,
            ymin=0, ymax=10,
            % Dimensioni
            width=10cm,
            height=8cm,
            % Disattiva il clipping per mostrare le label esterne
            clip=false
        ]

            % Disegno della curva (Power law + Exponential cutoff)
            % f(x) = -x + costante - termine_esponenziale
            \addplot [
                domain=0:7.75, 
                samples=100, 
                color=cyan, 
                line width=2pt
            ] {9 - x - exp(x-7.5)};

            % Linea tratteggiata rosa (Cut-off)
            \draw [dashed, magenta!60, line width=2pt] (axis cs:6.5, 0) -- (axis cs:6.5, 1.8);

            % Etichetta della pendenza (~ 1/x)
            % In scala log-log, 1/x è una retta con pendenza -1
            \node [magenta!60, font=\Large] at (axis cs:6.2, 3.8) {$\sim \frac{1}{x}$};

        \end{axis}
    \end{tikzpicture}
    \caption{Rappresentazione log-log della densità $\rho(x)$ che mostra un regime a legge di potenza seguito da un cut-off.}
    \label{fig:log_log_plot}
\end{figure}


\newpage
\section{La Legge di Zipf}

La \textbf{Legge di Zipf} è una legge empirica che descrive la frequenza di un elemento in funzione del suo rango (la sua posizione in una classifica ordinata per frequenza). Matematicamente, è una legge di potenza della forma $f(R) \propto R^{-\alpha}$, dove $R$ è il rango e $\alpha$ è un esponente caratteristico.

Per essere una distribuzione di probabilità, deve essere normalizzata. Assumendo che il rango vada da 1 a $R_{max}$:
$$\int_{1}^{R_{max}} f(R) dR = 1$$

La costante di normalizzazione $C$ dipende in modo cruciale dal valore dell'esponente $\alpha$.

\begin{itemize}
    \item \textbf{Caso $\alpha < 1$}: L'integrale è dominato dal limite superiore $R_{max}$. La costante di normalizzazione dipende da $R_{max}$:
    $$C \int_{1}^{R_{max}} R^{-\alpha} dR = \frac{C}{1-\alpha}(R_{max}^{1-\alpha} - 1) \approx \frac{C}{1-\alpha}R_{max}^{1-\alpha} = 1$$
    Da cui si ottiene $C \approx (1-\alpha)R_{max}^{\alpha-1}$. La distribuzione è:
    $$f(R) = (1-\alpha)R_{max}^{\alpha-1} R^{-\alpha}$$
    
    \item \textbf{Caso $\alpha > 1$}: L'integrale converge anche per $R_{max} \to \infty$. In questo limite, la costante di normalizzazione è $C = \alpha - 1$. La distribuzione è:
    $$f(R) = (\alpha - 1) R^{-\alpha}$$
    
    \item \textbf{Caso $\alpha = 1$}: L'integrale diverge logaritmicamente. La normalizzazione richiede un $R_{max}$ finito.
    $$C \int_{1}^{R_{max}} \frac{dR}{R} = C \log R_{max} = 1 \implies C = \frac{1}{\log R_{max}}$$
    La distribuzione diventa:
    $$f(R) = \frac{1}{\log R_{max}} R^{-1}$$
    
\end{itemize}

\section{Legge di Heaps}

La \textbf{Legge di Heap}  descrive il numero di parole distinte (vocaboli) $D(N)$ in funzione del numero totale di parole $N$ scansionate in un testo. Questo è un esempio di processo di innovazione.



Il processo conta $+1$ ogni volta che si incontra una parola che non è mai apparsa prima (una "novità").

\begin{itemize}
    \item Se l'insieme di parole fosse finito (modello della scatola con un numero finito di palline), $D(N)$ tenderebbe a saturare esponenzialmente.
    \item Nei testi reali, $D(N)$ segue asintoticamente una \textbf{legge di potenza}:
    \begin{equation}
    D(N) \sim N^{r}
    \end{equation}
\end{itemize}

L'esponente $r$ dipende dal corpus e dalla lingua, ma deve essere sempre minore di 1. Non può essere maggiore di 1, poiché non si può contare più di una nuova parola per unità di tempo intrinseco (una parola scansionata).


Il \textbf{tasso di innovazione} è la derivata del numero di parole distinte, $dD(N)/dN$.
\begin{itemize}
    \item All'inizio, il tasso di innovazione è $\approx 1$ (tutte le parole sono nuove).
    \item Asintoticamente, il tasso di innovazione decresce con una legge di potenza:
    \begin{equation}
    \frac{dD(N)}{dN} \sim N^{r-1}
    \end{equation}
\end{itemize}


Esiste una profonda connessione tra l'esponente della legge di Zipf, $\alpha$, e quello della legge di Heaps, $r$. Questa relazione può essere derivata imponendo la condizione che la frequenza dell'elemento meno comune (quello con rango $R_{max}$) sia dell'ordine di $\frac{1}{N}$. Poiché $D(N)$ è il numero totale di elementi unici, possiamo identificare $R_{max}$ con $D(N)$.
$$f(R_{max}) = \frac{1}{N}$$

Sostituendo le diverse forme della legge di Zipf, si ottengono le seguenti relazioni:
\begin{itemize}
    \item \textbf{Caso $\alpha > 1$}:
    $$(\alpha - 1) R_{max}^{-\alpha} = \frac{1}{N} \implies R_{max} \sim N^{1/\alpha}$$
    Confrontando con $D(N) \sim N^r$, otteniamo $R_{max} = D(N)$ e quindi la \textbf{relazione di Heaps-Zipf}:
    $$r = \frac{1}{\alpha}$$
    Ad esempio, per $\alpha=2$ si ha $r=1/2$, quindi $D(N) \sim N^{0.5}$.

    \item \textbf{Caso $\alpha = 1$}:
    $$\frac{R_{max}^{-1}}{\log R_{max}} = \frac{1}{N} \implies R_{max}(N) \sim \frac{N}{\log N}$$

    \item \textbf{Caso $\alpha < 1$}:
    $$(1-\alpha)R_{max}^{-1} = \frac{1}{N} \implies R_{max}(N) \sim N$$
\end{itemize}

\begin{figure}[h!]
    \centering
    
    % --- PRIMO GRAFICO: D(N) ---
    \begin{minipage}{0.48\textwidth}
        \centering
        \begin{tikzpicture}[scale=0.7]
            \begin{axis}[
                axis lines = left,
                xlabel={$N$},
                ylabel={$D(N)$},
                xlabel style={at={(ticklabel* cs:1)}, anchor=north west},
                ylabel style={at={(ticklabel* cs:1)}, anchor=south east},
                label style={font=\Large},
                grid=both,
                minor tick num=1,
                grid style={line width=.1pt, draw=gray!30},
                major grid style={line width=.2pt, draw=gray!50},
                xtick=\empty, ytick=\empty,
                xmin=0, xmax=10,
                ymin=0, ymax=10,
                width=\textwidth,
                clip=false
            ]
                % Funzione smooth: passa da lineare a radice quadrata
                % f(x) = x / (1 + (x/5)^2)^0.25 approssima bene l'andamento
                \addplot [
                    domain=0:9.5, 
                    samples=100, 
                    color=cyan!70!blue, 
                    line width=1.8pt
                ] {x / (1 + (x/4.5)^4)^0.18 * 1.1};

                % Linea tratteggiata rosa al punto di flesso
                \draw [dashed, magenta!60, line width=1.2pt] (axis cs:5.2, 0) -- (axis cs:5.2, 5.5);
                
                % Label rosa posizionate come nel disegno
                \node [magenta!60, font=\Large, anchor=north west] at (axis cs:2.5, 2.5) {$\sim N$};
                \node [magenta!60, font=\Large, anchor=south west] at (axis cs:6.5, 4.) {$\sim N^{0.5}$};
            \end{axis}
        \end{tikzpicture}
    \end{minipage}
    \hfill
    % --- SECONDO GRAFICO: Innovation rate ---
    \begin{minipage}{0.48\textwidth}
        \centering
        \begin{tikzpicture}[scale=0.7]
            \begin{axis}[
                axis lines = left,
                xlabel={$N$},
                ylabel={Innovation rate},
                xlabel style={at={(ticklabel* cs:1)}, anchor=north west},
                ylabel style={at={(ticklabel* cs:1)}, anchor=south east},
                label style={font=\Large},
                grid=both,
                minor tick num=1,
                grid style={line width=.1pt, draw=gray!30},
                major grid style={line width=.2pt, draw=gray!50},
                xtick=\empty, ytick=\empty,
                xmin=0, xmax=10,
                ymin=0, ymax=10,
                width=\textwidth,
                clip=false
            ]
% Plateau costante
                \addplot [
                    domain=0:4, samples=10, color=cyan!60!blue, line width=2pt
                ] {8};
                
                % NUOVO DECADIMENTO: "Rivolto verso il basso"
                % Usiamo una parabola invertita per dare la concavità verso il basso
                \addplot [
                    domain=4:11, samples=100, color=cyan!60!blue, line width=2pt
                ] {8 - 0.13*(x-4)^2};

                % Linea tratteggiata
                \draw [dashed, magenta!50, line width=2pt] (axis cs:4, 0) -- (axis cs:4, 8);
              
                % Label rosa
                \node [magenta!60, font=\Large, anchor=south west] at (axis cs:4.5, 4.3) {$\sim N^{-0.5}$};
            \end{axis}
        \end{tikzpicture}
    \end{minipage}

    \vspace{1em}
    \caption{Evoluzione del numero di elementi distinti $D(N)$ e del relativo tasso di innovazione.}
    \label{fig:confronto_corretto}
\end{figure}


\section{Processo di Poisson e Legge di Taylor}

Un modo alternativo per modellare la comparsa di novità è attraverso un \textbf{processo di Poisson}. Se la probabilità di scoprire una novità in un piccolo intervallo di tempo è costante e pari a $p$, allora il numero di novità $D(t)$ accumulate fino al tempo $t$ segue una distribuzione di Poisson con media $\mu = pt$.
$$P\{D(t)=m\} = \frac{(pt)^m e^{-pt}}{m!}$$

Possiamo calcolare media e varianza di questo processo utilizzando la funzione generatrice dei momenti, $M(z) = \mathbb{E}[e^{zD(t)}]$.
$$M(z) = e^{-pt} e^{pte^z}$$

La media $\mu$ è la derivata prima di $M(z)$ calcolata in $z=0$:
$$\mu = \mathbb{E}[D(t)] = \left. \frac{d}{dz} M(z) \right|_{z=0} = pt$$

Il momento secondo è la derivata seconda calcolata in $z=0$:
$$\mathbb{E}[D^2(t)] = \left. \frac{d^2}{dz^2} M(z) \right|_{z=0} = pt + (pt)^2$$

La varianza $\sigma^2$ è quindi:
$$\sigma^2 = \mathbb{E}[D^2(t)] - \mu^2 = (pt + (pt)^2) - (pt)^2 = pt$$

Otteniamo quindi una relazione fondamentale per il processo di Poisson:
$$\sigma^2 = \mu \implies \sigma = \mu^{1/2}$$
Questa è una forma specifica della \textbf{Legge di Taylor}, che lega la varianza alla media con una legge di potenza.

\subsubsection{Derivazione della Distribuzione di Poisson}

La distribuzione di Poisson può essere derivata come caso limite di una \textbf{distribuzione binomiale}. La distribuzione binomiale descrive la probabilità di ottenere $m$ successi in $N$ prove indipendenti, dove la probabilità di successo in ogni prova è $p$.
$$P(m) = \binom{N}{m} p^m (1-p)^{N-m}$$

Consideriamo il limite per $N \to \infty$ e $p \to 0$ in modo tale che il numero medio di successi $\mu = Np$ rimanga costante.
Sostituendo $p = \frac{\mu}{N}$ nella formula binomiale:
$$P(m) = \frac{N!}{m!(N-m)!} \left(\frac{\mu}{N}\right)^m \left(1-\frac{\mu}{N}\right)^{N-m}$$

Analizzando i termini per $N$ molto grande:
\begin{itemize}
    \item $\frac{N!}{ (N-m)! N^m } = \frac{N(N-1)\cdots(N-m+1)}{N^m} \approx 1$
    \item $\left(1-\frac{\mu}{N}\right)^N \approx e^{-\mu}$
    \item $\left(1-\frac{\mu}{N}\right)^{-m} \approx 1$
\end{itemize}

Mettendo insieme questi limiti, si ottiene la formula della distribuzione di Poisson:
$$P(m) = \frac{\mu^m e^{-\mu}}{m!}$$



\chapter{Lezione 5}
\label{chap:lezione_05} 

\begin{flushright}
\textit{Data: 15/10/2025}
\end{flushright}

% Comandi personalizzati
\newcommand{\deriv}[2]{\frac{d #1}{d #2}}
\newcommand{\pderiv}[2]{\frac{\partial #1}{\partial #2}}


\section{Introduzione allo Scaling e Leggi di Potenza}

In questa lezione si affrontano i concetti di variabili di scaling, leggi di potenza e modelli generativi. Sebbene nelle lezioni precedenti siano stati introdotti strumenti per trattare le distribuzioni a coda lunga (long-tails) e le leggi di potenza, non è stata ancora fornita una spiegazione teorica fondamentale sul \textit{perché} tali distribuzioni siano così ubiquitarie in natura.

\subsection{Argomentazioni di Galileo sullo Scaling}
Il concetto di scaling non è recente; Galileo Galilei si interrogava su queste problematiche in modo molto concreto, analizzando le dimensioni delle ossa negli animali (si veda il "Dialogo sopra i due massimi sistemi del mondo").

Consideriamo la relazione tra lunghezza caratteristica $L$, massa $M$ e resistenza strutturale $S$ (Strength).
Se scaliamo la dimensione lineare di un oggetto di un fattore (ad esempio raddoppiandola):
\begin{equation}
    L \to 2L \implies M \propto L^3 \to 8M
\end{equation}
La massa scala con il cubo della lunghezza (assumendo densità costante). Tuttavia, la resistenza di un osso dipende dalla sua sezione trasversale $A \propto R^2$ (dove $R$ è il raggio o diametro dell'osso):
\begin{equation}
    S \propto R^2 \quad (\text{o } D^2)
\end{equation}
Se scaliamo geometricamente l'animale mantenendo le proporzioni ($R \propto L$):
\begin{equation}
    S \to 4S
\end{equation}
Si crea una discrepanza: il peso aumenta di un fattore 8, ma la capacità di sostenerlo aumenta solo di un fattore 4.
\begin{equation}
    \frac{\text{Peso}}{\text{Resistenza}} \propto \frac{L^3}{L^2} = L
\end{equation}
Questa è la ragione per cui non possono esistere giganti geometricamente simili agli umani: le loro ossa cederebbero sotto il proprio peso. Affinché la struttura regga (ovvero, per mantenere costante lo stress $S/M \approx cost$), la sezione dell'osso dovrebbe scalare come la massa:
\begin{equation}
    R^2 \sim M \sim L^3 \implies R \sim L^{3/2}
\end{equation}
Le ossa dovrebbero quindi diventare sproporzionatamente grosse rispetto alla lunghezza.
Galileo fornì una prova sperimentale utilizzando due lenze da pesca con diametri $d_1$ e $d_2$. Verificò che il rapporto del carico di rottura corrispondeva esattamente al rapporto dei quadrati dei diametri:
\begin{equation}
    \frac{S_1}{S_2} = \frac{d_1^2}{d_2^2}
\end{equation}

\subsection{Scaling in Biologia: Relazioni Allometriche}
In biologia, le leggi di scaling sono onnipresenti e prendono il nome di \textbf{Relazioni Allometriche}. Assumono la forma generale:
\begin{equation}
    Y = Y_0 M^b
\end{equation}
dove $Y$ è una variabile biologica e $M$ è la massa.
\begin{itemize}
    \item \textbf{Legge di Kleiber:} Il metabolismo basale $B$ scala come $M^{3/4}$. Questo esponente $3/4$ (invece del $2/3$ geometrico superficiale) si applica sorprendentemente anche a interi ecosistemi.
    \item \textbf{Circolazione:} Il tempo di circolazione sanguigna scala come $T \sim M^{1/4}$.
    \item \textbf{Frequenza cardiaca:} Scala con esponente negativo $\omega \sim M^{-1/4}$. Il cuore di un elefante batte molto più lentamente di quello di un topo.
\end{itemize}
Esistono dibattiti teorici aperti sull'origine esatta della legge dei $3/4$ (si vedano i lavori di West, Brown, Enquist). È cruciale distinguere tra il dato empirico misurato e la teoria ("la lente") usata per interpretarlo.

\subsection{Scaling nei Sistemi Sociali}
Nei sistemi sociali e urbani, osserviamo diverse classi di scaling rispetto alla popolazione $N$:
\begin{equation}
    Y \sim N^{\beta}
\end{equation}
\begin{itemize}
    \item \textbf{$\beta > 1$ (Superlineare):} Quantità legate all'interazione sociale (innovazione, brevetti, crimine). Qui "$1+1 > 2$".
    \item \textbf{$\beta \approx 1$ (Lineare):} Quantità legate ai bisogni individuali (es. numero di case).
    \item \textbf{$\beta < 1$ (Sublineare):} Infrastrutture (strade, reti elettriche). Ottimizzazione e economie di scala.
\end{itemize}

\section{Derivazione della Legge di Potenza dallo Scaling}

Un modello generativo deve spiegare perché queste leggi emergono. Partiamo dall'ipotesi di invarianza di scala per una distribuzione di probabilità $P(x)$.
La relazione di scaling impone che riscalando la variabile $x$ di un fattore $\lambda$, la distribuzione rimanga la stessa a meno di un pre-fattore:
\begin{equation}
    P(\lambda x) = g(\lambda) P(x)
\end{equation}
Vogliamo determinare la forma funzionale di $P(x)$.
Per $x=1$, l'equazione diventa $P(\lambda) = g(\lambda)P(1)$, da cui ricaviamo la funzione di scaling:
\begin{equation}
    g(\lambda) = \frac{P(\lambda)}{P(1)}
\end{equation}
Sostituendo nell'equazione originale:
\begin{equation}
    P(\lambda x) = \frac{P(\lambda)}{P(1)} P(x)
\end{equation}
Differenziamo entrambi i membri rispetto a $\lambda$:
\begin{equation}
    \pderiv{}{\lambda} P(\lambda x) = \frac{P(x)}{P(1)} \deriv{P(\lambda)}{\lambda}
\end{equation}
Applicando la regola della catena a sinistra ($y=\lambda x \to dy/d\lambda = x$):
\begin{equation}
    x P'(\lambda x) = \frac{P(x)}{P(1)} P'(\lambda)
\end{equation}
Poniamo ora $\lambda = 1$:
\begin{equation}
    x P'(x) = \frac{P'(1)}{P(1)} P(x)
\end{equation}
Il termine $\frac{P'(1)}{P(1)}$ è una costante. Definiamola come $-\beta$.
Otteniamo l'equazione differenziale:
\begin{equation}
    \frac{P'(x)}{P(x)} = -\frac{\beta}{x} \implies \deriv{(\ln P)}{x} = -\frac{\beta}{x}
\end{equation}
Integrando:
\begin{tcolorbox}[colback=yellow!25, colframe=yellow!75!orange, coltitle=black, title=\textbf{Legge di Potenza}]
\begin{equation}
    P(x) = P(0) x^{-\beta}
\end{equation}
L'invarianza di scala (o "zoom-in/zoom-out invariance") implica necessariamente una distribuzione a legge di potenza (scale-free).
\end{tcolorbox}

\section{Combinazione di Esponenziali}

Analizziamo un primo modello generativo ("toy model") che ottiene una legge di potenza combinando distribuzioni esponenziali.
Sia $X$ una variabile casuale con distribuzione esponenziale:
\begin{equation}
    \rho_X(x) = c_1 e^{ax}
\end{equation}
Definiamo una nuova variabile $Y$ attraverso una trasformazione esponenziale:
\begin{equation}
    Y = c_2 e^{bX}
\end{equation}
Vogliamo trovare la distribuzione $\tilde{\rho}(y)$.
Invertiamo la relazione:
\begin{equation}
    \frac{Y}{c_2} = e^{bX} \implies bX = \ln\left(\frac{Y}{c_2}\right) \implies X = \frac{1}{b} \ln\left(\frac{Y}{c_2}\right)
\end{equation}
Usiamo la conservazione della probabilità: $\tilde{\rho}(y) dy = \rho_X(x) dx$.
Calcoliamo lo Jacobiano della trasformazione inversa:
\begin{equation}
    \frac{dX}{dY} = \frac{d}{dY} \left[ \frac{1}{b} (\ln Y - \ln c_2) \right] = \frac{1}{bY}
\end{equation}
Sostituiamo nella formula:
\begin{equation}
    \tilde{\rho}(y) = \rho_X(X(y)) \left| \frac{dX}{dY} \right| = c_1 e^{a \left[ \frac{1}{b} \ln\left(\frac{y}{c_2}\right) \right]} \cdot \frac{1}{bY}
\end{equation}
Definiamo $K = c_1/b$. Manipoliamo il termine esponenziale:
\begin{equation}
    e^{\ln\left( (y/c_2)^{a/b} \right)} = \left( \frac{y}{c_2} \right)^{\frac{a}{b}} \propto y^{\frac{a}{b}}
\end{equation}
Quindi:
\begin{equation}
    \tilde{\rho}(y) = K y^{\frac{a}{b}} \cdot y^{-1} = K y^{\frac{a}{b} - 1} = K y^{-\left(1 - \frac{a}{b}\right)}
\end{equation}
Abbiamo ottenuto una legge di potenza $P(y) \sim y^{-\beta}$ con esponente $\beta = 1 - \frac{a}{b}$.

\section{Monkey Typewriter}
Sia $m$ il numero di caratteri distinti disponibili (es. 26 lettere).
Sia $p$ la probabilità di premere la barra spaziatrice (generazione di un fine parola).
La probabilità di premere un qualsiasi altro tasto specifico è $p_l = \frac{1-p}{m}$. Una parola di lunghezza $S$ è composta da $S$ caratteri specifici seguiti da uno spazio. Poiché l'ordine conta (le lettere sono fissate), la probabilità di generare una \textit{specifica} parola di lunghezza $S$ è:
\begin{equation}
    P(\text{word}_S) = \left( \frac{1-p}{m} \right)^S \cdot p
\end{equation}
Possiamo riscrivere questa probabilità come un esponenziale della lunghezza:
\begin{equation}
    P(\text{word}_S) = p \cdot e^{S \ln\left(\frac{1-p}{m}\right)} = p \cdot  e^{bS}
\end{equation}
con $b = \ln\left(\frac{1-p}{m}\right) < 0$.


Dunque, questo significa che più lunga è la parola, minore è la sua probabilità, dato che $b < 0$. In altre parole, la distribuzione delle lunghezze delle parole decade esponenzialmente. Tuttavia, per ogni lunghezza $S$ ci sono $m^S$ possibili combinazioni distinte di lettere, quindi il \textit{numero di parole} di tale lunghezza cresce esponenzialmente con $S$. Abbiamo pertanto due esponenziali opposti: uno decrescente (la probabilità di ogni sequenza) e uno crescente (il numero di sequenze possibili).

\subsection{Relazione Rango-Frequenza}
Vogliamo trovare la relazione tra la frequenza della parola (probabilità) e il suo rango $r$ (classifica dalla più frequente alla meno frequente).
Il numero totale di parole possibili con lunghezza $\le S$ determina il rango massimo associato alla lunghezza $S$.
Somma della serie geometrica delle parole possibili (m parole lunghe 1, $m^2$ lunghe 2, ...):
\begin{equation}
    r(S) \approx \sum_{i=1}^{S} m^i = \frac{m^{S+1}-m}{m-1} \approx \frac{m^{S+1}}{m-1} \approx m^S
\end{equation}
Invertendo per grandi $S$:
\begin{equation}
    S \approx \log_m r
\end{equation}
Sostituiamo $S$ nell'espressione della probabilità (frequenza $f(r)$):
\begin{equation}
    f(r) \sim \left( \frac{1-p}{m} \right)^{\log_m r}
\end{equation}
Usiamo l'identità $x^{\log_m r} = r^{\log_m x}$:
\begin{equation}
    f(r) \sim r^{\log_m \left( \frac{1-p}{m} \right)} = r^{\log_m(1-p) - \log_m m} = r^{\log_m(1-p) - 1}
\end{equation}
Questa è la Legge di Zipf $f(r) \sim r^{-\alpha}$ con esponente:
\begin{equation}
    \alpha = 1 - \log_m(1-p)
\end{equation}
Pertanto, anche questo processo puramente casuale, una ``scimmia'' che digita casualmente lettere e spazi, produce una \textit{distribuzione a legge di potenza} nella relazione rango-frequenza delle parole. 

Inserendo valori reali ($m=26, p \approx 0.2$), si ottiene $\alpha \approx 1.01$, che è sorprendentemente vicino al valore reale della legge di Zipf ($\alpha \approx 1$).
Tuttavia, il modello è concettualmente errato per il linguaggio umano (eccetto forse le lingue agglutinanti), poiché prevede che le parole più corte siano sempre le più frequenti in modo monotono, mentre i dati reali mostrano una distribuzione log-normale o simile per le lunghezze.
Più in generale, tutto ciò a cui applichiamo solitamente le combinazioni di esponenziali finisce per funzionare in un intervallo molto specifico prima che sia necessario considerare qualche tipo di correzione empirica.

\begin{figure}[h]
\centering
\begin{tikzpicture}[scale=0.9]
    \begin{axis}[
        width=10cm, height=6cm,
        xlabel={Lunghezza Parola $S$},
        ylabel={Probabilità $P(S)$},
        ymode=log,
        domain=1:20, % Partiamo da 1, S=0 non ha senso fisico
        xmin=1, xmax=20,
        legend pos=north east, % Spesso copre meno dati in alto a destra o sud-ovest
        grid=major,
        samples=100 % Aumentato per curve più morbide
    ]
    
    % Modello Monkey: Decadimento esponenziale puro
    % P(S) ~ exp(-bS). Normalizzato a occhio per il grafico.
    \addplot[blue, thick, dashed] {0.8*exp(-0.4*x)};
    \addlegendentry{\small Modello Monkey}

    % Realtà: Distribuzione Gamma/Lognormale
    % x^alpha * exp(-beta*x). Picco a x = alpha/beta.
    % Se alpha=5 e beta=1, picco a 5.
    \addplot[red, thick] {0.05 * x^5 * exp(-1*x)}; 
    \addlegendentry{\small Dati Reali}
    
    % Annotazione
    \node[coordinate, pin=above:{Picco tipico ($\sim 5$ chars)}] at (axis cs: 5, 0.05 * 3125 * 0.0067) {};
    
    \end{axis}
\end{tikzpicture}
\caption{Confronto tra la previsione teorica e la distribuzione reale della lunghezza delle parole.}
\end{figure}

\section{Inversione dei Quantili}

Sebbene si tratti di un passaggio elementare, è opportuno esaminarlo; essenzialmente, partendo da un processo di tipo legge di potenza (\textit{power law}), si osserverà conseguentemente una legge di potenza. Si consideri una variabile stocastica distribuita secondo $\rho(y)$ e una variabile stocastica derivata $x$ tale che $x = 1/y$. Si ha:

\begin{equation}
    \tilde{\rho}(x) = \rho(y) \left| \frac{dy}{dx} \right| = \rho\left(\frac{1}{x}\right) \frac{1}{x^2} \xrightarrow[x \gg 1]{y \ll 1} \rho(0) x^{-2} \sim x^{-2}.
    \label{eq:quantile_inversion}
\end{equation}

Di conseguenza, indipendentemente dalla distribuzione della variabile $y$, il comportamento asintotico della coda segue una legge di potenza. Ad esempio, nel modello di Ising, analizzando le fluttuazioni della magnetizzazione attorno alla media, definite come $\delta = \Delta m$, si osserva che $\delta \sim m^{-1}$, da cui ne deriva che $\rho(\delta) \sim m^{-2}$ nelle code della distribuzione.

\section{Processo di Yule-Simon}
Dedicheremo un'analisi approfondita a questo modello generativo di leggi di potenza, dato il suo notevole interesse teorico e la sua vasta applicabilità nei processi di innovazione: dalla biologia alla generazione del linguaggio e, più in generale, nella modellizzazione della frequenza con cui l'innovazione emerge nelle dinamiche sociali. Tali caratteristiche rendono questo modello uno standard di riferimento per la generazione di distribuzioni a coda pesante (\textit{heavy-tailed}).

Il modello giustifica un comportamento osservato in una molteplicità di sistemi, noto informalmente come \textbf{Rich-Get-Richer}  o \textbf{Preferential Attachment}. Ad esempio, nella distribuzione delle specie, è ragionevole attendersi che le specie più antiche siano anche le più diffuse. Il medesimo principio si applica a numerosi altri domini: parole già frequenti tendono ad apparire con frequenza ancora maggiore, articoli scientifici già citati hanno una maggiore probabilità di ricevere ulteriori citazioni...

In sintesi, il modello Yule–Simon fornisce il fondamento matematico per questo processo di auto-rinforzo, in cui la probabilità di acquisire nuovi membri (o occorrenze) cresce in funzione del numero già accumulato. Questa dinamica conduce naturalmente a distribuzioni a coda pesante (simili a leggi di potenza), dove un numero esiguo di entità domina il sistema, mentre la maggioranza rimane rara.


\subsubsection*{Il Processo di Yule (1922)}
Sebbene l'analisi di Simon sia successiva di circa trent'anni, è Yule a definire le dinamiche fondamentali del processo. Originariamente formulato in ambito biologico, introdurremo qui i postulati utilizzando la terminologia linguistica. Le regole generali che governano il processo sono le seguenti:

\begin{itemize}
    \item Nuove specie possono emergere nel sistema, ma non possono estinguersi.
    \item Una categoria (o genere, in senso biologico) a cui sono associate $k$ specie, acquisisce nuove specie con un tasso proporzionale al numero di specie già possedute (meccanismo \textit{rich-get-richer}).
    \item Occasionalmente, una parola diverge talmente da un'altra, o una mutazione apporta variazioni così significative, che la nuova specie non risulta classificabile nel genere di origine, determinando così la creazione di un nuovo genere o di una nuova categoria.
\end{itemize}

\noindent Mentre i principi sono attribuiti a Yule, lo sviluppo formale del processo è ascritto a Simon (1955). 

\subsubsection{Definizione e Master Equation}

Si consideri lo stato iniziale al tempo $t = 0$: il sistema è costituito da un singolo elemento (una parola in un dizionario o una specie in un genere). Ad ogni passo temporale $t \to t + 1$, viene aggiunto una nuova specie secondo le seguenti regole. Se definiamo $E(t)$ come l'evento di assegnazione al tempo $t$:

\begin{equation}
    E(t+1) = 
    \begin{cases} 
    \text{Raggruppato in un nuovo genere} & \text{Prob} = p \\
    \text{Raggruppato in un genere preesistente} & \text{Prob} = 1-p 
    \end{cases}
    \label{eq:assignment_prob}
\end{equation}

Si osservi che, indipendentemente dall'esito, un nuovo elemento viene introdotto nel sistema; la questione dirimente è se tale elemento costituisca un nuovo genere o se vada ad accrescere un genere preesistente.
Sia $D(t)$ il numero di generi distinti presenti al tempo $t$. Poiché la probabilità di apparizione di un nuovo genere è $p$, ci attendiamo che il tasso di variazione di $D(t)$ sia pari a $p$:

\begin{equation}
    \frac{d}{dt}D(t) = p \implies D(t) = D_{0} + p t.
\end{equation}
Si nota una dipendenza lineare in $t$, analoga alla legge di Heaps.


\noindent Ci interroghiamo ora sulla distribuzione del numero di generi aventi $k$ specie al tempo $t$, denotata come $P_{k,t}$. Per determinare tale quantità, utilizzeremo il formalismo della \textbf{Master Equation}, combinato con l'approssimazione al continuo. Nel passaggio temporale $t \to t+1$, l'aggiunta di una nuova parola (o specie) può comportare una delle seguenti tre transizioni per la variabile $P_{k,t}$:

\begin{enumerate}
    \item Un genere con $k-1$ specie acquisisce una nuova specie, incrementando il contatore $P_{k,t}$ di un'unità.
    \item Un genere con $k$ specie acquisisce una nuova specie (diventando un genere con $k+1$ specie), decrementando il contatore $P_{k,t}$ di un'unità.
    \item Un nuovo genere appare con una singola specie, incrementando il contatore $P_{1,t}$ (questo termine influenza solo il caso $k=1$).
\end{enumerate}

\noindent Formalizziamo matematicamente questi contributi. Ricordando che il fattore $(1-p)$ funge da selettore per l'assegnazione a generi esistenti, il termine di transizione è dato da:

\begin{equation}
\Delta P =    \underbrace{(1-p) \frac{(k-1)P_{k-1,t}}{t}}_{\text{Caso 1}} - \underbrace{(1-p) \frac{k P_{k,t}}{t}}_{\text{Caso 2}} + \underbrace{p \delta_{k,1}}_{\text{Caso 3}}
\end{equation}

\noindent Analizziamo i termini: nel Caso 1, dato che il sistema contiene $t$ specie totali al tempo $t$, e assumendo il meccanismo di attaccamento preferenziale, la probabilità che la nuova specie sia assegnata a un genere specifico di dimensione $k-1$ è proporzionale a $(k-1)/t$. Poiché vi sono $P_{k-1,t}$ generi di tale dimensione, il contributo totale è:

\begin{equation}
    \frac{k-1}{t} P_{k-1,t}.
\end{equation}
Applicando il medesimo ragionamento al Caso 2 e considerando il termine sorgente del Caso 3, otteniamo la Master Equation completa:

\begin{equation}
    P_{k,t+1} = P_{k,t} + (1-p)\frac{(k-1)P_{k-1,t}}{t} - (1-p)\frac{k P_{k,t}}{t} + p \delta_{k,1}.
\end{equation}
Possiamo riordinare i termini per evidenziare la differenza finita temporale:

\begin{equation}
    P_{k,t+1} - P_{k,t} = \frac{(1-p)}{t} \left[ (k-1)P_{k-1,t} - k P_{k,t} \right] + p \delta_{k,1}.
\end{equation}

\subsubsection*{Limite del Continuo}

Applichiamo ora l'approssimazione al continuo per grandi $t$ e $k$:

\begin{equation}
    \begin{cases}
    P_{k,t+1} - P_{k,t} \simeq \frac{\partial P_{k}}{\partial t} \\
    P_{k,t} - P_{k-1,t} \simeq \frac{\partial P_{k}}{\partial k}
    \end{cases}
\end{equation}
Si noti che l'approssimazione della differenza finita rispetto a $k$ implica:
\begin{equation}
    \frac{\partial f(x)}{\partial x} \approx f(x) - f(x-1).
\end{equation}
Pertanto, il termine al lato destro dell'equazione può essere riscritto come una derivata rispetto a $k$:

\begin{equation}
    (k-1)P_{k-1,t} - k P_{k,t} \simeq - \frac{\partial (k P_{k,t})}{\partial k}.
\end{equation}
La Master Equation assume quindi la forma di un'equazione differenziale alle derivate parziali (dove la delta di Kronecker è trattata come una delta di Dirac):

\begin{equation}
    \frac{\partial P_{k,t}}{\partial t} = - \frac{(1-p)}{t} \frac{\partial (k P_{k,t})}{\partial k} + p \delta(k-1).
\end{equation}
Ipotizziamo che per $t \to \infty$ il rapporto $P_{k,t}/t$ tenda a una distribuzione stazionaria $q_k$, dipendente unicamente dal numero di specie. Effettuando la sostituzione $P_{k,t} = q_k t$, si ottiene:

\begin{equation}
    q_k = -(1-p)(q_k + k q_k') + p \delta(k-1) \implies (1-p) k q_k' = p \delta(k-1) - (2-p)q_k.
\end{equation}
Per risolvere l'equazione, consideriamo il dominio $k > 1$ (dove la delta è nulla):

\begin{equation}
    \frac{dq_k}{q_k} = - \frac{2-p}{1-p} \frac{dk}{k} \implies q_k = C k^{-\frac{2-p}{1-p}}.
\end{equation}
La costante $C$ viene determinata tramite la condizione di normalizzazione nell'intervallo $[1, \infty)$:

\begin{equation}
    C \int_{1}^{\infty} k^{-\alpha} dk = \left. \frac{C}{1-\alpha} k^{1-\alpha} \right|_{1}^{\infty} = 1.
\end{equation}
Dato che l'esponente $\alpha = \frac{2-p}{1-p}$ è strettamente maggiore di 1, l'integrale converge e si ottiene:

\begin{equation}
    C = \alpha - 1 = \frac{1}{1-p}.
\end{equation}
Pertanto, la distribuzione asintotica per $k \ge 1$ è:

\begin{equation}
    q_k = \frac{1}{1-p} k^{-\frac{2-p}{1-p}}.
\end{equation}
Possiamo esprimere il risultato come una legge di potenza $q_k \propto k^{-\alpha}$, con esponente:

\begin{equation}
    \alpha = \frac{2-p}{1-p} = 1 + \frac{1}{1-p}.
\end{equation}
Richiamando la relazione tra la funzione di densità di probabilità (PDF) e la legge di Zipf $f(R) \sim R^{-\gamma}$, con $\gamma = 1/(\alpha-1)$, otteniamo:

\begin{equation}
    \gamma = 1-p.
\end{equation}
Se la probabilità di introdurre un nuovo genere è molto piccola ($p \ll 1$), allora $\gamma \sim 1$. Questo risultato è coerente con le osservazioni empiriche di Zipf sui testi, dimostrando che il processo di Yule-Simon nello stato stazionario riproduce la legge di Zipf.

\subsection{Approccio Alternativo}

Esiste un metodo alternativo per giungere alla medesima conclusione, ponendosi in un contesto più realistico.
Sia $n_i(t)$ il numero di copie della $i$-esima parola al tempo $t$. Definiamo $p$ come la probabilità di generare una parola inedita e $1-p$ come la probabilità di generare una parola già osservata.
La frequenza relativa della parola $i$ al tempo $t$ è:

\begin{equation}
    f_i(t) = \frac{n_i(t)}{t}.
\end{equation}
Assumiamo l'\textit{ansatz} secondo cui il tasso di variazione del numero di occorrenze sia proporzionale alla frequenza stessa:

\begin{equation}
    \frac{d}{dt}n_i(t) = (1-p)\frac{n_i(t)}{t}.
\end{equation}
Integrando tra il tempo di prima apparizione $t_i$ (dove $n_i(t_i)=1$) e il tempo attuale $t$:

\begin{equation}
    \ln\left(\frac{n_i(t)}{1}\right) = (1-p) \ln\left(\frac{t}{t_i}\right) \implies n_i(t) = \left( \frac{t}{t_i} \right)^{1-p}.
\end{equation}
Per determinare la distribuzione di probabilità di $n_i$, passiamo attraverso la distribuzione cumulativa. Calcoliamo la probabilità che $n_i$ sia minore o uguale a un valore $n$:

\begin{equation}
    n_i(t) \le n \implies \left(\frac{t}{t_i}\right)^{1-p} \le n \implies t_i \ge t n^{-\frac{1}{1-p}} \equiv \tau.
\end{equation}
La probabilità $P(t_i \ge \tau)$ rappresenta la frazione di parole introdotte dopo il tempo $\tau$. Utilizzando la funzione $D(t)$ (numero di parole distinte) introdotta precedentemente:

\begin{equation}
    P(t_i \ge \tau) \approx 1 - \frac{D(\tau)}{D(t)} = 1 - \frac{N_0 + p\tau}{N_0 + pt}.
\end{equation}
Sostituendo $\tau$ e considerando il limite asintotico per grandi $t$ (trascurando $N_0$):

\begin{equation}
    P(n_i \le n) \approx 1 - \frac{p t n^{-\frac{1}{1-p}}}{p t} = 1 - n^{-\frac{1}{1-p}}.
\end{equation}
Infine, otteniamo la distribuzione di probabilità derivando la cumulativa rispetto a $n$:

\begin{equation}
    P(n) = \frac{d}{dn} P(n_i \le n) = \frac{1}{1-p} n^{-\frac{1}{1-p}-1} \sim n^{-\beta},
\end{equation}
con $\beta = 1 + \frac{1}{1-p}$, confermando esattamente il risultato ottenuto con la Master Equation.

\chapter{Lezione 6}
\label{chap:lezione_06} 

\begin{flushright}
\textit{Data: 16/10/2025}
\end{flushright}

\section{Il Modello di Yule-Simon per Generi e Specie}

In questa sezione analizziamo il modello di Yule-Simon (1925) nella sua accezione biologica originale, distinta dall'applicazione alla Legge di Zipf per i testi vista precedentemente. L'obiettivo del modello, presentato originariamente in un articolo su \textit{Nature} del 1922 relativo alla flora della Nuova Zelanda, è spiegare la distribuzione statistica di generi e specie.

\subsubsection{Definizione e Dinamica del Modello}
Consideriamo un sistema in evoluzione in cui:
\begin{itemize}
    \item Il \textbf{tempo} $n$ è scandito dal numero di \textit{Generi} presenti nel sistema.
    \item Ad ogni passo temporale (da $n$ a $n+1$), avvengono due eventi deterministici:
    \begin{enumerate}
        \item Viene introdotto esattamente \textbf{1 nuovo Genere} (innovazione tassonomica).
        \item Vengono introdotte \textbf{$m$ nuove Specie}.
    \end{enumerate}
\end{itemize}

\noindent Questo è un \textbf{Quenched Model}: il tasso di innovazione non è probabilistico ma strutturale ($p = \frac{1}{m+1}$ costante).
L'elemento stocastico risiede nell'assegnazione delle $m$ nuove specie ai generi esistenti. Questa assegnazione segue la regola dell'\textit{attaccamento preferenziale} (\textbf{Rich-get-Richer}): la probabilità che un genere $i$ acquisisca una nuova specie è proporzionale al numero di specie $K_i$ che già possiede.

\noindent Sia $N_{tot}(n) = n(m+1)$ il numero totale di specie al tempo $n$. La probabilità $P_i$ di assegnare una nuova specie al genere $i$ è:
\begin{equation}
    P_i = \frac{K_i}{\sum_j K_j} = \frac{K_i}{n(m+1)}
\end{equation}
Poiché ad ogni step inseriamo $m$ specie, dobbiamo ripetere questo processo $m$ volte. La probabilità totale che il genere $i$ acquisisca una nuova specie nell'intervallo temporale unitario è dunque:
\begin{equation}
    \mathbb{P}(K_i \to K_i + 1) \approx m \cdot \frac{K_i}{n(m+1)} = \frac{m}{m+1} \frac{K_i}{n}
\end{equation}

\subsection{Master Equation}
Definiamo $P_{k,n}$ come la frazione di generi aventi $k$ specie al tempo $n$. Il numero assoluto di tali generi è $N_{k,n} = n P_{k,n}$.
Scriviamo la Master Equation che governa l'evoluzione di $N_{k,n}$.
Per $k > 1$, la variazione del numero di generi con $k$ specie è data dal flusso in entrata (da generi con $k-1$ specie) meno il flusso in uscita (generi con $k$ specie che ne acquisiscono un'altra):
\begin{equation}
    (n+1)P_{k,n+1} - n P_{k,n} = \frac{m}{m+1} \left[ (k-1)P_{k-1,n} - k P_{k,n} \right]
\end{equation}
A sinistra abbiamo la variazione temporale discreta, a destra i termini di guadagno e perdita.
Per $k=1$, dobbiamo considerare anche il termine sorgente (il nuovo genere introdotto ad ogni step, che nasce con 1 specie):
\begin{equation}
    (n+1)P_{1,n+1} - n P_{1,n} = 1 - \frac{m}{m+1} \cdot 1 \cdot P_{1,n}
\end{equation}
Il termine "$1$" rappresenta il nuovo genere; il termine sottrattivo rappresenta i generi con 1 specie che ne acquisiscono una seconda, diventando generi con 2 specie.

\subsection{Soluzione Asintotica}
Nel limite di tempi lunghi ($n \to \infty$), assumiamo una soluzione stazionaria $P_{k,n} \to P_k$.
L'equazione per $k=1$ diventa:
\begin{equation}
    (m+1)P_1 = m P_1 + 1 - \frac{m}{m+1} P_1
\end{equation}
Isoliamo $P_1$ attraverso i seguenti passaggi algebrici:
\begin{equation}
    P_1 \left( m+1 - m + \frac{m}{m+1} \right) = 1
\end{equation}
\begin{equation}
    P_1 \left( 1 + \frac{m}{m+1} \right) = 1 \implies P_1 \left( \frac{2m+1}{m+1} \right) = 1
\end{equation}
Da cui otteniamo il risultato fondamentale per i generi con una sola specie:
\begin{equation}
    P_1 = \frac{m+1}{2m+1}
\end{equation}
Per il caso generico $k$, l'equazione stazionaria è:
\begin{equation}
    P_k = \frac{m}{m+1} \left[ (k-1)P_{k-1} - k P_k \right]
\end{equation}
Risolviamo per $P_k$:
\begin{equation}
    P_k \left( 1 + \frac{mk}{m+1} \right) = \frac{m}{m+1}(k-1)P_{k-1}
\end{equation}
Sviluppiamo il termine tra parentesi a sinistra:
\begin{equation}
    1 + \frac{mk}{m+1} = \frac{m+1+mk}{m+1} = \frac{m(1 + \frac{1}{m} + k)}{m+1}
\end{equation}
Sostituendo nell'equazione:
\begin{equation}
    P_k \frac{m(k + 1 + 1/m)}{m+1} = \frac{m(k-1)}{m+1} P_{k-1}
\end{equation}
Semplificando il fattore $\frac{m}{m+1}$ da entrambi i lati, otteniamo la relazione ricorsiva pulita:
\begin{equation}
    P_k = \frac{k-1}{k + 1 + \frac{1}{m}} P_{k-1}
\end{equation}

\subsubsection{Svolgimento della Ricorsione e Funzione Beta}
Iterando la relazione ricorsiva fino a $P_1$, otteniamo:
\begin{equation}
    P_k = P_1 \prod_{j=2}^{k} \frac{j-1}{j + 1 + \frac{1}{m}} = P_1 \frac{(k-1)!}{(k+1+\frac{1}{m})(k+\frac{1}{m})\cdots(3+\frac{1}{m})}
\end{equation}
Per esprimere questo prodotto in termini di funzioni Gamma, moltiplichiamo e dividiamo per completare i termini mancanti nel denominatore, specificamente $\Gamma(2 + \frac{1}{m})$:
\begin{equation}
    P_k = P_1 \frac{\Gamma(k) \Gamma(2+\frac{1}{m})}{\Gamma(k + 2 + \frac{1}{m})}
\end{equation}
Ricordando che $P_1 = \frac{m+1}{2m+1}$, che possiamo riscrivere manipolando i termini $\Gamma$, si giunge alla forma compatta in termini della funzione Beta di Eulero $B(a,b) = \frac{\Gamma(a)\Gamma(b)}{\Gamma(a+b)}$:
\begin{equation}
    P_k \propto B\left(k, 2 + \frac{1}{m}\right)
\end{equation}
Utilizzando l'approssimazione di Stirling per $k$ grandi ($B(k, \beta) \sim k^{-\beta}$), otteniamo il comportamento asintotico a legge di potenza:
\begin{equation}
    P_k \sim k^{-\left(2 + \frac{1}{m}\right)}
\end{equation}
L'esponente della distribuzione è quindi $\beta = 2 + \frac{1}{m} = 1 + \frac{1}{1-p}$, confermando il risultato già ottenuto in precedenza.

\section{Random Walks}

Perché analizzare nuovamente i Random Walks (RW) in un corso di Sistemi Complessi, dato che sono stati trattati in corsi precedenti? La motivazione risiede nella ricerca di meccanismi universali in grado di generare \textbf{leggi di potenza} (power laws). Abbiamo visto come il modello di Yule-Simon generi tali leggi attraverso meccanismi di crescita; vedremo ora come anche un semplice Random Walk, pur privo di correlazioni complesse o memoria a lungo termine, produca naturalmente leggi di potenza nella distribuzione dei tempi di primo ritorno.

\subsection{Il Modello 1D a Tempo Discreto}
Consideriamo un camminatore aleatorio unidimensionale.
\begin{itemize}
    \item La variabile $x(t)$ indica la posizione al tempo $t$.
    \item Si parte dall'origine: $x(0) = 0$.
    \item Ad ogni passo temporale, la posizione varia di $\Delta x = \pm 1$ con probabilità simmetrica $p = q = 1/2$.
\end{itemize}

Una proprietà fondamentale della topologia di questo cammino è che, partendo dall'origine, ogni ritorno all'origine ($x(t)=0$) deve necessariamente avvenire a un tempo \textbf{pari} ($t=2m$). Infatti, ogni passo a destra deve essere compensato da un passo a sinistra per tornare a zero, implicando che il numero totale di passi deve essere la somma di due numeri uguali, e quindi pari. Definiamo due quantità probabilistiche distinte:

\begin{enumerate}
    \item \textbf{Probabilità di Ritorno ($u_{2n}$):}
    Indichiamo con $u_{2n}$ la probabilità che il camminatore si trovi all'origine al tempo $t=2n$.
    Questa definizione non specifica se sia la prima volta che il camminatore torna, o se vi siano stati ritorni intermedi.
    Per convenzione e coerenza logica, definiamo:
    \begin{equation}
        u_0 := 1
    \end{equation}
    poiché al tempo $t=0$ il camminatore è certamente nell'origine.

    \item \textbf{Probabilità di Primo Ritorno ($f_{2n}$):}
    Indichiamo con $f_{2n}$ la probabilità che il camminatore ritorni all'origine per la \textit{prima volta assoluta} esattamente al tempo $t=2n$.
    Evidentemente, per $t=0$, non si può parlare di "ritorno", quindi:
    \begin{equation}
        f_0 := 0
    \end{equation}
\end{enumerate}

\subsection{Relazione Ricorsiva tra $u_{2n}$ e $f_{2n}$}

Esiste una relazione causale stretta tra le probabilità di ritorno generico e quelle di primo ritorno. Affinché il camminatore si trovi all'origine al tempo $t=2n$, deve essersi verificato il \textbf{primo ritorno} in un qualche istante intermedio $t=2m$ (con $1 \le m \le n$), e successivamente il camminatore deve aver compiuto un percorso che, partendo da zero, lo ha riportato a zero dopo l'intervallo di tempo rimanente $2n - 2m$.

\noindent Poiché gli eventi di "primo ritorno al tempo $2m$" sono mutuamente esclusivi per diversi valori di $m$, possiamo scrivere la probabilità totale $u_{2n}$ come somma delle probabilità congiunte:
\begin{equation}
    u_{2n} = \sum_{m=1}^{n} \text{Prob}(\text{Primo ritorno a } 2m) \times \text{Prob}(\text{Ritorno a } 2n | \text{Primo ritorno a } 2m)
\end{equation}
Il secondo termine, data l'invarianza temporale del processo (il cammino "ricomincia" ogni volta che tocca lo zero), è semplicemente la probabilità di un ritorno (qualsiasi) in un intervallo di durata $2n-2m$. Otteniamo quindi l'equazione di convoluzione discreta:
\begin{equation} \label{eq:convoluzione}
    u_{2n} = \sum_{m=1}^{n} f_{2m} u_{2n-2m} \quad \quad \quad  \text{, per } n \ge 1
\end{equation}

\begin{itemize}
    \item Se $m=n$, il termine nella somma è $f_{2n} u_0 = f_{2n} \cdot 1$. Questo rappresenta il caso in cui il ritorno a $2n$ è esso stesso il primo ritorno.
    \item Se $m < n$, stiamo considerando traiettorie che hanno toccato lo zero per la prima volta a $2m$ e poi, potenzialmente, altre volte prima di arrivare a $2n$.
\end{itemize}

\subsection{Metodo delle Funzioni Generatrici}

Per risolvere l'equazione ricorsiva (\ref{eq:convoluzione}), che coinvolge una convoluzione, lo strumento matematico più potente è quello delle Funzioni Generatrici. Esse "codificano" l'intera sequenza di probabilità in una serie di potenze.
Definiamo le seguenti funzioni generatrici nella variabile complessa (o formale) $z$:

\begin{equation}
    F(z) := \sum_{m=1}^{\infty} f_{2m} z^m
\end{equation}

\begin{equation}
    U(z) := \sum_{n=0}^{\infty} u_{2n} z^n
\end{equation}
Il nostro obiettivo è trovare una relazione funzionale tra $F(z)$ e $U(z)$ sfruttando la convoluzione.

\subsection{Derivazione della Relazione $U(z) - F(z)$}
Partiamo dalla definizione di $U(z)$ e separiamo il termine $n=0$ (dove $u_0=1$) dalla somma:
\begin{equation}
    U(z) = u_0 + \sum_{n=1}^{\infty} u_{2n} z^n = 1 + \sum_{n=1}^{\infty} u_{2n} z^n
\end{equation}
Sostituiamo nell'argomento della somma l'espressione di $u_{2n}$ data dall'equazione (\ref{eq:convoluzione}):
\begin{equation}
    U(z) = 1 + \sum_{n=1}^{\infty} \left( \sum_{m=1}^{n} f_{2m} u_{2n-2m} \right) z^n
\end{equation}
Possiamo riscrivere $z^n$ come $z^m \cdot z^{n-m}$ per associare le potenze ai rispettivi termini:
\begin{equation}
    U(z) = 1 + \sum_{n=1}^{\infty} \sum_{m=1}^{n} (f_{2m} z^m) (u_{2n-2m} z^{n-m})
\end{equation}

A questo punto dobbiamo invertire l'ordine di sommatoria. Nello spazio bidimensionale $(n, m)$, la doppia somma copre un'area triangolare definita da $1 \le m \le n < \infty$. Possiamo coprire questo triangolo in due modi:
\begin{enumerate}
    \item \textbf{Fissando $n$ (fette verticali):} Si somma su $m$ che va da $1$ a $n$. È quello che stiamo facendo ora.
    \item \textbf{Fissando $m$ (fette orizzontali):} Se fissiamo prima $m$ (che può andare da $1$ a $\infty$), allora $n$ deve partire da $m$ e andare a $\infty$.
\end{enumerate}

Procediamo con il secondo metodo (scambio delle somme):
\begin{equation}
    \sum_{n=1}^{\infty} \sum_{m=1}^{n} [...] = \sum_{m=1}^{\infty} \sum_{n=m}^{\infty} (f_{2m} z^m) (u_{2n-2m} z^{n-m})
\end{equation}
Poiché il termine $f_{2m} z^m$ non dipende dall'indice $n$, possiamo portarlo fuori dalla somma interna:
\begin{equation}
    = \sum_{m=1}^{\infty} f_{2m} z^m \left[ \sum_{n=m}^{\infty} u_{2n-2m} z^{n-m} \right]
\end{equation}
Effettuiamo ora un cambio di variabile nella somma interna. Poniamo $k = n - m$.
Quando $n=m$, $k=0$. Quando $n \to \infty$, $k \to \infty$.
\begin{equation}
    = \sum_{m=1}^{\infty} f_{2m} z^m \left[ \sum_{k=0}^{\infty} u_{2k} z^k \right]
\end{equation}
Riconosciamo immediatamente le definizioni delle funzioni generatrici:
\begin{equation}
    = F(z) \cdot U(z)
\end{equation}
Abbiamo quindi dimostrato la relazione fondamentale:
\begin{equation}
    U(z) = 1 + F(z)U(z)
\end{equation}
Da cui possiamo isolare la funzione incognita $F(z)$:
\begin{equation} \label{eq:F_from_U}
    F(z) = 1 - \frac{1}{U(z)}
\end{equation}
Per calcolare la distribuzione dei primi ritorni $F(z)$, dobbiamo quindi prima calcolare esplicitamente $U(z)$.

\subsection{Calcolo Esplicito di $U(z)$}

La probabilità $u_{2n}$ è determinata dal rapporto tra il numero di percorsi favorevoli e il numero totale di percorsi possibili di lunghezza $2n$.
\begin{itemize}
    \item \textbf{Totale percorsi:} Ogni passo ha 2 possibilità ($+1$ o $-1$). In $2n$ passi, ci sono $2^{2n}$ configurazioni totali.
    \item \textbf{Percorsi favorevoli:} Per tornare all'origine, il numero di passi a destra ($n_{\text{up}}$) deve essere uguale al numero di passi a sinistra ($n_{\text{down}}$). Essendo $n_{\text{up}} + n_{\text{down}} = 2n$, deve essere $n_{\text{up}} = n$.
    Il numero di modi per scegliere quali $n$ passi siano "up" su $2n$ totali è dato dal coefficiente binomiale $\binom{2n}{n}$.
\end{itemize}
Quindi:
\begin{equation}
    u_{2n} = \binom{2n}{n} \left(\frac{1}{2}\right)^{2n} = \binom{2n}{n} \frac{1}{4^n}
\end{equation}

\textit{Digressione asintotica:} Usando l'approssimazione di Stirling ($n! \approx \sqrt{2\pi n} (n/e)^n$), si può mostrare che $u_{2n} \sim \frac{1}{\sqrt{\pi n}}$, ovvero decade come una legge di potenza $n^{-1/2}$. Tuttavia, per il calcolo della funzione generatrice useremo l'espressione esatta.

\hfill

\noindent Sostituiamo l'espressione combinatoria nella definizione di $U(z)$:
\begin{equation}
    U(z) = \sum_{n=0}^{\infty} \binom{2n}{n} \frac{1}{4^n} z^n = \sum_{n=0}^{\infty} \binom{2n}{n} \left(\frac{z}{4}\right)^n
\end{equation}
Vogliamo dimostrare che questa serie converge alla funzione algebrica:
\begin{equation}
    U(z) = \frac{1}{\sqrt{1-z}} 
\end{equation}
Per fare ciò, dobbiamo espandere in serie di Taylor la funzione $(1-z)^{-1/2}$ utilizzando il Teorema Binomiale Generalizzato (Binomio di Newton esteso agli esponenti reali negativi).
L'espansione è data da:
\begin{equation}
    (1-z)^{-1/2} = \sum_{k=0}^{\infty} \binom{-1/2}{k} (-z)^k = \sum_{k=0}^{\infty} \binom{-1/2}{k} (-1)^k z^k
\end{equation}
Dobbiamo procedere passo dopo passo per ricondurre il coefficiente binomiale $\binom{-1/2}{k}$ a una forma nota.


\subsubsection{Sviluppo del Coefficiente Binomiale Frazionario}
Per definizione generalizzata del coefficiente binomiale:
\begin{equation}
    \binom{\alpha}{k} = \frac{\alpha (\alpha-1) (\alpha-2) \dots (\alpha-k+1)}{k!}
\end{equation}
Ponendo $\alpha = -1/2$, il numeratore diventa un prodotto di $k$ termini:
\begin{equation}
    \text{Num} = \left(-\frac{1}{2}\right) \left(-\frac{3}{2}\right) \left(-\frac{5}{2}\right) \dots \left(-\frac{1}{2} - k + 1\right)
\end{equation}
L'ultimo termine della produttoria può essere scritto come:
\begin{equation}
    -\frac{1}{2} - k + 1 = -\frac{1}{2} - \frac{2k}{2} + \frac{2}{2} = -\frac{2k-1}{2}
\end{equation}
Quindi il numeratore è:
\begin{equation}
    \text{Num} = \left(-\frac{1}{2}\right) \left(-\frac{3}{2}\right) \left(-\frac{5}{2}\right) \dots \left(-\frac{2k-1}{2}\right)
\end{equation}

\textbf{Passo 1: Estrazione del segno e delle potenze di 2.}
Poiché ci sono $k$ termini nel prodotto, possiamo fattorizzare $(-1/2)$ da ciascuno di essi:
\begin{equation}
    \text{Num} = \left(-\frac{1}{2}\right)^k \underbrace{\left[ 1 \cdot 3 \cdot 5 \cdot \dots \cdot (2k-1) \right]}_{\text{Prodotto dei numeri dispari}}
\end{equation}

\textbf{Passo 2: Completamento del Fattoriale.}
Il termine tra parentesi quadre è il \textit{semifattoriale} dispari $(2k-1)!!$. Non sappiamo calcolarlo direttamente con i fattoriali standard.
L'artificio matematico necessario consiste nel moltiplicare e dividere per il prodotto di tutti i numeri \textbf{pari} mancanti fino a $2k$:
\begin{equation}
    [1 \cdot 3 \cdot \dots \cdot (2k-1)] \cdot \frac{2 \cdot 4 \cdot 6 \cdot \dots \cdot 2k}{2 \cdot 4 \cdot 6 \cdot \dots \cdot 2k}
\end{equation}
Il numeratore ora contiene tutti gli interi da 1 a $2k$, quindi diventa semplicemente $(2k)!$
Analizziamo il denominatore (prodotto dei pari). Possiamo estrarre un fattore 2 da ciascuno dei $k$ termini:
\begin{equation}
    \text{Denom Pari} = (2\cdot 1) \cdot (2\cdot 2) \cdot (2\cdot 3) \cdot \dots \cdot (2\cdot k) = 2^k (1 \cdot 2 \cdot 3 \cdot \dots \cdot k) = 2^k k!
\end{equation}

Sostituendo queste espressioni nel numeratore originale del binomiale:
\begin{equation}
    \text{Num del binomiale} = \left(-\frac{1}{2}\right)^k \frac{(2k)!}{2^k k!} = \frac{(-1)^k}{2^k} \frac{(2k)!}{2^k k!} = \frac{(-1)^k}{2^{2k}} \frac{(2k)!}{k!}
\end{equation}

\textbf{Passo 3: Ricostruzione del coefficiente finale.}
Ora reinseriamo questo risultato nella definizione completa di $\binom{-1/2}{k}$, dividendo per il $k!$ che era al denominatore fin dall'inizio:
\begin{equation}
    \binom{-1/2}{k} = \frac{\text{Num}}{k!} = \frac{1}{k!} \left[ \frac{(-1)^k}{4^k} \frac{(2k)!}{k!} \right] = \frac{(-1)^k}{4^k} \frac{(2k)!}{k! k!}
\end{equation}
Riconoscendo che $\frac{(2k)!}{k! k!} = \binom{2k}{k}$, otteniamo il risultato fondamentale:
\begin{equation}
    \binom{-1/2}{k} = \left(-\frac{1}{4}\right)^k \binom{2k}{k}
\end{equation}

\hfill

\noindent Sostituiamo il coefficiente appena calcolato nell'espansione binomiale di $(1-z)^{-1/2}$:
\begin{equation}
    (1-z)^{-1/2} = \sum_{k=0}^{\infty} \left[ \left(-\frac{1}{4}\right)^k \binom{2k}{k} \right] (-1)^k z^k
\end{equation}
I termini $(-1)^k$ e il segno meno dentro $(-1/4)^k$ si cancellano (poiché $(-1)^{2k} = 1$):
\begin{equation}
    = \sum_{k=0}^{\infty} \binom{2k}{k} \left(\frac{1}{4}\right)^k z^k = \sum_{k=0}^{\infty} \binom{2k}{k} \left(\frac{z}{4}\right)^k
\end{equation}
Questa serie è identica all'espressione di $U(z)$ che avevamo derivato combinatoricamente.
Dunque, abbiamo dimostrato che:
\begin{equation}
    U(z) = \frac{1}{\sqrt{1-z}}
\end{equation}

\subsection{Risultato Finale e Legge di Potenza}

Ora possiamo tornare all'equazione (\ref{eq:F_from_U}) per trovare la funzione generatrice dei primi ritorni:
\begin{equation}
    F(z) = 1 - \frac{1}{U(z)} = 1 - \sqrt{1-z}
\end{equation}

Sviluppando in serie di Taylor il termine $\sqrt{1-z}$ (analogamente a quanto fatto prima, ma con esponente $+1/2$), si ottengono i coefficienti $f_{2n}$.
L'analisi asintotica di questi coefficienti per $n$ grandi mostra che:
\begin{equation}
    f_{2n} \sim n^{-3/2}
\end{equation}
Abbiamo ottenuto che la distribuzione dei tempi di primo ritorno per un Random Walk 1D decade come una legge di potenza con esponente $\alpha = 3/2$.
Questa è nota come \textbf{Distribuzione di Levy}.
Il risultato è significativo perché dimostra che non è necessario introdurre dinamiche complesse o correlazioni strutturate per ottenere comportamenti di scaling (legge di potenza): anche un processo stocastico semplice e privo di memoria come il Random Walk genera naturalmente invarianza di scala nella statistica dei ritorni. Il RW funge quindi da modello fondamentale per i Sistemi Complessi.

\section{Leggi di Potenza in Fenomeni Critici}
I fenomeni critici rappresentano un esempio paradigmatico dell'emersione di invarianza di scala (\textit{scale invariance}) e leggi di potenza nei sistemi fisici. Un esempio classico è la transizione di fase paramagnetica-ferromagnetica.
Sia $m$ il parametro d'ordine (magnetizzazione).
\begin{itemize}
    \item Per $T > T_c$ (alta temperatura): il sistema è disordinato, $m=0$.
    \item Per $T < T_c$: emerge un ordine spontaneo, $m \neq 0$.
\end{itemize}
In prossimità della temperatura critica $T_c$, le grandezze fisiche non scalano in modo esponenziale, ma secondo leggi di potenza. Ad esempio, la funzione di correlazione spaziale delle fluttuazioni decade come:
\begin{equation}
    G(r) \sim \frac{1}{|r|^{d-2+\eta}}
\end{equation}
La lunghezza di correlazione $\xi$, che rappresenta la scala caratteristica delle fluttuazioni connesse, diverge al punto critico:
\begin{equation}
    \xi \sim |T - T_c|^{-\nu}
\end{equation}
Questa divergenza implica che il sistema diventa invariante per cambio di scala: non esiste più una lunghezza caratteristica che domini la fisica del sistema. Kenneth Wilson (Nobel 1982) formalizzò l'approccio per trattare questi sistemi attraverso il \textbf{Gruppo di Rinormalizzazione} (RG).

\subsection{Percolazione e Ipotesi di Scaling}

Analizziamo il fenomeno della percolazione come modello geometrico di transizione di fase. Consideriamo un reticolo quadrato 2D dove ogni sito è occupato con probabilità $p$. Esiste una probabilità critica $p_c$ (per il reticolo quadrato $p_c \approx 0.5927$) sopra la quale compare un cluster gigante che attraversa l'intero sistema.
Al punto critico $p_c$, la struttura dei cluster è frattale. Cerchiamo di determinare la forma della distribuzione delle dimensioni dei cluster, $P(s)$, definita come la probabilità che un sito scelto a caso appartenga a un cluster di area $s$.

\subsubsection{Analisi Dimensionale e Scaling}
Cerchiamo una relazione funzionale per $P(s)$. Supponiamo che essa dipenda dalle grandezze caratteristiche del sistema:
\begin{itemize}
    \item $s$: dimensione (area) del cluster in esame.
    \item $a$: area della cella elementare del reticolo (cutoff ultravioletto).
    \item $\langle s \rangle$: dimensione media dei cluster (che diverge a $p_c$).
\end{itemize}
Per costruire una funzione adimensionale, dobbiamo combinare queste grandezze. I possibili gruppi adimensionali sono $s/a$ e $s/\langle s \rangle$.
Assumiamo una forma di scaling generalizzata:
\begin{equation}
    P(s) = C s^{-\tau} f\left( \frac{s}{a}, \frac{a}{\langle s \rangle} \right)
\end{equation}
dove $\tau$ è un esponente critico (Fisher exponent) e $f$ è una funzione di scaling universale. La condizione di normalizzazione richiede $\sum_s P(s) = 1$.


\noindent Effettuiamo una trasformazione di scala spaziale, riducendo la risoluzione del sistema.
Riscaliamo l'unità di lunghezza (o area) di un fattore $\lambda$:
\begin{equation}
    a \to a' = \frac{a}{\lambda}
\end{equation}
Sotto questa trasformazione, ci aspettiamo che la distribuzione mantenga la stessa forma funzionale, ma con parametri rinormalizzati:
\begin{equation}
    P(s) = C' s^{-\tau} f\left( \frac{s}{a/\lambda}, \frac{a/\lambda}{\langle s \rangle'} \right) = C' s^{-\tau} f\left( \frac{s \lambda}{a}, \frac{a}{\lambda \langle s \rangle'} \right)
\end{equation}
Se osserviamo il sistema a una scala più grossolana (zoom out, quindi cella elementare effettiva più grande), la dimensione media dei cluster misurata in unità del reticolo deve cambiare.

\subsubsection{Condizione di Punto Critico}
Esiste una sola situazione in cui la dimensione media dei cluster rimane invariante sotto riscalamento (a parte fattori di scala banali), ovvero quando essa è infinita:
\begin{equation}
    \langle s \rangle \to \infty
\end{equation}
In questo regime (punto critico), il termine $a/\langle s \rangle$ tende a 0 e scompare dalla dipendenza della funzione $f$. La distribuzione diventa una pura legge di potenza (Power Law):
\begin{equation}
    P(s) \sim s^{-\tau}
\end{equation}
In tutti gli altri casi (fuori dal punto critico), la funzione di cutoff $f$ introduce una scala caratteristica che rompe l'invarianza di scala (spesso un cutoff esponenziale $e^{-s/s_{\xi}}$).
L'invarianza di scala, espressa dalla relazione $P(\lambda s) = g(\lambda) P(s)$, è possibile solo ed esclusivamente al punto critico.

\subsection{Rinormalizzazione nello Spazio Reale}

Applichiamo ora concretamente il Gruppo di Rinormalizzazione nello spazio reale (\textit{Real Space Renormalization Group}) a un sistema 1D.
Immaginiamo una catena di siti, ciascuno occupato (nero) con probabilità $p$ e vuoto (bianco) con probabilità $1-p$.

\subsubsection{La Trasformazione $\mathcal{R}_2(p)$}
Definiamo una trasformazione a blocchi (\textbf{Block Spin transformation}) con fattore di scala $b=2$:
\begin{enumerate}
    \item Raggruppiamo i siti a coppie vicine (blocchi di 2).
    \item Definiamo lo stato del blocco rinormalizzato (super-sito) basandoci su una regola deterministica.
    \item \textbf{Regola:} Il blocco è considerato "nero" (occupato) se e solo se \textbf{entrambi} i siti costituenti sono neri. Altrimenti è bianco.
\end{enumerate}
Questa regola è dettata dalla necessità di preservare la connettività: in 1D, basta un solo sito vuoto per interrompere il flusso.
La probabilità $p'$ che il blocco rinormalizzato sia nero è data dalla probabilità congiunta che i due siti originali siano neri. Essendo eventi indipendenti:
\begin{equation}
    p' = \mathcal{R}_2(p) = p \cdot p = p^2
\end{equation}
Questa è la nostra \textbf{Mappa di Rinormalizzazione}.

\subsubsection{Analisi dei Punti Fissi}
I punti fissi $p^*$ della trasformazione sono le soluzioni di $p^* = \mathcal{R}_2(p^*)$:
\begin{equation}
    p^* = (p^*)^2 \implies p^*(1 - p^*) = 0
\end{equation}
Abbiamo due punti fissi: $p^* = 0$ e $p^* = 1$.

\noindent Per determinare la natura dei punti fissi, linearizziamo la mappa nell'intorno di $p^*$. Sia $p_n$ il valore alla $n$-esima iterazione e consideriamo una piccola perturbazione $\epsilon_n$:
\begin{equation}
    p_n = p^* + \epsilon_n
\end{equation}
Sviluppando al primo ordine:
\begin{equation}
    p_{n+1} = \mathcal{R}(p_n) \approx \mathcal{R}(p^*) + \mathcal{R}'(p^*) \epsilon_n = p^* + \lambda \epsilon_n
\end{equation}
dove l'autovalore di stabilità è $\lambda = \frac{d\mathcal{R}}{dp}\big|_{p^*}$.
Calcoliamo la derivata:
\begin{equation}
    \frac{d}{dp}(p^2) = 2p
\end{equation}
Valutiamo la stabilità:
\begin{itemize}
    \item Per $p^* = 0$: $\lambda = 2(0) = 0$. Poiché $|\lambda| < 1$, il punto fisso è \textbf{Stabile}. Qualsiasi $p < 1$ tenderà a 0 iterando la mappa.
    \item Per $p^* = 1$: $\lambda = 2(1) = 2$. Poiché $|\lambda| > 1$, il punto fisso è \textbf{Instabile}.
\end{itemize}

\textbf{Interpretazione Fisica:} In 1D, il sistema è percolante solo se $p=1$. Qualsiasi imperfezione ($p < 1$), se vista su scala sufficientemente grande (dopo molte rinormalizzazioni), interrompe la connessione. Il sistema fluisce verso la fase disordinata ($p=0$). Il punto $p=1$ è il punto critico "banale".

\subsection{Percolazione 2D}

Consideriamo ora un reticolo triangolare 2D. Questo sistema presenta una transizione di fase non banale.
Eseguiamo una rinormalizzazione raggruppando 3 siti (un triangolo elementare) in un singolo super-sito (blocco).
\begin{itemize}
    \item \textbf{Fattore di scala spaziale ($b$):} La geometria del reticolo implica che se il lato del triangolo piccolo è $1$, il lato del triangolo che connette i centri dei blocchi adiacenti è $\sqrt{3}$. Quindi $b = \sqrt{3}$.
    \item \textbf{Regola di Rinormalizzazione (Majority Rule):} Definiamo il blocco come "nero" se la maggioranza dei suoi siti (2 o 3) è nera.
\end{itemize}

Calcoliamo la probabilità rinormalizzata $p'$ sommando le probabilità delle configurazioni che soddisfano la regola.
Le configurazioni possibili di 3 siti sono $2^3 = 8$. Quelle vincenti sono:
\begin{enumerate}
    \item \textbf{3 siti neri:} C'è 1 sola configurazione (N-N-N). Probabilità: $p^3$.
    \item \textbf{2 siti neri e 1 bianco:} Ci sono $\binom{3}{2} = 3$ configurazioni possibili (N-N-B, N-B-N, B-N-N). Ciascuna ha probabilità $p^2(1-p)$. Il contributo totale è $3p^2(1-p)$.
\end{enumerate}

La mappa di rinormalizzazione è dunque:
\begin{equation}
    p' = R(p) = p^3 + 3p^2(1-p)
\end{equation}
Sviluppando l'algebra:
\begin{equation}
    p' = p^3 + 3p^2 - 3p^3 = 3p^2 - 2p^3
\end{equation}

\subsubsection{Punti Fissi e Criticità}
Risolviamo l'equazione di punto fisso $p = R(p)$:
\begin{equation}
    3p^2 - 2p^3 = p
\end{equation}
Portiamo tutto a sinistra e fattorizziamo:
\begin{equation}
    2p^3 - 3p^2 + p = 0 \implies p(2p^2 - 3p + 1) = 0
\end{equation}
Risolvendo l'equazione quadratica $2p^2 - 3p + 1 = 0$:
\begin{equation}
    p_{1,2} = \frac{3 \pm \sqrt{9 - 8}}{4} = \frac{3 \pm 1}{4} \implies p \in \{1, \frac{1}{2}\}
\end{equation}
I punti fissi sono quindi tre:
\begin{itemize}
    \item $p^* = 0$: Fase vuota (Stabile).
    \item $p^* = 1$: Fase piena (Stabile).
    \item $p^* = 1/2$: \textbf{Punto Critico $p_c$} (Instabile).
\end{itemize}
Questo modello prevede correttamente l'esistenza di una soglia di percolazione non banale a $p_c = 0.5$.

\subsubsection{Esponente Critico}
Dalla teoria dello scaling, la lunghezza di correlazione diverge vicino al punto critico come:
\begin{equation}
    \xi \sim |p - p_c|^{-\nu}
\end{equation}
Vogliamo calcolare l'esponente critico $\nu$ usando la mappa RG.
Consideriamo la trasformazione della lunghezza di correlazione e della "distanza dal punto critico" $\delta p = p - p_c$ sotto un passo di rinormalizzazione.

\begin{enumerate}
    \item \textbf{Trasformazione spaziale:} Poiché abbiamo riscalato il sistema di un fattore $b = \sqrt{3}$, la nuova lunghezza di correlazione $\xi'$ (misurata nelle nuove unità) sarà più piccola:
    \begin{equation} \label{eq:xi_scaling}
        \xi' = \frac{\xi}{b}
    \end{equation}

    \item \textbf{Trasformazione del parametro:} Linearizziamo la mappa $R(p)$ attorno al punto instabile $p_c = 1/2$:
    \begin{equation}
        p' - p_c \approx \Lambda (p - p_c) \quad \text{con} \quad \Lambda = \frac{dR}{dp}\bigg|_{p_c}
    \end{equation}
    Calcoliamo la derivata (il Jacobiano 1D):
    \begin{equation}
        R'(p) = \frac{d}{dp}(3p^2 - 2p^3) = 6p - 6p^2 = 6p(1-p)
    \end{equation}
    Valutata nel punto critico $p_c = 1/2$:
    \begin{equation}
        \Lambda = R'(1/2) = 6 \cdot \frac{1}{2} \cdot \left(1 - \frac{1}{2}\right) = 6 \cdot \frac{1}{4} = \frac{3}{2}
    \end{equation}
    Quindi, ad ogni passo, la distanza dal punto critico viene amplificata di un fattore $\Lambda = 1.5$.
\end{enumerate}
Ora colleghiamo le due trasformazioni usando la definizione di scaling $\xi \sim (\delta p)^{-\nu}$.
Per il sistema rinormalizzato deve valere la stessa legge:
\begin{equation}
    \xi' \sim (\delta p')^{-\nu}
\end{equation}
Sostituendo le relazioni trovate ($\xi' = \xi/b$ e $\delta p' = \Lambda \delta p$):
\begin{equation}
    \frac{\xi}{b} \sim (\Lambda \delta p)^{-\nu} = \Lambda^{-\nu} (\delta p)^{-\nu}
\end{equation}
Poiché $(\delta p)^{-\nu} \sim \xi$, otteniamo:
\begin{equation}
    \frac{\xi}{b} = \Lambda^{-\nu} \xi \implies b = \Lambda^{\nu}
\end{equation}
Estraendo $\nu$ mediante i logaritmi:
\begin{equation}
    \nu = \frac{\ln b}{\ln \Lambda}
\end{equation}
Sostituendo i valori numerici del nostro modello ($b=\sqrt{3}$, $\Lambda=3/2$):
\begin{equation}
    \nu = \frac{\ln \sqrt{3}}{\ln (3/2)} = \frac{\frac{1}{2} \ln 3}{\ln 3 - \ln 2} \approx 1.35
\end{equation}
Il valore esatto dell'esponente $\nu$ per la percolazione 2D è $4/3 \approx 1.33$. Il nostro semplice modello di rinormalizzazione nello spazio reale fornisce un risultato sorprendentemente accurato (errore $\sim 1.5\%$), catturando l'essenza della fisica dei fenomeni critici senza calcoli complessi.
\chapter{Lezione 7}
\label{chap:lezione_07} 

\begin{flushright}
\textit{Data: 22/10/2025}
\end{flushright}

\section{Self-Organized Criticality e il Modello Sandpile}

In questa lezione analizziamo la \textit{Self-Organized Criticality} (SOC), un concetto fondamentale per comprendere come i sistemi complessi possano evolvere spontaneamente verso uno stato critico senza la necessità di sintonizzare parametri esterni. Questo approccio è utile per esplorare i metodi di calcolo degli esponenti critici. Il modello di riferimento è il \textbf{Sandpile Model} (o modello della pila di sabbia), introdotto nel 1987 come un "toy model" ispirato alle pile di sabbia reali.

\subsubsection{Dinamica del Modello}
Consideriamo una topologia specifica, in questo caso un reticolo rettangolare bidimensionale (2D). La dinamica è definita dalle seguenti regole:
\begin{enumerate}
    \item \textbf{Iniezione di energia:} Si lascia cadere "energia" (rappresentata da cubi o granelli) sul reticolo.
    \item \textbf{Instabilità e Valanghe:} Quando in un sito si accumulano 4 cubi, il sito diventa instabile. L'instabilità provoca una ridistribuzione dell'energia ai siti vicini.
    \item Se un sito possiede 3 cubi, è definito "instabile" nel senso che l'aggiunta di un solo ulteriore cubo innesca la valanga.
\end{enumerate}

Se si continua a iniettare energia nel sistema, questo raggiungerà asintoticamente uno \textbf{Stato Stazionario}. In questo stato, le valanghe (definite come la sequenza di eventi di ridistribuzione) seguono una distribuzione a legge di potenza (power-law) per dimensioni infinite. Possiamo misurare diverse osservabili fisiche relative alle valanghe:
\begin{itemize}
    \item La dimensione lineare o diametro ($r$).
    \item La dimensione totale ($s$, numero di siti coinvolti o granelli spostati).
    \item La durata temporale ($t$).
\end{itemize}
Sperimentalmente si osserva che le distribuzioni di probabilità di queste grandezze sono tutte leggi di potenza.
In questo contesto, distinguiamo tra:
\begin{itemize}
    \item \textit{Fast modes}: i modi veloci legati alla dinamica delle valanghe.
    \item \textit{Slow modes}: i modi lenti legati all'iniezione di energia dall'esterno.
\end{itemize}

\subsection{Relazioni di Scaling e Esponenti Critici}

Definiamo le distribuzioni di probabilità e le relazioni di scala tra le variabili. Dagli esperimenti e dalle simulazioni emerge che:

\begin{align}
    P(s) &\sim s^{-\tau} \label{eq:Ps} \\
    P(r) &\sim r^{-\lambda} \label{eq:Pr} \\
    P(t) &\sim t^{-\alpha} \label{eq:Pt}
\end{align}

Le variabili stesse sono collegate da relazioni di scala:
\begin{align}
    t &\sim s^x \label{eq:t_s} \\
    t &\sim r^z \label{eq:t_r} \\
    s &\sim r^D \label{eq:s_r}
\end{align}

Dove $z$ è definito come l'\textbf{esponente dinamico}.
Non tutti questi esponenti sono indipendenti. L'obiettivo è calcolare gli esponenti critici derivandoli l'uno dall'altro e risolvendo il modello. Per fare ciò, utilizzeremo il \textbf{Gruppo di Rinormalizzazione nello Spazio Reale}, analizzando come cambia il sistema al variare della scala.

\subsubsection{Derivazione delle Relazioni tra Esponenti}

Cerchiamo di correlare l'esponente $\tau$ con gli esponenti $D$ e $\lambda$.
Consideriamo la conservazione della probabilità tra le distribuzioni di dimensione $s$ e dimensione lineare $r$:
\begin{equation}
    P(s) ds = P(r) dr
\end{equation}
Dalla relazione $s \sim r^D$, differenziando rispetto a $r$ otteniamo:
\begin{equation}
    \frac{ds}{dr} \sim \frac{d}{dr}(r^D) \sim r^{D-1} \implies ds \sim r^{D-1} dr
\end{equation}
Sostituendo le leggi di potenza (\ref{eq:Ps}) e (\ref{eq:Pr}) nell'uguaglianza delle probabilità:
\begin{equation}
    s^{-\tau} ds = r^{-\lambda} dr
\end{equation}
Esprimiamo $s$ in termini di $r$ nel lato sinistro:
\begin{equation}
    (r^D)^{-\tau} \cdot r^{D-1} dr \sim r^{-\lambda} dr
\end{equation}
Semplificando le potenze di $r$:
\begin{equation}
    r^{-D\tau} \cdot r^{D-1} \sim r^{-\lambda} \implies r^{-D\tau + D - 1} \sim r^{-\lambda}
\end{equation}
Uguagliando gli esponenti:
\begin{equation}
    -D\tau + D - 1 = -\lambda
\end{equation}
Moltiplicando per $-1$ e riordinando, otteniamo la prima relazione fondamentale:
\begin{equation}
    \lambda = D(\tau - 1) + 1
\end{equation}

Ora correliamo gli esponenti temporali. Sappiamo che $t \sim s^x$. Possiamo anche scrivere $t \sim r^z$ e, dato che $s \sim r^D$, segue che:
\begin{equation}
    t \sim (r^D)^x = r^{Dx}
\end{equation}
Confrontando con $t \sim r^z$, otteniamo la relazione tra $z$, $D$ e $x$:
\begin{equation}
    z = Dx \implies x = \frac{z}{D}
\end{equation}
Usiamo la conservazione della probabilità tra $s$ e $t$:
\begin{equation}
    P(s) ds = P(t) dt
\end{equation}
Sostituendo le distribuzioni note:
\begin{equation}
    s^{-\tau} ds = t^{-\alpha} dt
\end{equation}
Dalla relazione $t \sim s^x$, differenziando otteniamo $dt \sim x s^{x-1} ds$. Sostituendo:
\begin{equation}
    s^{-\tau} ds \sim t^{-\alpha} (s^{x-1} ds)
\end{equation}
Esprimiamo $t$ in funzione di $s$ ($t \sim s^x$):
\begin{equation}
    s^{-\tau} \sim (s^x)^{-\alpha} \cdot s^{x-1} = s^{-\alpha x} \cdot s^{x-1} = s^{-\alpha x + x - 1}
\end{equation}
Uguagliando gli esponenti di $s$:
\begin{equation}
    -\tau = -x\alpha + x - 1
\end{equation}
Isoliamo $\alpha$:
\begin{equation}
    x\alpha = \tau + x - 1 \implies \alpha = \frac{\tau - 1}{x} + 1
\end{equation}
Sostituendo $x = z/D$:
\begin{equation}
    \alpha = 1 + \frac{D}{z}(\tau - 1)
\end{equation}
Abbiamo quindi un set di esponenti $\tau, z, D, \lambda, \alpha$.
L'analisi dimensionale suggerirebbe che per $d=2$ si abbia $D=2$, ma i sistemi frattali mostrano deviazioni. Se conosciamo $D, \tau, z$, possiamo derivare tutti gli altri.
Valori sperimentali o numerici tipici sono:
\begin{itemize}
    \item $D \simeq 2.75 \pm 0.05$.
    \item In altri casi $D \simeq 3.25 \pm 0.25$.
\end{itemize}
Si tratta di esponenti frattali (non interi). Il nostro obiettivo è trovare un metodo analitico per calcolarli.

\subsection{Procedura di Coarse-Graining}
Applichiamo il metodo del \textit{Coarse-Graining} (granulazione grossolana). Guardiamo il sistema a due scale differenti: la scala $k$ e la scala $k+1$.
L'idea è raggruppare insieme porzioni del sistema ("patches") per formare unità più grandi. Consideriamo un blocco di $2 \times 2$ siti alla scala $k$. Questi 4 siti diventano una singola unità alla scala $k+1$. Vogliamo trovare una dinamica invariante di scala (scale-invariant dynamics).
Si tratta di un processo combinatorio.

Definiamo il vettore delle probabilità alla scala microscopica più bassa ($k=0$):
\begin{equation}
    \vec{P}^{(0)} = (0, 0, 0, 1)
\end{equation}
Le componenti di questo vettore rappresentano la probabilità di scaricare energia su 1, 2, 3 o 4 siti vicini. Alla scala microscopica, in un reticolo quadrato, un sito scarica sempre su tutti e 4 i vicini, quindi $P_4^{(0)} = 1$ e gli altri sono nulli.

Passando alla scala $k$, avremo un vettore generico $\vec{P}^{(k)}$ con norma unitaria $\sum P_i^{(k)} = 1$.
Nel passaggio di coarse-graining alla scala $k+1$, definiamo $\rho^{(k+1)}$ come la densità dei siti critici (density of critical state).
Fisicamente, $\rho$ rappresenta la densità delle pile che stanno per crollare se viene aggiunto un altro granello.

\subsubsection{Equazioni di Rinormalizzazione Autoconsistenti}
Dobbiamo calcolare le probabilità di trasferimento alla scala successiva. Definiamo $W_\alpha$ come la probabilità di avere una configurazione critica su un blocco di dimensione maggiore. Le probabilità per le configurazioni critiche sono date da:
\begin{align}
    W_2 &= 4\rho^2(1-\rho)^2 \\
    W_3 &= 4\rho^3(1-\rho) \\
    W_4 &= \rho^4
\end{align}
Queste rappresentano le probabilità combinatorie che 2, 3 o 4 siti nel blocco siano critici (assumendo indipendenza in campo medio). Nello Stato Stazionario, la quantità di energia che entra deve essere uguale a quella che esce in media. L'energia è costante in media:
\begin{equation}
    \delta E^{(k+1)} = \rho^{(k+1)}
\end{equation}
L'energia trasferita a seguito di una valanga alla scala $k+1$ è proporzionale a $\rho^{(k+1)}$.
L'equazione di bilancio (punto fisso) che lega la densità critica alle probabilità di trasferimento è data da:
\begin{equation}
    \rho^{(k+1)} = \frac{1}{\sum_{n=1}^4 n P_n^{(k+1)}} = \frac{1}{P_1^{(k+1)} + 2P_2^{(k+1)} + 3P_3^{(k+1)} + 4P_4^{(k+1)}}
\end{equation}
Le probabilità $P_n^{(k+1)}$ sono a loro volta determinate da una somma pesata sui $W_\alpha$:
\begin{equation}
    P_n^{(k+1)} = \sum_{\alpha \in \{2,3,4\}} W_\alpha P_n^{(k+1)}(\alpha)
\end{equation}
Dove $P_n^{(k+1)}(\alpha)$ è la probabilità che il super-sito scarichi su $n$ vicini dato che internamente ha $\alpha$ sottositi critici.
Come indicato nelle note, il calcolo esplicito di questi termini coinvolge coefficienti binomiali (es. $\binom{4}{2}$ per scegliere quali siti sono critici).

\subsubsection{Calcolo dell'esponente $\tau$}
Definiamo la \textbf{Stop Probability} $R$ come la probabilità che, a una qualsiasi scala, la valanga si fermi (non propaghi alla scala successiva).
La valanga si ferma se, trasferendo energia su $n$ siti, nessuno di questi è critico. Quindi:
\begin{equation}
    R = P_1^* (1 - \rho^*) + P_2^* (1 - \rho^*)^2 + P_3^* (1 - \rho^*)^3 + P_4^* (1 - \rho^*)^4
\end{equation}
dove l'asterisco indica i valori al punto fisso della rinormalizzazione.

Possiamo collegare $R$ alla distribuzione delle valanghe. La probabilità che una valanga si fermi tra la scala $a(k)$ e la scala $a(k+1)$ è data dall'integrale della distribuzione $P(r)$ in quell'intervallo.
Definiamo la scala minima come $a(k) = a_0 2^k$.
Dalla relazione $P(r) \sim r^{-\lambda}$ e sapendo che $\lambda = 1 + D(\tau-1)$, scriviamo:
\begin{equation}
    P(r) \sim r^{-(1 + D(\tau - 1))}
\end{equation}
La probabilità di arresto $R$ corrisponde all'integrale della distribuzione nell'intervallo di scala considerato:
\begin{equation}
    R \approx \text{Pr}(\text{stop tra } k \text{ e } k+1) = \int_{a(k)}^{a(k+1)} P(r) dr
\end{equation}
Sostituendo la legge di potenza:
\begin{equation}
    R \sim \int_{a(k)}^{a(k+1)} r^{-(1 + D(\tau - 1))} dr
\end{equation}
Risolviamo l'integrale definito tra $a(k)$ e $a(k+1) = 2a(k)$:
\begin{equation}
    R \sim \left[ \frac{r^{-D(\tau-1)}}{-D(\tau-1)} \right]_{a(k)}^{2a(k)} \propto a(k)^{-D(\tau-1)} (1 - 2^{-D(\tau-1)})
\end{equation}
Analizzando il comportamento di scaling e invertendo la relazione per $\tau$, si ottiene la formula finale:
\begin{equation}
    \tau = 1 - \frac{1}{D} \frac{\ln(1 - R^*)}{\ln 2}
\end{equation}
Sostituendo i valori numerici ottenuti dal punto fisso, si trova $\tau \approx 1.253$.
Il valore esatto è $\tau = 5/4 = 1.25$.
Per $D$, gli esperimenti suggeriscono $D \approx 2$.

\subsubsection{Calcolo dell'esponente dinamico $z$}
Vogliamo ora determinare l'esponente dinamico $z$, che lega tempo e spazio ($t \sim r^z$).
Consideriamo la durata media di una valanga alla scala $k$, $\langle t \rangle^{(k)}$.
Nella rinormalizzazione temporale, cerchiamo di capire come scala il tempo tra $k$ e $k+1$.
L'esponente $z$ è dato da:
\begin{equation}
    z = \frac{\ln \langle t \rangle}{\ln 2}
\end{equation}
Dove il fattore 2 al denominatore deriva dal raddoppio della scala spaziale ($a(k+1)/a(k)~=~2$); $\langle t \rangle$ rappresenta il fattore di riscalamento temporale medio derivato dalle equazioni dinamiche del gruppo di rinormalizzazione.

\section{La Scienza delle Reti}
\label{chap:network_science}

La \textit{Network Science} è un campo di studio interdisciplinare che analizza la topologia delle interazioni nei sistemi complessi. A differenza della teoria dei grafi classica, che si occupa di proprietà astratte, la scienza delle reti si concentra sulla struttura e sulla dinamica di reti reali (tecnologiche, sociali, biologiche).

\subsection{Il Problema dei Ponti di Königsberg}
La nascita della teoria dei grafi è convenzionalmente datata 1736, con la soluzione di Eulero al problema dei ponti di Königsberg. La città prussiana presentava 7 ponti che collegavano due isole e le rive del fiume Pregel. Il quesito era: \textit{è possibile ideare un percorso che attraversi ogni ponte esattamente una volta?}

Eulero astrasse il problema trasformando la mappa geografica in un grafo:
\begin{itemize}
    \item \textbf{Nodi:} Le quattro zone di terra ferma.
    \item \textbf{Archi:} I sette ponti.
\end{itemize}
Eulero si concentrò sulla topologia e dimostrò che l'esistenza di tale percorso (oggi chiamato \textbf{cammino euleriano}) dipende esclusivamente dai gradi dei nodi.

\begin{tcolorbox}[colback=colorC!15, colframe=colorC!70!colorB,  title=\textbf{Teorema di Eulero-Hierholzer}]
Una condizione necessaria (e sufficiente per grafi connessi) affinché esista un cammino che attraversa ogni arco una sola volta è che il numero di nodi con grado dispari sia esattamente \textbf{zero} o \textbf{due}.
\end{tcolorbox}
\begin{itemize}
    \item Se non ci sono nodi di grado dispari, il cammino è un \textbf{circuito euleriano} (inizia e finisce nello stesso punto).
    \item Se ci sono due nodi di grado dispari, il cammino deve iniziare in uno di essi e terminare nell'altro.
\end{itemize}
Nel caso di Königsberg, i quattro nodi avevano tutti gradi dispari ed il cammino era impossibile.

\subsection{Classificazione e Tipologie di Reti}
Le reti possono essere categorizzate in base alla natura dei loro nodi e delle interazioni.

\subsubsection{Reti Tecnologiche}
\begin{itemize}
    \item \textbf{Internet:} Rappresenta l'infrastruttura fisica (cavi, router, fibre ottiche). La sua struttura è decentralizzata e gerarchica, composta da \textit{Backbone} (connessioni a lunga distanza ad alta larghezza di banda), \textit{ISP Regionali} e \textit{ISP Locali} che connettono gli utenti finali. 
    \item \textbf{World Wide Web (WWW):} È una rete "logica" di informazioni costruita sopra Internet. I nodi sono le pagine web e gli archi sono i collegamenti ipertestuali (URL).
    \item \textbf{Reti Infrastrutturali:} Esempio rilevante è la rete europea dei gasdotti. 
\end{itemize}


\subsubsection{Reti di Informazione}
\begin{itemize}
    \item \textbf{Wikipedia:} Una rete massiva di collaborazione.
    \item \textbf{Reti di Citazioni:} I nodi sono articoli scientifici e un arco diretto rappresenta una citazione da un articolo all'altro.
\end{itemize}

\subsubsection{Reti Sociali e Social Balance}
Le reti sociali mappano le interazioni tra individui (amicizia, collaborazioni, parentele).
Un esempio storico è la rete delle famiglie fiorentine del XV secolo. L'ascesa dei Medici, pur non essendo la famiglia più ricca, è stata attribuita alla loro posizione topologica strategica nella rete di matrimoni e affari, che ha permesso loro di agire come intermediari tra famiglie rivali disconnesse.

\hfill 

\noindent \textbf{Social Balance:} Le relazioni sociali possono essere positive (amicizia $+1$) o negative (inimicizia $-1$). Consideriamo una triade di nodi collegati $i, j, k$. La triade si dice \textbf{bilanciata} se il prodotto dei segni delle relazioni è positivo:
\begin{equation}
    S_{ij} S_{jk} S_{ki} = 1
\end{equation}
Le configurazioni stabili sono:
\begin{enumerate}
    \item Tre amici ($+, +, +$).
    \item Due nemici e un amico comune ($+, -, -$): "Il nemico del mio nemico è mio amico".
\end{enumerate}
Le configurazioni instabili (sbilanciate con prodotto $-1$) tendono a evolvere dinamicamente verso l'equilibrio.

\hfill 

\noindent \textbf{Small World Effect:} L'esperimento di Milgram (1967) ha introdotto il concetto di "Piccolo Mondo".
A partecipanti casuali nel Midwest USA fu chiesto di far recapitare un pacchetto a un target a Boston, passandolo solo a conoscenti diretti. La lunghezza media delle catene completate era di circa 6 passaggi.

Questo ha dato origine all'espressione "6 gradi di separazione": nonostante le grandi dimensioni della rete sociale, la distanza media tra due nodi qualsiasi è molto piccola.


\subsection{Struttura di un Grafo}
Un grafo $G(v, e)$ è un oggetto matematico composto da:
\begin{itemize}
    \item Un insieme di \textbf{nodi} (o veritici) $v$.
    
    Esempio: $\{a, b, c, d\}$.
    \item Un insieme di \textbf{archi} (o link/edges) $e$, costituiti da coppie di nodi. 
    
    Esempio: $\{(a, b), (a, c), (b, c), (b, d)\}$.
\end{itemize}


\noindent Definiamo le seguenti misure fondamentali:
\begin{itemize}
    \item \textbf{Grado ($k$):} Il numero di archi incidenti su un nodo specifico.
    \item \textbf{Adiacenza:} Due nodi sono adiacenti (distanza = 1) quando sono collegati da un arco.
    \item \textbf{Densità:} È definita come il numero di archi presenti in una rete semplice diviso per il numero massimo di archi possibili.
    \item \textbf{Rete Completa:} Una rete che possiede la massima densità possibile (tutti i nodi sono collegati tra loro).
\end{itemize}

\noindent Per i \textbf{grafi diretti}, il concetto di grado si suddivide in:
\begin{itemize}
    \item \textbf{In-degree ($k_{in}$):} Il numero di archi che il nodo riceve (frecce entranti).
    \item \textbf{Out-degree ($k_{out}$):} Il numero di archi che il nodo invia (frecce uscenti).
\end{itemize}

Il grado totale è la somma: $k = k_{in} + k_{out}$.

\hfill 

\noindent La distanza tra due nodi è definita come il numero minimo di archi necessari per collegarli. Questo percorso minimo è chiamato \textit{shortest path} o \textbf{cammino geodetico}.
Esistono algoritmi specifici per calcolare il cammino geodetico. Nelle reti reali si osserva spesso lo "Small-world effect", ovvero la distanza media tra i nodi è sorprendentemente piccola.

\subsection{Matrice di Adiacenza}
La struttura topologica di una rete è completamente descritta dalla sua \textbf{matrice di adiacenza} $A$. 

\subsubsection{\textcolor{colorA}{Grafi Non Diretti}}
Per un grafo non diretto con $n$ nodi, la matrice di adiacenza $A$ è una matrice quadrata $n \times n$ simmetrica. Gli elementi sono definiti come segue:
\begin{equation}
A_{ij} = 
\begin{cases} 
1 & \text{se esiste un arco tra } i \text{ e } j \\
0 & \text{altrimenti}
\end{cases}
\end{equation}

\noindent Proprietà della matrice:
\begin{itemize}
    \item \textbf{Simmetria:} Poiché il grafo è non diretto, la matrice è simmetrica rispetto alla diagonale principale ($A_{ij} = A_{ji}$).
    \item \textbf{Self-loops:} Se $A_{ii} \neq 0$, significa che esistono collegamenti ricorrenti (cappi) che connettono un nodo a se stesso.
    \item \textbf{Multi-edges:} Se $A_{ii} > 1$, significa che ci sono archi multipli tra la stessa coppia di nodi.
\end{itemize}

\noindent \textbf{{Esempio:}} Consideriamo un grafo con 4 nodi e le connessioni: $(1,2), (1,4), (2,4), (3,4)$.

\hfill

\noindent
\begin{minipage}{0.35\textwidth}
\begin{equation*}
A = 
\begin{pmatrix} 
0 & 1 & 0 & 1 \\ 
1 & 0 & 0 & 1 \\ 
0 & 0 & 0 & 1 \\ 
1 & 1 & 1 & 0 
\end{pmatrix}
\end{equation*}
\end{minipage}%
\hfill
\begin{minipage}{0.55\textwidth}
\centering
\begin{tikzpicture}[auto, node distance=1cm, thick, scale=0.5,
                    main node/.style={circle,draw,font=\bfseries}]
  \node[main node] (1) {1};
  \node[main node] (2) [below left of=1] {2};
  \node[main node] (4) [below right of=1] {4};
  \node[main node] (3) [below right of=4] {3};

  \path
    (1) edge (2)
    (1) edge (4)
    (2) edge (4)
    (4) edge (3);
\end{tikzpicture}
\end{minipage}



\subsubsection{\textcolor{colorA}{Grafi Pesati Non Diretti}}
La matrice di adiacenza non è necessariamente binaria. È possibile avere una \textbf{forza di interazione} o peso associato a ciascun arco.
La definizione diventa:
\begin{equation}
A_{ij} = 
\begin{cases} 
w_{ij} & \text{peso del collegamento tra } i \text{ e } j \\
0 & \text{se } i \text{ e } j \text{ non sono collegati}
\end{cases}
\end{equation}

\noindent \textbf{Esempio:} 
Consideriamo i seguenti pesi: $w_{1,2}=4$, $w_{1,4}=3.2$, $w_{3,4}=2$, $w_{2,4}=0.5$.

\hfill 

\noindent
\begin{minipage}{0.35\textwidth}
\begin{equation*}
A = 
\begin{pmatrix} 
0 & 4 & 0 & 3.2 \\ 
4 & 0 & 0 & 0.5 \\ 
0 & 0 & 0 & 2 \\ 
3.2 & 0.5 & 2 & 0 
\end{pmatrix}
\end{equation*}
\end{minipage}%
\hfill
\begin{minipage}{0.55\textwidth}
\centering
\begin{tikzpicture}[auto, node distance=1.6cm, thick, scale=0.5,
                    main node/.style={circle,draw,font=\bfseries}]
  \node[main node] (1) {1};
  \node[main node] (2) [below left of=1] {2};
  \node[main node] (4) [below right of=1] {4};
  \node[main node] (3) [below right of=4] {3};

  \path
    (1) edge node[above left] {4} (2)
    (1) edge node[above right] {3.2} (4)
    (2) edge node[below] {0.5} (4)
    (4) edge node[above right] {2} (3);
\end{tikzpicture}
\end{minipage}

\hfill \\


\noindent Dato un grafo non diretto con $N$ nodi, il grado $k_i$ del nodo $i$ è definito come il numero di connessioni che esso possiede. In termini di matrice di adiacenza:
\begin{tcolorbox}[colback=yellow!25, colframe=yellow!75!orange, coltitle=black, title=\textbf{Grado di un nodo}]
\begin{equation}
k_i = \sum_{j=1}^N A_{ij}
\end{equation}
\end{tcolorbox}
\noindent Il numero totale di archi $m$ nella rete è legato alla somma dei gradi dalla seguente relazione (Handshake Lemma):
\begin{equation}
\sum_{i=1}^N k_i = \sum_{i=1}^N \sum_{j=1}^N A_{ij} = 2m
\end{equation}
Il fattore 2 deriva dal fatto che ogni arco connette due nodi e viene quindi contato due volte nella sommatoria (una volta per $i$ e una volta per $j$).
Possiamo quindi scrivere:
\begin{tcolorbox}[colback=yellow!25, colframe=yellow!75!orange, coltitle=black, title=\textbf{Numero totale di archi}]
\begin{equation}
m = \frac{1}{2} \sum_{i,j} A_{ij}
\end{equation}
\end{tcolorbox}
\noindent Il \textbf{grado medio} $\langle k \rangle$ è dato da:
\begin{tcolorbox}[colback=yellow!25, colframe=yellow!75!orange, coltitle=black, title=\textbf{Grado medio}]
\begin{equation}
\langle k \rangle = \frac{1}{N} \sum_{i=1}^N k_i = \frac{2m}{N}
\end{equation}
\end{tcolorbox}

\newpage
\subsubsection{\textcolor{colorA}{Grafi Diretti}}
Nei grafi diretti, l'interazione ha una direzione. La matrice di adiacenza $A$ \textbf{non è simmetrica}.
È cruciale definire la convenzione utilizzata, poiché cambia la lettura della matrice. Adottiamo la seguente convenzione (\textbf{Input Convention}):
\begin{equation}
A_{ij} = 
\begin{cases} 
1 & \text{se c'è un link \textbf{da} } j \text{ A } i \quad (j \to i) \\
0 & \text{altrimenti}
\end{cases}
\end{equation}
In altre parole, la riga $i$ raccoglie gli input che il nodo $i$ riceve dagli altri nodi ($j$).

\noindent \textbf{Esempio:} 
Consideriamo archi diretti: $1 \to 2$, $1 \to 3$, $3 \to 2$, $3 \to 4$, $4 \to 3$.

\hfill 

\noindent
\begin{minipage}{0.35\textwidth}
\begin{equation*}
A = 
\begin{pmatrix} 
0 & 0 & 0 & 0 \\ 
1 & 0 & 1 & 0 \\ 
1 & 0 & 0 & 1 \\ 
0 & 0 & 1 & 0 
\end{pmatrix}
\end{equation*}

\end{minipage}%
\hfill
\begin{minipage}{0.55\textwidth}
\centering
\begin{tikzpicture}[->, >=stealth, auto, node distance=1.5cm, thick, scale=0.5,
                    main node/.style={circle,draw,font=\bfseries}]
  \node[main node] (1) {1};
  \node[main node] (2) [right of=1] {2};
  \node[main node] (3) [below of=1] {3};
  \node[main node] (4) [below of=2] {4};

  \path
    (1) edge (2)
    (1) edge (3)
    (3) edge (2)
    (3) edge[bend left] (4)
    (4) edge[bend left] (3);
\end{tikzpicture}
\end{minipage}

\hfill

\noindent Analisi della matrice:
\begin{itemize}
    \item Riga 1: tutta zeri $\rightarrow$ Il nodo 1 non riceve link da nessuno.
    \item Riga 2: colonne 1 e 3 $\rightarrow$ Il nodo 2 riceve da 1 e da 3.
    \item Riga 3: colonne 1 e 4 $\rightarrow$ Il nodo 3 riceve da 1 e da 4.
    \item Riga 4: colonna 3 $\rightarrow$ Il nodo 4 riceve da 3.
\end{itemize}


È possibile avere processi dinamici su reti fisse (dinamica \textit{sopra} la topologia) oppure reti con archi dipendenti dal tempo (dinamica \textit{della} topologia). Inoltre, si possono avere feedback loops (cicli di retroazione) tra i due processi.

\hfill \\

\noindent Nei grafi diretti, distinguiamo tra grado entrante ($k_{in}$) e grado uscente ($k_{out}$). Utilizzando la convenzione dell'input ($A_{ij}=1 \iff j \to i$):

\begin{itemize}
    \item \textbf{In-degree ($k_i^{in}$):} Somma sulla riga $i$. Rappresenta il numero di archi che arrivano al nodo $i$.
    \begin{tcolorbox}[colback=yellow!25, colframe=yellow!75!orange, coltitle=black, title=\textbf{Grado Entrante}]
    \begin{equation}
    k_i^{in} = \sum_{j=1}^N A_{ij}
    \end{equation}
    \end{tcolorbox}
    
    
    
    \item \textbf{Out-degree ($k_j^{out}$):} Somma sulla colonna $j$. Rappresenta il numero di archi che partono dal nodo $j$.
     \begin{tcolorbox}[colback=yellow!25, colframe=yellow!75!orange, coltitle=black, title=\textbf{Grado Uscente}]
     \begin{equation}
    k_j^{out} = \sum_{i=1}^N A_{ij}
    \end{equation}
    \end{tcolorbox}
\end{itemize}
Per il numero totale di archi $m$, non vi è il doppio conteggio presente nei grafi non diretti, poiché ogni arco ha una direzione univoca. La somma di tutti i gradi entranti deve eguagliare la somma di tutti i gradi uscenti, ed entrambe sono uguali a $m$:
\begin{equation}
m = \sum_{i=1}^N k_i^{in} = \sum_{j=1}^N k_j^{out} = \sum_{i,j} A_{ij}
\end{equation}
Di conseguenza, il grado medio entrante è uguale al grado medio uscente:
\begin{equation}
\langle k_{in} \rangle = \frac{1}{N}\sum_{i=1}^N k_i^{in} = \frac{m}{N}
\end{equation}
\begin{equation}
\langle k_{out} \rangle = \frac{1}{N}\sum_{j=1}^N k_j^{out} = \frac{m}{N}
\end{equation}
Quindi vale sempre:
 \begin{tcolorbox}[colback=yellow!25, colframe=yellow!75!orange, coltitle=black, title=\textbf{Grado Medio}]
     \begin{equation}
    \langle k_{in} \rangle = \langle k_{out} \rangle
    \end{equation}
\end{tcolorbox}

\noindent La \textbf{densità} $\rho$ della rete è il rapporto tra gli archi esistenti e il numero massimo di archi possibili (che è $\binom{N}{2}$):
\begin{tcolorbox}[colback=yellow!25, colframe=yellow!75!orange, coltitle=black, title=\textbf{Densità}]
\begin{equation}
\rho = \frac{m}{N(N-1)/2} = \frac{\langle k \rangle}{N-1}
\end{equation}
\end{tcolorbox}

Analizziamo il comportamento della densità nel limite di reti di grandi dimensioni, ovvero per $n \gg 1$.
In questo regime, possiamo approssimare $n-1 \approx n$. L'espressione della densità si semplifica in:

\begin{equation}
\rho \approx \frac{\langle k \rangle}{n}
\end{equation}
Essendo una probabilità (o frazione di archi esistenti), la densità è sempre limitata nell'intervallo:
\begin{equation}
0 \le \rho \le 1
\end{equation}
A seconda di come il grado medio $\langle k \rangle$ scala al crescere della dimensione della rete $n$, possiamo distinguere diversi regimi:

\begin{itemize}
    \item \textbf{Reti Dense:} Se il grado medio scala linearmente con la dimensione della rete $\langle k \rangle \propto n$, allora il rapporto si stabilizza e la densità rimane costante al crescere di $n$:
    \begin{equation}
    \rho = \text{const}
    \end{equation}

    \item \textbf{Reti Sparse (Comportamento Sub-lineare):} Se il grado medio cresce, ma più lentamente di $n$ (crescita sub-lineare), allora la densità tende asintoticamente a zero:
    \begin{equation}
    \text{Se } \langle k \rangle \text{ cresce sub-linearmente} \implies \rho \to 0
    \end{equation}

    \item \textbf{Reti Estremamente Sparse (Extremely Sparse Networks):} Questo è il caso più comune nelle reti reali di grandi dimensioni. Qui il grado medio è costante e non dipende dalla dimensione della rete $\langle k \rangle = \text{const}$. Di conseguenza, la densità decresce come l'inverso di $n$:
    \begin{equation}
    \langle k \rangle = \text{const} \implies \rho \sim \frac{1}{n} \to 0
    \end{equation}
\end{itemize}


\subsection{Grafi Diretti Aciclici (DAG)}
Un \textbf{DAG} (Directed Acyclic Graph) è un grafo diretto che non contiene cicli chiusi.
Una proprietà fondamentale dei DAG è legata alla struttura della loro matrice di adiacenza.
\begin{tcolorbox}[colback=colorC!15, colframe=colorC!70!colorB,  title=\textbf{Teorema:}]
Per ogni DAG, esiste almeno un'etichettatura (labeling) dei nodi tale che la matrice di adiacenza risulti \textbf{strettamente triangolare superiore}.
\end{tcolorbox}

Questo significa che:
\begin{equation}
A_{ij} = 0 \quad \text{per tutti gli } i \ge j
\end{equation}
In una matrice triangolare superiore, tutti gli elementi sulla diagonale principale e al di sotto di essa sono nulli. Questo implica che non ci sono link che "tornano indietro" rispetto all'ordine stabilito dei nodi.


\noindent
\begin{minipage}{0.35\textwidth}
{Esempio di Matrice \\ Triangolare Superiore:}
\begin{equation*}
A = 
\begin{pmatrix} 
0 & 1 & 1 & 0 & 1 \\ 
0 & 0 & 0 & 0 & 1 \\ 
0 & 0 & 0 & 0 & 1 \\ 
0 & 0 & 0 & 0 & 0 \\ 
0 & 0 & 0 & 0 & 0 
\end{pmatrix}
\end{equation*}
\end{minipage}
\hfill
\begin{minipage}{0.55\textwidth}
\centering
\begin{tikzpicture}[->, >=stealth, auto, node distance=1cm, thick,
                    main node/.style={circle,draw,font=\bfseries,fill=white, minimum size=0.8cm}]

  % Layout gerarchico allargato per evitare sovrapposizioni
  \node[main node] (5) at (0, 4) {5};    % Sorgente
  \node[main node] (3) at (-2, 2) {3};   % Sx
  \node[main node] (4) at (2, 2) {4};    % Dx
  \node[main node] (2) at (0, 0) {2};    % Centro
  \node[main node] (1) at (0, -2.5) {1}; % Pozzo

  \path
    % Da 5 (e)
    (5) edge (3)
    (5) edge (4)
    (5) edge (2) % Passa in mezzo a 3 e 4
    (5) edge[out=10, in=10, looseness=1.5] (1) % Curva larga esterna a sinistra
    
    % Da 4 (d)
    (4) edge (2)
    
    % Da 3 (c)
    (3) edge (2)
    (3) edge (1) % Curva leggera per evitare il nodo 2
    
    % Da 2 (b)
    (2) edge (1);

\end{tikzpicture}
\end{minipage}



\subsection{Reti Bipartite}
Una \textbf{rete bipartita} è un grafo $G=(U,V,E)$ i cui nodi sono divisi in due insiemi disgiunti, $U$ e $V$, tali che ogni arco collega un nodo di $U$ a uno di $V$. Non esistono collegamenti tra nodi appartenenti allo stesso insieme.

Esempi classici includono: Attori e Film ($U$ sono gli attori, $V$ sono i film. Un arco esiste se l'attore recita nel film), Autori e Articoli ($U$ sono gli autori, $V$ sono le pubblicazioni), Entità e Gruppi...


\noindent Poiché la rete connette due classi distinte di oggetti, non si utilizza la classica matrice di adiacenza quadrata, bensì la \textbf{Matrice di Incidenza} $B$.
Se abbiamo $g$ gruppi e $n$ entità, la matrice $B$ avrà dimensione $g \times n$.
La definizione formale è:
\begin{equation}
B_{ij} = 
\begin{cases} 
1 & \text{se l'entità } j \text{ appartiene al gruppo } i \\
0 & \text{altrimenti}
\end{cases}
\end{equation}
\textbf{Rappresentazione Grafica:}
Di seguito è riportato un esempio di rete bipartita con la relativa matrice di incidenza.

\vspace{0.5cm}
\noindent
\begin{minipage}{0.35\textwidth}
\begin{equation*}
B = 
\begin{pmatrix} 
1 & 0 & 1 & 0 & 0 \\ 
0 & 0 & 1 & 0 & 0 \\ 
0 & 1 & 0 & 1 & 0 \\ 
1 & 0 & 0 & 1 & 1 
\end{pmatrix}
\end{equation*}
\end{minipage}%
\hfill
\begin{minipage}{0.55\textwidth}
\centering
\begin{tikzpicture}[thick,
    group node/.style={rectangle,draw,fill=colorC,minimum size=8mm,font=\small\bfseries, text=white,},
    entity node/.style={circle,draw,fill=colorE,minimum size=8mm,font=\small\bfseries}, text=white,]
    
    % Gruppi a sinistra
    \node[group node] (g1) at (0,3) {G1};
    \node[group node] (g2) at (0,2) {G2};
    \node[group node] (g3) at (0,1) {G3};
    \node[group node] (g4) at (0,0) {G4};
    
    % Entità a destra
    \node[entity node] (e1) at (3,3.2) {E1};
    \node[entity node] (e2) at (3,2.4) {E2};
    \node[entity node] (e3) at (3,1.6) {E3};
    \node[entity node] (e4) at (3,0.8) {E4};
    \node[entity node] (e5) at (3,0) {E5};

    % Archi basati sulla matrice B
    \draw (g1) -- (e1);
    \draw (g1) -- (e3);
    \draw (g2) -- (e3);
    \draw (g3) -- (e2);
    \draw (g3) -- (e4);
    \draw (g4) -- (e1);
    \draw (g4) -- (e4);
    \draw (g4) -- (e5);
\end{tikzpicture}
\end{minipage}

\subsubsection{Proiezioni One-Mode}
Spesso si è interessati alle relazioni dirette tra nodi dello stesso tipo (es. "quali attori hanno recitato insieme?"). Per ottenere queste informazioni, si effettua una proiezione della rete bipartita su uno dei due insiemi.

\begin{enumerate}
    \item \textbf{Proiezione Entità-Entità:}
Questa proiezione crea una rete composta solo dalle $n$ entità. Due entità sono collegate se condividono almeno un gruppo.
Matematicamente, si ottiene moltiplicando la trasposta di $B$ per $B$ stessa:
\begin{equation}
P = B^T B
\end{equation}
$P$ è una matrice quadrata $n \times n$ simmetrica. L'elemento generico $P_{ij}$ è dato da:
\begin{equation}
P_{ij} = \sum_{k=1}^g (B^T)_{ik} B_{kj} = \sum_{k=1}^g B_{ki} B_{kj}
\end{equation}
L'interpretazione fisica dei termini è la seguente:
\begin{itemize}
    \item \textbf{Elementi diagonali ($P_{ii}$):} Rappresentano il numero di gruppi a cui partecipa l'entità $i$.
    \item \textbf{Elementi fuori diagonale ($P_{ij}$ con $i \neq j$):} Rappresentano il numero di gruppi che le entità $i$ e $j$ condividono.
\end{itemize}
Per ottenere la matrice di adiacenza classica (binaria o pesata senza self-loops), si impone solitamente $A_{ii} = 0$.
\item \textbf{Proiezione Gruppo-Gruppo:}
Analogamente, possiamo proiettare sui gruppi per vedere quali gruppi condividono entità.
La matrice risultante $P'$ è quadrata di dimensione $g \times g$:
\begin{equation}
P' = B B^T
\end{equation}
L'elemento generico è:
\begin{equation}
P'_{ij} = \sum_{k=1}^n B_{ik} (B^T)_{kj} = \sum_{k=1}^n B_{ik} B_{jk}
\end{equation}
Qui, $P'_{ij}$ indica quante entità sono presenti contemporaneamente sia nel gruppo $i$ che nel gruppo $j$.
\end{enumerate}




\subsection{Alberi}
Un \textbf{albero} è una tipologia fondamentale di grafo. Si definisce albero un grafo connesso e non diretto che non contiene cicli (loops).
\begin{enumerate}
    \item In un albero esiste \textbf{uno e un solo cammino} che collega qualsiasi coppia di vertici.
    \item Un albero con $N$ vertici possiede esattamente $N-1$ archi.
    \item Se si aggiunge un arco a un albero, si crea necessariamente un ciclo.
    \item Se si rimuove un qualsiasi arco da un albero, il grafo diventa sconnesso.
\end{enumerate}

\begin{center}
\begin{tikzpicture}[thick, scale=1.3, every node/.style={circle,draw,fill=white,font=\small}]
  \node (root) at (0,0) {Root}
    child {node {A}
      child {node {A1}}
      child {node {A2}}
    }
    child {node {B}
      child {node {B1}}
    }
    child {node {C}
      child {node {C2}}
      child {node {C3}}
    };
\end{tikzpicture}
\end{center}

\chapter{Lezione 8}
\label{chap:lezione_08} 

\begin{flushright}
\textit{Data: 23/10/2025}
\end{flushright}

\section{Passeggiate e Cammini}

	\begin{description}
		\item[Passeggiata (Walk):] Una passeggiata su una rete è una qualsiasi sequenza di nodi tale che ogni coppia consecutiva di nodi sia connessa da un arco (edge). In termini matematici, una passeggiata può visitare lo stesso nodo e lo stesso arco più volte. È la definizione con il minor numero di vincoli possibili.
		
		\item[Cammino (Path):] Un cammino è una passeggiata auto-evitante (\textit{self-avoiding}). Non visita mai lo stesso nodo o lo stesso arco due volte.
	\end{description}
	
	\subsection{Passeggiate e Matrice di Adiacenza}
	La matrice di adiacenza $A$ può essere utilizzata per contare il numero di passeggiate tra i nodi. Adottiamo la convenzione per cui $A_{ij} = 1$ se esiste un collegamento \textbf{da} $j$ \textbf{a} $i$.
	
	\begin{itemize}
		\item \textbf{Lunghezza 2:} Il numero di passeggiate di lunghezza 2 dal nodo $j$ al nodo $i$ è dato dalla somma su tutti i possibili nodi intermedi $k$:
		$$N_{ij}^{(2)} = \sum_{k=1}^{n} A_{ik} A_{kj} = [A^2]_{ij}$$
		Questo perché il prodotto $A_{ik}A_{kj}$ è non nullo (uguale a 1) solo se esiste la passeggiata $j \to k \to i$. Si noti che la somma avviene su tutti i $k$, senza imporre che il percorso sia auto-evitante.
		
		\item \textbf{Lunghezza 3:} Analogamente, per le passeggiate di lunghezza 3 da $j$ a $i$ (attraverso i passaggi intermedi $l$ e $k$):
		$$N_{ij}^{(3)} = \sum_{k,l=1}^{n} A_{ik} A_{kl} A_{lj} = [A^3]_{ij}$$
		
		\item \textbf{Lunghezza r:} In generale, il numero di passeggiate di lunghezza $r$ da $j$ a $i$ corrisponde all'elemento $(i, j)$ della matrice $A^r$:
		$$N_{ij}^{(r)} = [A^r]_{ij}$$
	\end{itemize}
	
	\subsection{Cicli}
	Un ciclo è una passeggiata che inizia e termina nello stesso nodo.
	\begin{itemize}
		\item Il numero di cicli di lunghezza $r$ che iniziano e finiscono nel nodo $i$ è dato dall'elemento diagonale:
		$$N_{ii}^{(r)} = [A^r]_{ii}$$
		
		\item Il numero totale di cicli di lunghezza $r$ nell'intera rete è la somma degli elementi diagonali, ovvero la traccia di $A^r$:
		$$L_r = \sum_{i=1}^{n} N_{ii}^{(r)} = \sum_{i=1}^{n} [A^r]_{ii} = \text{Tr}(A^r)$$
	\end{itemize}
	
	Contare il numero preciso di percorsi o cicli \textit{distinti} non è banale. Il metodo della traccia conta le permutazioni (es. $1 \to 2 \to 3 \to 1$ e $2 \to 3 \to 1 \to 2$) e (per grafi non orientati) le direzioni opposte (es. $1 \to 2 \to 3 \to 1$ e $1 \to 3 \to 2 \to 1$) come passeggiate distinte, anche se tracciano lo stesso ciclo geometrico.
	
	\subsection*{Esempi}
	\subsubsection{Grafo Triangolare Non Orientato}
	Si consideri un semplice grafo non orientato di 3 nodi, completamente connesso (un triangolo).
	La matrice di adiacenza $A$ e il suo quadrato $A^2$ sono:
	$$A = \begin{pmatrix} 0 & 1 & 1 \\ 1 & 0 & 1 \\ 1 & 1 & 0 \end{pmatrix} \quad \implies \quad A^2 = \begin{pmatrix} 2 & 1 & 1 \\ 1 & 2 & 1 \\ 1 & 1 & 2 \end{pmatrix}$$
	
	\begin{itemize}
		\item $A^2_{11} = 2$: Esistono 2 passeggiate di lunghezza 2 dal nodo 1 a se stesso. Queste sono $1 \to 2 \to 1$ e $1 \to 3 \to 1$ 
		\item $A^2_{12} = 1$: Esiste 1 passeggiata di lunghezza 2 dal nodo 2 al nodo 1. Questa è $2 \to 3 \to 1$.
	\end{itemize}

	Il cubo di $A$ è:
	$$A^3 = \begin{pmatrix} 2 & 3 & 3 \\ 3 & 2 & 3 \\ 3 & 3 & 2 \end{pmatrix}$$
	
	\begin{itemize}
		\item $A^3_{11} = 2$: Esistono 2 cicli di lunghezza 3 che iniziano al nodo 1: $1 \to 2 \to 3 \to 1$ e $1 \to 3 \to 2 \to 1$.
		\item $\text{Tr}(A^3) = 2+2+2 = 6$. Questi sono i 6 cicli totali di lunghezza 3 (2 per ogni nodo di partenza).
	\end{itemize}
	
	\subsubsection{Grafo Triangolare Orientato}
	Si consideri un ciclo orientato $2 \to 1$, $1 \to 3$, $3 \to 2$.
	La matrice di adiacenza $A$ (con $A_{ij}$ definita come $j \to i$) e il suo quadrato $A^2$ sono:
	$$A = \begin{pmatrix} 0 & 1 & 0 \\ 0 & 0 & 1 \\ 1 & 0 & 0 \end{pmatrix} \quad \implies \quad A^2 = \begin{pmatrix} 0 & 0 & 1 \\ 1 & 0 & 0 \\ 0 & 1 & 0 \end{pmatrix}$$
	
	\begin{itemize}
		\item La diagonale di $A^2$ è composta interamente da zeri, il che significa che non esistono loop di lunghezza 2 (il loop più piccolo possibile è di lunghezza 3).
		\item $A^2_{13} = 1$: Esiste una passeggiata di lunghezza 2 dal nodo 3 al nodo 1, che corrisponde a $3 \to 2 \to 1$
	\end{itemize}
	
	%==================================================================
	% PARTE 2: LAPLACIANO, RANDOM WALKS
	%==================================================================
	
	\section{La Matrice Laplaciana}
	Per un \textbf{grafo non orientato e \color{colorF}{non pesato}}, la \textbf{matrice Laplaciana} $L$ è definita come:
    \begin{tcolorbox}[colback=yellow!25, colframe=yellow!75!orange, coltitle=black, title=\textbf{Matrice Laplaciana}]
\begin{equation}
L_{ij} = \begin{cases}
		k_i & \text{se } i = j \\
		-1  & \text{se } i \neq j \text{ e } A_{ij} \neq 0 \\
		0   & \text{altrimenti}
	\end{cases}
\end{equation}
\end{tcolorbox}

	Dove $k_i$ è il grado del nodo $i$.
	
	\noindent In forma matriciale, si può scrivere come:
	\begin{equation}
	    L=D-A
	\end{equation}
	dove $D$ è la \textbf{matrice diagonale dei gradi} ($D_{ii} = k_i$) e $A$ è la \textbf{matrice di adiacenza}. Ciò è equivalente alla definizione compatta:
    \begin{equation}
        L_{ij} = \delta_{ij} k_i - A_{ij}
    \end{equation}
	
	
	\subsection{Il Laplaciano Discreto}
	Il nome deriva dalla sua relazione con l'operatore Laplaciano continuo $\nabla^2$.
	In uno spazio continuo 1D, la derivata seconda è $\frac{\partial^2 \phi}{\partial x^2}$. Un'approssimazione alle differenze finite è:
	\begin{equation}
	    \frac{\partial^2 \phi}{\partial x^2} \approx [\phi(x+1) - \phi(x)] - [\phi(x) - \phi(x-1)] = \phi(x+1) + \phi(x-1) - 2\phi(x)
	\end{equation}
	In generale, l'operatore Laplaciano discreto è spesso definito come la somma dei vicini meno $k$ volte il valore centrale, con $k$ numero totale di vicini:
    \begin{equation}
        \nabla^2 \phi_i \approx \left( \sum_{\text{vicini } j} \phi_j \right) - k_i \phi_i
    \end{equation}
	Consideriamo ora un campo scalare $\vec{\phi} = (\phi_1, ..., \phi_n)^T$ definito sui nodi di una rete. Vediamo come agisce la matrice $L$ su di esso. L'$i$-esima componente del vettore $L\vec{\phi}$ è:
    \begin{equation}
        (L\vec{\phi})_i = \sum_{j} L_{ij}\phi_j = \sum_{j} (\delta_{ij}k_i - A_{ij})\phi_j = k_i \phi_i - \sum_{j} A_{ij}\phi_j
    \end{equation}
	Poiché $A_{ij}$ è non nullo solo per i vicini di $i$, l'espressione diventa:
    \begin{equation}
(L\vec{\phi})_i = k_i \phi_i - \sum_{j, \text{ vicini di } i} \phi_j
    \end{equation}
	Questa può essere riscritta come:
    \begin{equation}
        (L\vec{\phi})_i = \sum_{j, \text{ vicini di } i} (\phi_i - \phi_j)
    \end{equation}
	Si noti che $(L\vec{\phi})_i = -(\nabla^2 \phi_i)$. Il Laplaciano del grafo $L = D-A$ corrisponde al \textbf{negativo} dell'operatore Laplaciano discreto standard. Esso misura la differenza totale tra il valore del campo nel nodo $i$ e i valori nei suoi vicini.
	
	\subsubsection{Esempio}
	Consideriamo il seguente grafo quasi-bipartito:
    \begin{figure}[h!]
    \centering
    \begin{tikzpicture}
        \tikzset{
            node_style/.style={circle, fill=colorC, text=white, font=\large, minimum size=0.9cm},
            edge_style/.style={color=colorE, line width=1.2pt}
        }
        % Nodi
        \node[node_style] (1) at (0, 2.5) {1};
        \node[node_style] (2) at (3, 2.5) {2};
        \node[node_style] (3) at (-0.5, 0) {3};
        \node[node_style] (4) at (1.5, -0.5) {4};
        \node[node_style] (5) at (3.5, 0) {5};
        % Archi
        \foreach \i in {1,2} {
            \foreach \j in {3,4,5} {
                \draw[edge_style] (\i) -- (\j);
            }
        }
    \end{tikzpicture}
\end{figure}

La matrice di adiacenza $A$ e la matrice dei gradi $D$ sono:
	$$ A = \begin{pmatrix}
		0 & 0 & 1 & 1 & 1 \\
		0 & 0 & 1 & 1 & 1 \\
		1 & 1 & 0 & 0 & 0 \\
		1 & 1 & 0 & 0 & 0 \\
		1 & 1 & 0 & 0 & 0 
	\end{pmatrix} \quad
	D = \begin{pmatrix}
		3 & 0 & 0 & 0 & 0 \\
		0 & 3 & 0 & 0 & 0 \\
		0 & 0 & 2 & 0 & 0 \\
		0 & 0 & 0 & 2 & 0 \\
		0 & 0 & 0 & 0 & 2 
	\end{pmatrix} $$
	
	Il Laplaciano $L = D - A$ risulta:
	$$ L = \begin{pmatrix}
		3 &  0 & -1 & -1 & -1 \\
		0 &  3 & -1 & -1 & -1 \\
		-1 & -1 &  2 &  0 &  0 \\
		-1 & -1 &  0 &  2 &  0 \\
		-1 & -1 &  0 &  0 &  2 
	\end{pmatrix} $$
	
	\subsection{Proprietà del Laplaciano}
	\begin{enumerate}
		\item \textbf{Somma di Riga/Colonna:} La somma di tutti gli elementi in qualsiasi riga (e colonna) è 0.
        \begin{equation}
            \sum_{j} L_{ij} = \sum_{j} (\delta_{ij}k_i - A_{ij}) = k_i - \sum_{j} A_{ij} = k_i - k_i = 0
        \end{equation}
		Essendo $L$ simmetrica per un grafo non orientato, anche la somma su colonna è nulla.
		
		\item \textbf{Autovalori Reali:} $L$ è una matrice \textbf{reale e simmetrica}, quindi è \textbf{hermitiana} 
        $L^\dagger = L^T = L$
        Di conseguenza, \textbf{tutti i suoi autovalori sono reali}.
		
		\textit{Dimostrazione:}
        
		Si parte dall'equazione agli autovalori $L\vec{v} = \lambda\vec{v}$.
		Si prende la trasposta coniugata di entrambi i membri: $(L\vec{v})^\dagger = (\lambda\vec{v})^\dagger$, che fornisce $\vec{v}^\dagger L^\dagger = \lambda^* \vec{v}^\dagger$.
		Poiché $L$ è hermitiana, $L^\dagger = L$, quindi $\vec{v}^\dagger L = \lambda^* \vec{v}^\dagger$.
		Moltiplicando l'equazione originale a sinistra per $\vec{v}^\dagger$:
        \begin{equation}
            \vec{v}^\dagger L \vec{v} = \lambda \vec{v}^\dagger \vec{v}
        \end{equation}
		Moltiplicando l'equazione trasposta a destra per $\vec{v}$:
        \begin{equation}
           \vec{v}^\dagger L \vec{v} = \lambda^* \vec{v}^\dagger \vec{v} 
        \end{equation}
		Uguagliando le due espressioni si ottiene:
        \begin{equation}
            \lambda \vec{v}^\dagger \vec{v} = \lambda^* \vec{v}^\dagger \vec{v} \implies (\lambda - \lambda^*) \vec{v}^\dagger \vec{v} = 0
        \end{equation}
		Poiché $\vec{v}$ è un autovettore, è non nullo, e $\vec{v}^\dagger \vec{v} = ||\vec{v}||^2$ è un numero reale positivo.
		Pertanto, deve essere $\lambda - \lambda^* = 0$, il che implica $\lambda = \lambda^*$, ovvero $\lambda$ è reale.
		
		\item \textbf{Semidefinita Positiva:} Tutti gli autovalori di $L$ sono non negativi ($\lambda \ge 0$).
		
		\textit{Dimostrazione:}
        
        Espandiamo la forma quadratica $\vec{v}^T L \vec{v}$.
\begin{equation}
	\vec{v}^T L \vec{v} = \sum_{i,j} L_{ij} v_i v_j = \sum_{i} L_{ii} v_i^2 + \sum_{i \neq j} L_{ij} v_i v_j
\end{equation}
Usando la definizione $L_{ij} = \delta_{ij}k_i - A_{ij}$:
\begin{equation}
	= \sum_{i} k_i v_i^2 - \sum_{i \neq j} A_{ij} v_i v_j
\end{equation}
Assumendo assenza di self-loops, $A_{ii}=0$, quindi $\sum_{i \neq j} A_{ij} v_i v_j = \sum_{i,j} A_{ij} v_i v_j$
\begin{equation}
	= \sum_{i} k_i v_i^2 - \sum_{i,j} A_{ij} v_i v_j
\end{equation}
Consideriamo ora l'espressione $\frac{1}{2} \sum_{i,j} A_{ij} (v_i - v_j)^2$:
\begin{equation}
	= \frac{1}{2} \sum_{i,j} A_{ij} (v_i^2 + v_j^2 - 2v_i v_j)
\end{equation}
\begin{equation}
	= \frac{1}{2} \left( \sum_{i,j} A_{ij} v_i^2 + \sum_{i,j} A_{ij} v_j^2 - 2\sum_{i,j} A_{ij} v_i v_j \right)
\end{equation}
Poiché $A$ è simmetrica, $\sum_{j} A_{ij} = k_i$ e $\sum_{i} A_{ij} = k_j$:
\begin{equation}
	= \frac{1}{2} \left( \sum_{i} v_i^2 \left(\sum_{j} A_{ij}\right) + \sum_{j} v_j^2 \left(\sum_{i} A_{ij}\right) - 2\sum_{i,j} A_{ij} v_i v_j \right)
\end{equation}
\begin{equation}
	= \frac{1}{2} \left( \sum_{i} k_i v_i^2 + \sum_{j} k_j v_j^2 - 2\sum_{i,j} A_{ij} v_i v_j \right)
\end{equation}
Essendo $i$ e $j$ indici muti, $\sum_{i} k_i v_i^2 = \sum_{j} k_j v_j^2$.
\begin{equation}
	= \frac{1}{2} \left( 2\sum_{i} k_i v_i^2 - 2\sum_{i,j} A_{ij} v_i v_j \right) = \sum_{i} k_i v_i^2 - \sum_{i,j} A_{ij} v_i v_j
\end{equation}
Questa è esattamente $\vec{v}^T L \vec{v}$.
Abbiamo quindi dimostrato:
\begin{equation}
	\vec{v}^T L \vec{v} = \frac{1}{2} \sum_{i,j} A_{ij} (v_i - v_j)^2
\end{equation}
Dall'equazione agli autovalori, $\vec{v}^T L \vec{v} = \lambda ||\vec{v}||^2$.
\begin{equation}
	\lambda = \frac{\frac{1}{2} \sum_{i,j} A_{ij} (v_i - v_j)^2}{||\vec{v}||^2}
\end{equation}
Poiché $A_{ij} \ge 0$ e $(v_i - v_j)^2 \ge 0$, il numeratore è non negativo. Il denominatore è positivo. Di conseguenza, $\lambda \ge 0$.
		
		
		\item \textbf{Autovalore Nullo:} Esiste almeno un autovalore uguale a 0.
		Consideriamo il vettore $\vec{v}_0 = (1, 1, ..., 1)^T$.
		\begin{equation}
		     (L \vec{v}_0)_i = \sum_{j} L_{ij} (1) = 0
		\end{equation}
       Questo perché ogni riga di $L$ somma a zero.
		Pertanto, $L \vec{v}_0 = 0 \cdot \vec{v}_0$. Ciò dimostra che $\lambda_1 = 0$ è un autovalore con autovettore corrispondente $\vec{v}_0$ (o la sua versione normalizzata $(1/\sqrt{n}, ..., 1/\sqrt{n})^T$).
		
		\item \textbf{Non Invertibile:} Poiché $L$ ha almeno un autovalore nullo, il suo determinante è 0 ($\det(L) = \prod \lambda_i = 0$). Ciò implica che $L$ non è invertibile.
		
		\item \textbf{Molteplicità di $\mathbf{\lambda=0}$:} La molteplicità algebrica dell'autovalore nullo corrisponde al numero di componenti disconnesse nel grafo.
		Se la rete ha $C$ componenti disconnesse, $L$ sarà una matrice a blocchi diagonali (dopo un riordinamento dei nodi). Ogni blocco rappresenta il Laplaciano per quella componente, e ogni componente possiede il proprio autovalore nullo. È possibile costruire $C$ autovettori linearmente indipendenti per $\lambda=0$ (es. $(1,1,..,0,0)$ per la componente 1, $(0,0,..,1,1)$ per la componente 2, ecc.).
		Se esiste un solo autovalore nullo (molteplicità 1), la rete è completamente connessa.
	\end{enumerate}
	
	\section{Random Walks}
	Consideriamo ora un processo dinamico su un grafo connesso non orientato: la passeggiata aleatoria (Random Walk). Il processo è definito come segue:
	\begin{itemize}
		\item Si inizia da un nodo casuale.
		\item A ogni passo temporale discreto, il camminatore situato sul nodo $j$ (con grado $k_j$) si sposta su uno dei suoi $k_j$ vicini.
		\item La scelta del vicino avviene uniformemente a caso, quindi la probabilità di spostarsi verso uno specifico vicino $i$ è $1/k_j$.
		\item Il camminatore non può rimanere sullo stesso nodo; deve necessariamente muoversi.
	\end{itemize}
	
	\subsection{Master Equation}
	Sia $P_i(t)$ la probabilità che il camminatore si trovi al nodo $i$ al tempo $t$.
	Per trovarsi al nodo $i$ al tempo $t$, il camminatore deve essere stato in un nodo vicino $j$ al tempo $t-1$ e poi essersi spostato da $j$ a $i$. La probabilità di questo evento è la somma su tutti i nodi $j$:
	\begin{equation}
	    P_i(t) = \sum_{j} \mathbb{P}(j \to i) P_j(t-1)
	\end{equation}
	La probabilità di transizione $\mathbb{P}(j \to i)$ è $1/k_j$ se $j$ è connesso a $i$ ($A_{ij}=1$), e 0 altrimenti.
	\begin{equation}
	    P_i(t) = \sum_{j=1}^{n} A_{ij} \frac{1}{k_j} P_j(t-1)
	\end{equation}
	La somma è su tutti i $j$, ma $A_{ij}$ (o $A_{ji}$ dato che non è orientato) la vincola ai soli vicini di $i$.
	
	\noindent In forma matriciale, sia $\vec{P}(t)$ il vettore delle probabilità:
	\begin{equation}
	    \vec{P}(t) = (A D^{-1}) \vec{P}(t-1)
	\end{equation}
	dove $D^{-1}$ è la matrice diagonale con elementi $(D^{-1})_{jj} = \frac{1}{k_j}$
	
	\subsection{Stato Stazionario}
	Cerchiamo la distribuzione stazionaria $\vec{P} = \vec{P}(\infty)$ per $t \to \infty$. 
    
    \noindent In questo limite, $\vec{P}(t) = \vec{P}(t-1) = \vec{P}$. L'equazione diventa un problema agli autovalori:
	\begin{equation}
	    \vec{P} = (A D^{-1}) \vec{P} \quad \longrightarrow \quad (\mathbb{I} - A D^{-1}) \vec{P} = 0
	\end{equation}
	Moltiplicando per $D$ a sinistra (assumendo $k_i \neq 0$ per ogni $i$):
    \begin{equation}
        D(\mathbb{I} - A D^{-1}) \vec{P} = (D - A) D^{-1} \vec{P} = 0 \quad \longrightarrow \quad (D^{-1} \vec{P}) = 0
    \end{equation}
	Questa equazione implica che il vettore $(D^{-1} \vec{P})$ sia l'autovettore del Laplaciano $L$ corrispondente all'autovalore $\lambda = 0$.
	Assumendo che il grafo sia connesso, esiste un solo autovettore per $\lambda=0$ (a meno di una costante $a$), che sappiamo essere $\vec{v}_0 = (1, 1, ..., 1)^T$.
	\begin{equation}
	    D^{-1} \vec{P} = a \vec{v}_0 = (a, a, ..., a)^T \quad \longrightarrow \quad \vec{P} = a D \vec{v}_0
	\end{equation}
	In componenti, questo corrisponde a:
	\begin{equation}
    P_i = a k_i
    \end{equation}
    La probabilità di trovarsi in uno specifico nodo $i$ nello stato stazionario è proporzionale al suo grado $k_i$.
	
	\noindent Troviamo la costante di normalizzazione $a$ imponendo $\sum P_i = 1$:
	\begin{equation}
	    \sum_{i} P_i = a \sum_{i} k_i = a (2m) = 1 \implies a = \frac{1}{2m}
	\end{equation}
	Dove $m$ è il numero totale di archi.
	
	\noindent La distribuzione di probabilità stazionaria è dunque:
     \begin{tcolorbox}[colback=yellow!25, colframe=yellow!75!orange, coltitle=black, title=\textbf{Distribuzione di Probabilità Stazionaria}]
\begin{equation}
	    P_i (\infty)= \frac{k_i}{\sum_j k_j} = \frac{k_i}{2m}
	\end{equation}
\end{tcolorbox}
	Questa rappresenta una sorta di \textbf{attaccamento preferenziale} nel tempo: \textbf{il camminatore aleatorio trascorre più tempo sui nodi altamente connessi}.
	
	\subsubsection{Probabilità di transizione nello Stato Stazionario}
	Per attraversare $i \to j$, il camminatore deve prima trovarsi al nodo $i$ (Prob. $P_i$) e poi scegliere di spostarsi verso $j$ (Prob. $1/k_i$).
	
	\begin{equation}
	    \mathbb{P}(i \to j) = P_i \cdot \frac{1}{k_i} = \left( \frac{k_i}{2m} \right) \cdot \frac{1}{k_i} = \frac{1}{2m}
	\end{equation}

     \begin{tcolorbox}[colback=yellow!25, colframe=yellow!75!orange, coltitle=black, title=\textbf{Risultato Notevole:}]
in una rete non orientata, la probabilità di attraversare un \textbf{qualsiasi} collegamento è identica e indipendente dai gradi dei nodi. Tutti gli archi hanno la stessa probabilità di essere attraversati.
	
\end{tcolorbox}
	%==================================================================
	% PARTE 3: MISURE DI CENTRALITÀ
	%==================================================================
	
	\section{Misure di Centralità}
	Le \textbf{misure di centralità} sono metodi per classificare i nodi in base alla loro "importanza". La definizione di importanza varia a seconda del contesto. Una misura semplice è la \textbf{Degree Centrality} (Centralità di Grado): si ordinano i nodi in base al loro grado $k_i$.

Esploreremo ora misure più complesse basate su equazioni autoconsistenti.

	
	\subsection{Eigenvector Centrality (Centralità di Autovettore)}
	L'idea fondamentale è che l'importanza di un nodo è determinata dall'importanza dei suoi vicini. Un nodo è importante se è connesso ad altri nodi importanti. Possiamo esprimere questo concetto tramite un'equazione autoconsistente. 
    
    Sia $x_i$ la centralità del nodo $i$.
	\begin{equation}
	    x_i = \frac{1}{\lambda} \sum_{j} A_{ij} x_j
	\end{equation}
	Qui $A_{ij}$ è la matrice di adiacenza (usando la convenzione $j \to i$) e $\lambda$ è una costante di proporzionalità.
	In forma matriciale, questa è l'equazione agli autovalori:
    \begin{equation}
        \lambda x_i = \sum_{j} A_{ij} x_j \implies \lambda \vec{x} = A \vec{x}
    \end{equation}
	Il vettore di centralità $\vec{x}$ è un autovettore della matrice di adiacenza $A$.
	
	\textit{Quale autovettore?} La matrice di adiacenza possiede $n$ autovettori. Necessitiamo di una misura di centralità in cui tutti gli $x_i$ siano reali e positivi (o almeno non negativi).
	
	Il \textbf{Teorema di Perron-Frobenius} fornisce la risposta:
    \begin{tcolorbox}[colback=colorC!15, colframe=colorC!70!colorB,  title=\textbf{Teorema di Perron-Frobenius:}]

		Per ogni matrice $A$ irriducibile e non negativa (come la matrice di adiacenza di un grafo connesso), il suo autovalore più grande $\lambda_1$ è reale e positivo. Le componenti dell'autovettore corrispondente $\vec{v}_1$ sono anch'esse tutte positive.

\end{tcolorbox}
	
	
	Pertanto, il vettore della \textbf{Eigenvector Centrality} $\vec{x}$ è definito come l'autovettore principale della matrice di adiacenza (quello corrispondente all'autovalore massimo $\lambda_1$).
    
    \hfill
    
	\noindent \textit{Dimostrazione:}
    
	Si considera un vettore casuale $\vec{x}(0)$ con tutte le componenti non negative.
	Iteriamo l'applicazione di $A$: $\vec{x}(t) = A \vec{x}(t-1) = A^t \vec{x}(0)$.
	Possiamo decomporre $\vec{x}(0)$ nella base degli autovettori $\vec{v}_i$ di $A$:
    \begin{equation}
        \vec{x}(0) = \sum_{i=1}^n c_i \vec{v}_i
    \end{equation}
    Applicando $A^t$ si ottiene:
    \begin{equation}
        \vec{x}(t) = A^t \sum_{i=1}^n c_i \vec{v}_i = \sum_{i=1}^n c_i (A^t \vec{v}_i) = \sum_{i=1}^n c_i \lambda_i^t \vec{v}_i
    \end{equation}
	Sia $\lambda_1$ l'autovalore massimo ($\lambda_1 > |\lambda_i|$ per $i \ge 2$). Possiamo fattorizzarlo:
    \begin{equation}
      \vec{x}(t) = \lambda_1^t \left( c_1 \vec{v}_1 + \sum_{i=2}^n c_i \left(\frac{\lambda_i}{\lambda_1}\right)^t \vec{v}_i \right)  
    \end{equation}
	Per $t \to \infty$, il termine $\left(\frac{\lambda_i}{\lambda_1}\right)^t \to 0$ per tutti gli $i \ge 2$.
	Di conseguenza, per grandi $t$, il vettore $\vec{x}(t)$ converge in direzione all'autovettore principale:
    \begin{equation}
        \vec{x}(t) \approx c_1 \lambda_1^t \vec{v}_1
    \end{equation}
	Poiché $A$ è una matrice non negativa e $\vec{x}(0)$ è non negativo, $\vec{x}(1) = A\vec{x}(0)$ deve essere non negativo. Per induzione, $\vec{x}(t)$ deve essere non negativo per ogni $t$.
	Essendo $\vec{x}(t)$ non negativo e $c_1 \lambda_1^t$ uno scalare, l'autovettore $\vec{v}_1$ deve avere tutte le componenti non negative. Il teorema completo garantisce che siano strettamente positive. Questo giustifica il suo utilizzo come misura di centralità.

\hfill 

   \noindent  \textbf{Limitazione della Eigenvector Centrality:} Questo metodo fallisce per i grafi orientati, specialmente quelli con o nodi senza collegamenti in ingresso.
	
	Si consideri un nodo A senza link in ingresso. La riga $i=A$ nella matrice di adiacenza (usando $A_{ij}$ come $j \to i$) sarà composta da zeri.
	Allora la centralità $x_A$ sarà:
	$$x_A = \frac{1}{\lambda_1} \sum_{j} A_{Aj} x_j = \frac{1}{\lambda_1} \sum_{j} (0) x_j = 0$$
	
	Se un nodo B ha un link in ingresso solo da A, anche la sua centralità $x_B$ sarà 0, poiché $x_A=0$. Questo problema si propaga, portando molti nodi ad avere centralità nulla semplicemente a causa della natura orientata della rete.
    
	\subsection{Katz Centrality}
	Per risolvere il problema della Eigenvector Centrality, la \textbf{Katz Centrality} introduce una piccola costante positiva $\beta$ per ogni nodo. Questo assicura una centralità di "base", anche per nodi senza collegamenti entranti. L'equazione diventa (usando $\alpha$ come parametro):
    \begin{equation}
        x_i = \alpha \sum_{j} A_{ij} x_j + \beta
    \end{equation}
    In forma matriciale, usando $\vec{1}$ come vettore unitario:
    \begin{equation}
        \vec{x} = \alpha A \vec{x} + \beta \vec{1}
    \end{equation}
	Possiamo risolvere per $\vec{x}$:
    \begin{equation}
       (\mathbb{I} - \alpha A) \vec{x} = \beta \vec{1} \quad 
	\longrightarrow \quad \vec{x} = \beta (\mathbb{I} - \alpha A)^{-1} \vec{1} 
    \end{equation}
	
	Il rapporto $\alpha / \beta$ è l'unico parametro rilevante, quindi spesso si pone $\beta = 1$. Il parametro $\alpha$ regola ora l'importanza relativa della componente dell'autovettore (struttura della rete) rispetto alla componente costante $\beta$ (importanza di base).
	\begin{itemize}
		\item Se $\alpha \to 0$, tutti i nodi hanno la stessa centralità $x_i = \beta$.
		\item Se $\beta \to 0$, recuperiamo la Eigenvector Centrality.
	\end{itemize}
	
	\subsubsection{Condizione di Invertibilità}
	La soluzione $\vec{x} = \beta (\mathbb{I} - \alpha A)^{-1} \vec{1}$ richiede che la matrice $(\mathbb{I} - \alpha A)$ sia invertibile.
	Ciò significa che il suo determinante deve essere non nullo:
	\begin{equation}
	  \det(\mathbb{I} - \alpha A) \neq 0 
	\end{equation}	
	Questo è equivalente a:
	\begin{equation}
	    \det(\alpha^{-1} \mathbb{I} - A) \neq 0
	\end{equation}
	Il determinante è zero se $\alpha^{-1}$ è un autovalore di $A$. Aumentando $\alpha$ da 0, le centralità divergeranno quando $\alpha^{-1}$ incontra l'autovalore \textit{massimo} $\lambda_1$ di $A$.
	
	Pertanto, dobbiamo scegliere $\alpha$ tale che:
	\begin{equation}
	    \alpha^{-1} > \lambda_1 \quad \implies \quad \alpha < \frac{1}{\lambda_1}
	\end{equation}
	Questo assicura che la matrice sia invertibile e che i punteggi di centralità siano finiti e positivi. 
    
\hfill

    \noindent \textbf{Limitazione della Katz Centrality:} Un nodo con centralità molto alta (come Amazon) che linka a milioni di altri nodi (es. venditori) trasferisce la sua intera centralità a tutti loro. Questo "sovraccarica" la centralità dei nodi puntati da un grande hub.
	
	\subsection{PageRank Centrality}
	La \textbf{PageRank Centrality} (che prende il nome da Larry Page) risolve il problema della Katz Centrality dividendo la centralità trasferita per il \textbf{grado uscente} ($k_j^{\text{out}}$) del nodo $j$ che sta inviando "l'importanza".
	
	L'equazione del PageRank è:
    \begin{equation}
        x_i = \alpha \sum_{j} A_{ij} \frac{x_j}{k_j^{\text{out}}} + \beta
    \end{equation}
	Il termine $k_j^{\text{out}}$ deve essere gestito se $k_j^{\text{out}} = 0$. Una correzione comune è usare $k_j^{\text{out}} = \max(k_j^{\text{out}}, 1)$
	
	\noindent Questa è una distribuzione "democratica" dell'importanza: il nodo $j$ divide la sua importanza $x_j$ equamente tra tutti i $k_j^{\text{out}}$ nodi a cui si collega.
	In forma matriciale:
    \begin{equation}
        \vec{x} = \alpha (A D_{\text{out}}^{-1}) \vec{x} + \beta \vec{1}
    \end{equation}
	dove $D_{\text{out}}^{-1}$ è la matrice diagonale di $1/k_j^{\text{out}}$
	
	La soluzione è:
	\begin{equation}
	    \vec{x} = \beta (\mathbb{I} - \alpha A D_{\text{out}}^{-1})^{-1} \vec{1}
	\end{equation}
	
	Nuovamente, dobbiamo assicurare che la matrice $(\mathbb{I} - \alpha A D_{\text{out}}^{-1})$ sia invertibile.
	Ciò richiede $\alpha < 1/\lambda_1(A D_{\text{out}}^{-1})$, dove $\lambda_1$ è l'autovalore massimo della matrice $W = A D_{\text{out}}^{-1}$ (la matrice di transizione della passeggiata aleatoria).
	
	\begin{itemize}
		\item Per una rete \textbf{non orientata}, $k_j^{\text{out}} = k_j$, e si può dimostrare che $\lambda_1(A D^{-1}) = 1$. La condizione diventa $\alpha < 1$.
		\item Per una rete \textbf{orientata}, non esiste una regola generale. Il valore usato da Google era notoriamente $\alpha = 0.85$.
	\end{itemize}
	
	\subsubsection{Esempio PageRank: Grafo a Percorso Non Orientato}
	Consideriamo la seguente rete:
    \begin{figure}[h]
        \centering
    \begin{tikzpicture}[scale=0.5]
    % Definizione dello stile per i nodi
    \tikzset{
        node_style/.style={
            circle,
            fill=colorC,      
            text=white,      % Colore del testo b
            minimum size=1cm, % Dimensione minima per uniformità
            inner sep=2pt
        },
        edge_style/.style={
            color=colorD,       
            line width=1.5pt % Spessore della linea
        }
    }

    % Posizionamento dei nodi per formare la 'V'
    \node[node_style] (2) at (0, 0) {2}; % Nodo centrale
    \node[node_style] (1) at (4, 1.5) {1}; % Nodo in alto a destra
    \node[node_style] (3) at (4, -1.5) {3}; % Nodo in basso a destra

    % Disegno degli archi
    \draw[edge_style] (1) -- (2);
    \draw[edge_style] (2) -- (3);

\end{tikzpicture}
    \end{figure}
	
	$$ A = \begin{pmatrix} 0 & 1 & 0 \\ 1 & 0 & 1 \\ 0 & 1 & 0 \end{pmatrix} \quad
	D = \begin{pmatrix} 1 & 0 & 0 \\ 0 & 2 & 0 \\ 0 & 0 & 1 \end{pmatrix} \quad
	D^{-1} = \begin{pmatrix} 1 & 0 & 0 \\ 0 & 1/2 & 0 \\ 0 & 0 & 1 \end{pmatrix} $$
	
	Dobbiamo calcolare $M = (\mathbb{I} - \alpha A D^{-1})$:
	$$M = \begin{pmatrix} 1 & 0 & 0 \\ 0 & 1 & 0 \\ 0 & 0 & 1 \end{pmatrix} - \alpha \begin{pmatrix} 0 & 1 & 0 \\ 1 & 0 & 1 \\ 0 & 1 & 0 \end{pmatrix} \begin{pmatrix} 1 & 0 & 0 \\ 0 & 1/2 & 0 \\ 0 & 0 & 1 \end{pmatrix} = \begin{pmatrix} 1 & -\alpha/2 & 0 \\ -\alpha & 1 & -\alpha \\ 0 & -\alpha/2 & 1 \end{pmatrix}$$
	
	Il determinante è $\det(M) = 1(1 - \alpha^2/2) - (-\alpha/2)(-\alpha) = 1 - \alpha^2/2 - \alpha^2/2 = 1 - \alpha^2$.
	Il determinante è 0 per $\alpha = \pm 1$. Come previsto per un grafo non orientato, dobbiamo scegliere $\alpha < 1$.
	
	Il vettore di centralità è $\vec{x} = \beta (M)^{-1} \vec{1}$. Ponendo $\beta=1$:
	$$M^{-1} = \frac{1}{\det(M)} (\text{Cof}(M))^T = \frac{1}{1-\alpha^2} \begin{pmatrix} 1-\alpha^2/2 & \alpha/2 & \alpha^2/2 \\ \alpha & 1 & \alpha \\ \alpha^2/2 & \alpha/2 & 1-\alpha^2/2 \end{pmatrix}$$
	
	Moltiplicando per $\vec{1} = (1, 1, 1)^T$ (ovvero sommando le righe di $M^{-1}$):
	$$x_1 = \frac{1}{1-\alpha^2} (1 - \alpha^2/2 + \alpha/2 + \alpha^2/2) = \frac{1 + \alpha/2}{1 - \alpha^2}$$
	$$x_2 = \frac{1}{1-\alpha^2} (\alpha + 1 + \alpha) = \frac{1 + 2\alpha}{1 - \alpha^2}$$
	$$x_3 = \frac{1}{1-\alpha^2} (\alpha^2/2 + \alpha/2 + 1 - \alpha^2/2) = \frac{1 + \alpha/2}{1 - \alpha^2}$$
	
	Da questo risultato osserviamo:
	\begin{itemize}
		\item $x_1 = x_3$, il che è ragionevole data la simmetria della rete.
		\item $x_2 > x_1$ (poiché $\alpha > 0$), quindi il nodo 2 possiede la centralità più alta.
		\item Questo ordinamento ($2 > 1 = 3$) corrisponde al grado dei nodi ($k_2=2, k_1=k_3=1$). Per reti più grandi e complesse, le classifiche ottenute tramite PageRank e semplice grado differiranno spesso.
	\end{itemize}

    
	\subsection{Betweenness Centrality}
	Questa misura segue una filosofia diversa, basata sul flusso di informazione e sui cammini minimi. Un nodo è importante se si trova su molti \textbf{cammini minimi} (geodetiche) tra altre coppie di nodi nella rete. Misura quanto un nodo sia "in mezzo" agli altri.

	\begin{enumerate}
		\item Per ogni coppia di nodi $(s, t)$ nella rete, si calcolano tutti i cammini minimi tra di essi (es. usando l'algoritmo di Dijkstra).
		\item Sia $g_{st}$ il numero totale di cammini minimi tra $s$ e $t$. (Può esserci più di un percorso della stessa lunghezza minima).
		\item Sia $n_{st}(i)$ il numero di quei cammini minimi che passano attraverso il nodo $i$.
		\item La betweenness centrality $x_i$ del nodo $i$ è la somma delle frazioni dei cammini minimi che passano per $i$, sommata su tutte le possibili coppie $(s, t)$ dove $s \neq t$:
        \begin{equation}
            x_i = \sum_{s \neq t} \frac{n_{st}(i)}{g_{st}}
        \end{equation}
		
	\end{enumerate}
	
	Per normalizzare questo valore, si può dividere per il numero totale di coppie, che è $n(n-1)$ (per grafi orientati) o $n(n-1)/2$ (per grafi non orientati). Una normalizzazione comune per grafi non orientati è:
    \begin{equation}
        x_i^{\text{norm}} = \frac{2}{n(n-1)} \sum_{s < t} \frac{n_{st}(i)}{g_{st}}
    \end{equation}
	
	Questa misura è computazionalmente onerosa ($O(nm)$ o $O(n^3)$) ma è molto efficace per identificare nodi critici o "colli di bottiglia" (bottlenecks) in una rete, ovvero nodi cruciali per la comunicazione o il trasporto.
	

\chapter{Lezione 9}
\label{chap:lezione_09} 

\begin{flushright}
\textit{Data: 29/10/2025}
\end{flushright}


\section{Transitività e Coefficiente di Clustering}
Una proprietà importante delle reti è la \textbf{transitività}, espressa colloquialmente per le reti sociali come: "l'amico del mio amico è mio amico". Questo concetto misura la tendenza dei nodi a formare triangoli (triadi).

\begin{itemize}
\item \textbf{Transitività Perfetta:} Si verifica nelle \textbf{cliques} (o cricche), che sono sottografi in cui tutti i nodi sono completamente connessi tra loro (es. una 3-clique è un triangolo).
\item \textbf{Transitività Parziale:} Il caso più generale in cui un amico di un amico non è necessariamente un amico.
\end{itemize}

Questa proprietà è quantificata dal \textbf{Coefficiente di Clustering ($C$)}.

\subsection{Coefficiente di Clustering Globale}
Il coefficiente di clustering globale $C$ è una misura riferita all'intera rete. È definito come la frazione di "triadi aperte" (cammini di lunghezza 2) che sono "chiuse" (ovvero, formano un triangolo).

\begin{equation}
C = \frac{\text{\# di cammini chiusi di lunghezza 2}}{\text{\# di cammini di lunghezza 2}}
\end{equation}
Alternativamente, poiché ogni triangolo contiene 6 cammini chiusi di lunghezza 2 (ad es., $1\to2\to3$, $1\to3\to2$, $2\to1\to3$, ecc.), questo può essere scritto come:

\begin{equation}
C = \frac{\text{\# di triangoli} \times 6}{\text{\# di cammini di lunghezza 2}}
\end{equation}

\begin{center}
\begin{tikzpicture}[scale=0.5]
    \tikzset{
        simple_node/.style={
            text=black
        },
        black_edge/.style={
            line width=1.5pt,
            colorB
        },
        red_dotted/.style={
            color=colorE,
            line width=2pt,
            loosely dotted, % Punti distanziati
            line cap=round  % Punti rotondi
        }
    }

    % Posizionamento nodi
    \node[simple_node] (2) at (0, -1) {2};       % Nodo sinistro
    \node[simple_node] (1) at (3, 2) {1};       % Nodo in alto
    \node[simple_node] (3) at (4.5, -0.8) {3};  % Nodo intermedio destro
    \node[simple_node] (4) at (4, -2.5) {4};    % Nodo in basso

    % Archi neri solidi
    \draw[black_edge] (2) -- (1);
    \draw[black_edge] (2) -- (3);
    \draw[black_edge] (2) -- (4);
    \draw[black_edge] (3) -- (4);

    % Linea rossa tratteggiata tra 1 e 3 (o tra 1 e il gruppo in basso)
    \draw[red_dotted] (1) -- (4);
    \draw[red_dotted] (1) -- (3);


\end{tikzpicture}
\end{center}

\subsubsection{Esempi}
\begin{enumerate}
\item Triangolo Completamente Connesso (3-clique):
I cammini di lunghezza 2 sono: $1\to2\to3$, $1\to3\to2$, $2\to1\to3$, $2\to3\to1$, $3\to1\to2$, $3\to2\to1$.
\begin{itemize}
\item \# di cammini di lunghezza 2 = 6.
\item Tutti e 6 questi cammini sono chiusi (es. per $1\to2\to3$, esiste il collegamento $1-3$).
\item $C = \frac{6}{6} = 1$.
\end{itemize}

\item Grafo a 4 Nodi (Triangolo + 1 nodo):
Consideriamo i nodi $\{i, j, k, l\}$, dove $\{i, k, l\}$ formano un triangolo, e $j$ è connesso solo a $i$.
% FIGURE PLACEHOLDER: Un grafo a triangolo (i, k, l) con un nodo extra (j) attaccato solo al nodo i
I cammini di lunghezza 2 sono:
\begin{itemize}
\item Da $j$: $j \to i \to k$, $j \to i \to l$ (2 cammini)
\item Da $k$: $k \to i \to j$, $k \to i \to l$, $k \to l \to i$ (3 cammini)
\item Da $l$: $l \to i \to j$, $l \to i \to k$, $l \to k \to i$ (3 cammini)
\item Da $i$: $i \to k \to l$, $i \to l \to k$ (2 cammini)
\end{itemize}
\begin{itemize}
\item Numero totale di cammini di lunghezza 2: $2+3+3+2 = 10$.
\item Solo i cammini all'interno del triangolo $\{i, k, l\}$ sono chiusi. Questi sono i 6 cammini dell'Esempio 1 (es. $i \to k \to l$ è chiuso da $i-l$, $k \to i \to l$ è chiuso da $k-l$, ecc.).
\item $C = \frac{6}{10} = \frac{3}{5}$.
\end{itemize}
\end{enumerate}

\subsubsection{Valori Empirici}
Le reti del mondo reale tendono ad avere alti coefficienti di clustering, ben lontani da quanto previsto in un grafo casuale semplice.
\begin{itemize}
\item Collaborazioni attori cinematografici: $C \approx 0.20$
\item Collaborazioni biologi: $C \approx 0.09$
\item Reti di email: $C \approx 0.16$
\end{itemize}
In un grafo casuale semplice e grande con $n$ nodi e grado medio $\langle k \rangle$, la probabilità che due nodi qualsiasi siano connessi è $p = \langle k \rangle / (n-1)$. La probabilità che due amici di un nodo siano anche amici tra loro è questa stessa $p$.
Pertanto, per un grafo casuale $C \approx p = \frac{\langle k \rangle}{n-1}$. Per $n \to \infty$, $C \to 0$. Ciò dimostra che i grafi casuali non sono buoni modelli per l'alto clustering osservato nelle reti reali.

\subsection{Coefficiente di Clustering Locale}
Possiamo definire anche un coefficiente di clustering $C_i$ per un singolo nodo $i$.
\begin{equation}
C_i = \frac{\text{\# di coppie di vicini di } i \text{ che sono connesse}}{\text{\# di coppie di vicini di } i}
\end{equation}
Il denominatore è il numero totale di possibili triangoli centrati in $i$, che è il numero di modi per scegliere 2 vicini tra $k_i$, ovvero $\binom{k_i}{2} = \frac{k_i(k_i-1)}{2}$.

\noindent Il numeratore è il numero di triangoli che esistono effettivamente coinvolgendo il nodo $i$. Questo può essere trovato usando la matrice di adiacenza. Un triangolo $i-j-k-i$ esiste se $A_{ij}=1$, $A_{jk}=1$, e $A_{ki}=1$. Il numero di triangoli connessi a $i$ è $\frac{1}{2}\sum_{j,k} A_{ij} A_{jk} A_{ki}$. Il fattore $\frac{1}{2}$ è presente perché la somma conta $i\to j\to k\to i$ e $i\to k\to j\to i$ come distinti.

\noindent Il coefficiente di clustering locale è:
\begin{equation}
C_i = \frac{\frac{1}{2}\sum_{j,k} A_{ij} A_{jk} A_{ki}}{\frac{1}{2}k_i(k_i-1)} = \frac{\sum_{j,k} A_{ij} A_{jk} A_{ki}}{k_i(k_i-1)}
\end{equation}

\noindent Il coefficiente di clustering globale $C$ può quindi essere definito come la media dei $C_i$ locali, pesata dal numero di possibili triangoli per ogni nodo:
\begin{equation}
C = \frac{\sum_i C_i \binom{k_i}{2}}{\sum_i \binom{k_i}{2}} = \frac{\sum_i \left( \frac{\sum_{j,k} A_{ij} A_{jk} A_{ki}}{k_i(k_i-1)} \right) \frac{k_i(k_i-1)}{2}}{\sum_i \frac{k_i(k_i-1)}{2}} = \frac{\sum_{i,j,k} A_{ij} A_{jk} A_{ki}}{\sum_i k_i(k_i-1)}
\end{equation}
Questa formula mostra che $C$ è calcolabile semplicemente dalla matrice di adiacenza $A$ (poiché $k_i = \sum_j A_{ij}$).

\section{Assortatività}
L'assortatività descrive la tendenza dei nodi a connettersi ad altri nodi che sono simili a loro. La "similitudine" può essere definita sulla base di varie caratteristiche.
\begin{itemize}
\item \textbf{Grado:} Questa è la più comune.
\begin{itemize}
\item \textbf{Assortativa:} Nodi ad alto grado si connettono a nodi ad alto grado (es. reti sociali).
\item \textbf{Disassortativa:} Nodi ad alto grado (hub) si connettono a nodi a basso grado (es. reti tecnologiche come Internet, o reti di interazione proteica).
\end{itemize}
\item \textbf{Caratteristiche non ordinate:} Nazionalità, genere, razza.
\item \textbf{Caratteristiche ordinate:} Età, reddito.
\end{itemize}

L'assortatività è una correlazione, quindi la misuriamo usando la covarianza. Sia $x_i$ il valore di una caratteristica per il nodo $i$. Il valore atteso $\mu$ di questa caratteristica, mediato sugli estremi di tutti i collegamenti, è:
\begin{equation}
\mu = \frac{1}{2m} \sum_i k_i x_i
\end{equation}
dove $2m = \sum_i k_i = \sum_{ij} A_{ij}$ è la somma di tutti i gradi.

\noindent La covarianza per questa caratteristica attraverso tutte le coppie di nodi connessi è:
\begin{equation}
\text{cov}(x, x') = \frac{1}{2m} \sum_{i,j} A_{ij} (x_i - \mu)(x_j - \mu)
\end{equation}
Il termine $A_{ij}$ assicura che sommiamo solo sui nodi connessi $i$ e $j$.

\subsection{Assortatività per Grado}
Analizziamo il caso specifico in cui la caratteristica è il grado stesso, $x_i = k_i$.
Il valore atteso $\mu$ diventa:
\begin{equation}
 \mu = \frac{1}{2m} \sum_i k_i (k_i) = \frac{1}{2m} \sum_i k_i^2
\end{equation}

\noindent Calcoliamo la covarianza.
\begin{align*}
\text{cov}(k, k') &= \frac{1}{2m} \sum_{i,j} A_{ij} (k_i - \mu)(k_j - \mu) \\
&= \frac{1}{2m} \sum_{i,j} A_{ij} (k_i k_j - \mu k_i - \mu k_j + \mu^2) \\
&= \frac{1}{2m} \left[ \sum_{i,j} A_{ij} k_i k_j - \mu \sum_{i,j} A_{ij} k_i - \mu \sum_{i,j} A_{ij} k_j + \mu^2 \sum_{i,j} A_{ij} \right]
\end{align*}
Semplifichiamo i termini:
\begin{itemize}
\item $\sum_{i,j} A_{ij} k_i = \sum_i k_i \left( \sum_j A_{ij} \right) = \sum_i k_i (k_i) = \sum_i k_i^2$
\item $\sum_{i,j} A_{ij} k_j = \sum_j k_j \left( \sum_i A_{ij} \right) = \sum_j k_j (k_j) = \sum_j k_j^2$
\item $\sum_{i,j} A_{ij} = 2m$
\end{itemize}
Sostituiamo questi termini:
\begin{align*}
\text{cov}(k, k') &= \frac{1}{2m} \left[ \sum_{i,j} A_{ij} k_i k_j - \mu \sum_i k_i^2 - \mu \sum_j k_j^2 + 2m \mu^2 \right] \\
\text{Usando } \sum_i &k_i^2 = 2m \mu \text{ (dalla def. di } \mu \text{)} \\
&= \frac{1}{2m} \left[ \sum_{i,j} A_{ij} k_i k_j - \mu (2m \mu) - \mu (2m \mu) + 2m \mu^2 \right] \\
&= \frac{1}{2m} \sum_{i,j} A_{ij} k_i k_j - \mu^2 - \mu^2 + \mu^2 \\
&= \frac{1}{2m} \sum_{i,j} A_{ij} k_i k_j - \mu^2
\end{align*}
Questa espressione può essere riscritta come:
\begin{equation}
 \text{cov}(k, k') = \frac{1}{2m} \sum_{i,j} \left( A_{ij} - \frac{k_i k_j}{2m} \right) k_i k_j
\end{equation}

Per ottenere un \textbf{coefficiente di assortatività $r$} normalizzato (correlazione di Pearson), dividiamo la covarianza per la massima covarianza possibile (la varianza):
\begin{equation}
r = \frac{\sum_{ij} (A_{ij} - \frac{k_i k_j}{2m}) k_i k_j}{\sum_{ij} (k_i \delta_{ij} - \frac{k_i k_j}{2m}) k_i k_j}
\end{equation}
Questo valore $r$ è limitato: $-1 \le r \le 1$.
\begin{itemize}
\item $r > 0$: Assortativo
\item $r < 0$: Disassortativo
\item $r = 0$: Non-assortativo
\end{itemize}
Questa formula è computazionalmente costosa ($O(n^2)$). Una formula equivalente più veloce ($O(m)$) è:
\begin{tcolorbox} [colback=yellow!25, colframe=yellow!75!orange, coltitle=black, title=\textbf{Coefficiente di Assortatività}]
\begin{equation}
r = \frac{S_1 S_e - S_2^2}{S_1 S_3 - S_2^2}
\end{equation}
dove $S_n = \sum_i k_i^n$ e $S_e = \sum_{ij} A_{ij} k_i k_j = 2 \sum_{\text{edges } i-j} k_i k_j$.
\end{tcolorbox}

\section{Modelli di Rete I: Grafi Casuali}
Passiamo ora ai modelli per generare reti. Il primo e più semplice è il grafo casuale, come definito da Erdős e Rényi. Ci concentreremo sul modello $G(n,p)$.

\begin{itemize}
\item \textbf{Ensemble $G(n,p)$:} Un insieme di grafi con $n$ nodi.
\item \textbf{Processo:} Per ogni possibile coppia di nodi, viene tracciato un arco con una probabilità fissa $p$.
\end{itemize}

\subsubsection{Proprietà di $G(n,p)$}
\begin{itemize}
\item \textbf{Numero Medio di Link $\langle m \rangle$:} Il numero di possibili coppie è $\binom{n}{2}$.
\begin{equation}
 \langle m \rangle = \binom{n}{2} p = \frac{n(n-1)}{2} p
 \end{equation}

\item \textbf{Grado Medio $\langle k \rangle$:} Il grado medio è $\langle k \rangle = \frac{2\langle m \rangle}{n}$.
\begin{equation}
 \langle k \rangle = \frac{2}{n} \left( \frac{n(n-1)}{2} p \right) = (n-1)p
 \end{equation}
\end{itemize}

\subsection{Distribuzione di Grado $p_k$}
La distribuzione di grado $p_k$ (la probabilità che un nodo abbia grado $k$) segue una \textbf{distribuzione binomiale}. Per un dato nodo, ci sono $n-1$ altri nodi a cui esso può connettersi.
\begin{equation}
p_k = \binom{n-1}{k} p^k (1-p)^{n-1-k}
\end{equation}
Per $n$ grande e $p$ piccolo (tale che $\langle k \rangle = (n-1)p$ rimanga finito), questa distribuzione binomiale converge a una \textbf{distribuzione di Poisson}.



\subsubsection{Derivazione: Da Binomiale a Poisson}
Vogliamo mostrare che $p_k \to \frac{\langle k \rangle^k e^{-\langle k \rangle}}{k!}$ quando $n \to \infty$.

\noindent Partiamo da $p_k = \frac{(n-1)!}{k!(n-1-k)!} p^k (1-p)^{n-1-k}$.

\begin{enumerate}
    \item Approssimiamo il coefficiente binomiale:
$$ \binom{n-1}{k} = \frac{(n-1)(n-2)\dots(n-k)}{k!} $$
Per $n \gg k$, ogni termine al numeratore è approssimativamente $n-1$.
$$ \binom{n-1}{k} \approx \frac{(n-1)^k}{k!} $$
\item Approssimiamo i termini di probabilità:
$$ p^k = \left( \frac{\langle k \rangle}{n-1} \right)^k = \frac{\langle k \rangle^k}{(n-1)^k} $$
$$ (1-p)^{n-1-k} = \left( 1 - \frac{\langle k \rangle}{n-1} \right)^{n-1-k} $$
Usando il limite $\lim_{x\to\infty} (1 - c/x)^x = e^{-c}$, e notando che $n-1-k \approx n-1$:
$$ (1-p)^{n-1-k} \approx \exp\left(-\langle k \rangle\right) $$
\item Combiniamo:
$$ p_k \approx \left( \frac{(n-1)^k}{k!} \right) \left( \frac{\langle k \rangle^k}{(n-1)^k} \right) \left( e^{-\langle k \rangle} \right) = \frac{\langle k \rangle^k e^{-\langle k \rangle}}{k!} $$
\end{enumerate}
 


\begin{tcolorbox}[colback=colorB!25,colframe=colorB!75!colorA,title={Distribuzione di Grado di G(n,p)}]
\begin{equation} \label{eq:poisson}
p_k = \frac{\langle k \rangle^k e^{-\langle k \rangle}}{k!}
\end{equation}
\end{tcolorbox}

Questa distribuzione ha un picco attorno a $\langle k \rangle$. Tutti i nodi sono statisticamente simili. Questo è un grafo "democratico", ma non riproduce gli "hub" (leggi di potenza) osservati in molte reti reali.

\subsubsection{Coefficiente di Clustering di $G(n,p)$}
Il coefficiente di clustering è la probabilità che due vicini di un nodo siano anche vicini tra loro. In $G(n,p)$, la probabilità che due nodi qualsiasi siano connessi è $p$.
\begin{equation}
 C = p = \frac{\langle k \rangle}{n-1}
\end{equation}
Per $n$ grande, $C \to 0$. I grafi casuali hanno un clustering molto basso.

\subsection{Diametro e Proprietà Small-World}
Il diametro $l$ (il cammino minimo più lungo) di un grafo $G(n,p)$ è molto piccolo.
Un argomento euristico:
\begin{itemize}
\item Il numero di primi vicini è $\sim \langle k \rangle$.
\item Il numero di secondi vicini è $\sim \langle k \rangle^2$.
\item Il numero di $l$-esimi vicini è $\sim \langle k \rangle^l$.
\end{itemize}
Il grafo è "coperto" quando questo numero è dell'ordine dei nodi totali $n$:
\begin{equation}
 \langle k \rangle^l \approx n \implies l \approx \frac{\log n}{\log \langle k \rangle}
\end{equation}
Il diametro di un grafo casuale cresce solo con il \textit{logaritmo} della sua dimensione. Questa è la \textbf{proprietà small-world}.

\subsection{La Componente Gigante}
Non è garantito che $G(n,p)$ sia completamente connesso. Può consistere in componenti disconnesse multiple. Siamo interessati alla dimensione della \textbf{Componente Gigante (GC)}, $S$, che è la frazione di nodi nella componente connessa più grande.

\subsubsection{Derivazione: Dimensione della Componente Gigante}
Sia $u$ la probabilità che un nodo scelto a caso \textbf{non} appartenga alla GC.

\noindent $S = 1 - u$ è la probabilità che vi \textit{appartenga}.

\noindent Affinché un nodo $i$ \textit{non} sia nella GC, deve fallire nel connettersi alla GC attraverso \textit{ciascuno} dei suoi potenziali $n-1$ vicini.
Consideriamo un altro nodo $j$:
\begin{itemize}
\item Prob. che $i$ e $j$ non siano connessi: $1-p$
\item Prob. che $i$ e $j$ siano connessi, MA $j$ non è nella GC: $p \cdot u$
\end{itemize}
La probabilità che $i$ non sia connesso alla GC \textit{tramite il nodo $j$} è la somma di questi due eventi mutualmente esclusivi:
$$ P(\text{non connesso tramite } j) = (1-p) + pu = 1 - p(1-u) $$
Poiché tutte le $n-1$ connessioni sono indipendenti, la probabilità $u$ che $i$ non sia connesso alla GC tramite \textit{nessun} nodo è:
$$ u = [1 - p(1-u)]^{n-1} $$
Sostituendo $p = \langle k \rangle / (n-1)$:
$$ u = \left[ 1 - \frac{\langle k \rangle}{n-1} (1-u) \right]^{n-1} $$
Usando il limite $\lim_{x\to\infty} (1 - c/x)^x = e^{-c}$, quando $n \to \infty$:
$$ u = \exp(-\langle k \rangle (1-u)) $$
Sostituendo $u = 1-S$:
$$ 1-S = \exp(-\langle k \rangle S) \implies S = 1 - \exp(-\langle k \rangle S) $$


\begin{tcolorbox} [colback=yellow!25, colframe=yellow!75!orange, coltitle=black, title=\textbf{Equazione di Autoconsistenza della Componente Gigante}]
\begin{equation} 
\label{eq:gc}
S = 1 - e^{-\langle k \rangle S}
\end{equation}
\end{tcolorbox}

\noindent Questa equazione è risolta graficamente trovando le intersezioni di $y=S$ e $y=1-e^{-\langle k \rangle S}$


\begin{figure}[htbp]
    \centering
    \begin{tikzpicture}
        % Definiamo la variabile Ssol all'inizio per usarla ovunque
        \pgfmathsetmacro{\Ssol}{0.5828}

        % =========================================
        % GRAFICO (a): Soluzione Grafica
        % =========================================
        \begin{axis}[
            name=plotA,
            width=7cm, height=7cm,
            xmin=0, xmax=1,
            ymin=0, ymax=1,
            xlabel={$S$},
            ylabel={$y$},
            tick pos=both,
            tick align=inside,
            xtick={0,0.2,0.4,0.6,0.8,1.0},
            ytick={0,0.2,0.4,0.6,0.8,1.0},
            clip=false
        ]
            \node[anchor=north west] at (rel axis cs:0.02, 0.98) {(a)};
            \addplot[dashed, domain=0:1, samples=2, gray!80] {x};

            \addplot[domain=0:1, samples=100, thick] {1 - exp(-0.5*x)} coordinate[pos=0.7] (c05_pos);
            \addplot[domain=0:1, samples=100, thick] {1 - exp(-1.0*x)} coordinate[pos=0.65] (c1_pos);
            \addplot[domain=0:1, samples=100, thick] {1 - exp(-1.5*x)};

            % Etichette
           \node[right, font=\footnotesize] (c5_label) at (axis cs: 0.6, 0.2) {$c=0.5$};
            \node[right, font=\footnotesize] (c1_label) at (axis cs: 0.7, 0.45) {$c=1$};
            \node[font=\footnotesize, anchor=east] (c15_label) at (axis cs: 0.3, 0.45) {$c=1.5$};
            

            % Punti rossi
            \node[circle, fill=colorE, inner sep=2.5pt] at (axis cs:0,0) {};
            
            % Definiamo il punto di intersezione e lo salviamo come 'intersectionA'
            \coordinate (intersectionA) at (axis cs:\Ssol, \Ssol);
            \node[circle, fill=colorE, inner sep=2.5pt] at (intersectionA) {};
            \draw[dotted] (intersectionA) -- (axis cs:\Ssol, 0);
        \end{axis}

        % =========================================
        % GRAFICO (b): Transizione di Fase
        % =========================================
        \begin{axis}[
            name=plotB,
            at={(plotA.south east)}, anchor=south west, xshift=1.5cm,
            width=8cm, height=7cm,
            xmin=0, xmax=3.5,
            ymin=0, ymax=1,
            xlabel={Average degree $c$},
            ylabel={Size of giant component $S$},
            yticklabel pos=right,
            ylabel near ticks,
            tick pos=both,
            tick align=inside,
            clip=false
        ]
            \node[anchor=north west] at (rel axis cs:0.02, 0.98) {(b)};
            
            % Curva transizione di fase
            \addplot[thick, domain=0.001:0.965, samples=200, variable=\t] ({-ln(1-\t)/\t}, {\t});
            \addplot[thick, domain=0:1] {0};

            % Definiamo i punti per la linea rossa QUI DENTRO, dove 'axis cs' funziona
            % Punto di partenza nel grafico B (c=1.5)
            \coordinate (pointB) at (axis cs: 1.5, \Ssol);
            % Punto di fine linea (estensione fino a c=2.8 per esempio)
            \coordinate (lineEnd) at (axis cs: 1, \Ssol);

           
        \end{axis}

        % =========================================
        % CONNESSIONI ESTERNE
        % =========================================
        % Ora usiamo solo coordinate già calcolate, niente matematica complessa qui fuori
        \draw[colorE!80, thick] (intersectionA) -- (pointB) -- (lineEnd);

    \end{tikzpicture}
    \caption{Soluzione dell'equazione della Componente Gigante nei grafi casuali. \textbf{(a)} Soluzione grafica dell'equazione di autoconsistenza $S = 1 - e^{-cS}$. Le curve mostrano il comportamento per diversi valori del grado medio $c = \langle k \rangle$. L'intersezione con la bisettrice $y=S$ fornisce le soluzioni. Oltre alla soluzione banale $S=0$, una soluzione non banale $S>0$  emerge quando $c > 1$. \textbf{(b)} Dimensione della componente gigante $S$ in funzione del grado medio $c$. La transizione di fase avviene a $c=1$; per $c \le 1$, non esiste una componente gigante ($S=0$), mentre per $c > 1$, una frazione finita dei nodi appartiene alla GC.}
    \label{fig:giant_component_transition}
\end{figure}

\begin{itemize}
\item $S=0$ è sempre una soluzione.
\item Una soluzione non banale $S > 0$ esiste solo se la pendenza di $y=1-e^{-\langle k \rangle S}$ in $S=0$ è maggiore di 1.
\item Pendenza in $S=0$: $\frac{d}{dS} \left( 1 - e^{-\langle k \rangle S} \right) \bigg|_{S=0} = \left( -e^{-\langle k \rangle S} \cdot (-\langle k \rangle) \right) \bigg|_{S=0} = \langle k \rangle$
\end{itemize}
Una GC non banale esiste solo se $\langle k \rangle > 1$. Questa è una \textbf{transizione di fase}.

\newpage
\section{Modelli di Rete II: Il Modello di Configurazione}
Il modello $G(n,p)$ genera una distribuzione di grado di Poisson. Cosa succede se vogliamo creare un grafo casuale con una \textit{specifica} distribuzione di grado $p_k$? Questo è lo scopo del \textbf{Modello di Configurazione}.

\subsubsection{Processo}
\begin{enumerate}
\item Si inizia con $n$ nodi.
\item Si assegna un grado $k_i$ a ciascun nodo $i$, estratto dalla $p_k$ desiderata. (Vincolo: $\sum k_i$ deve essere pari).
\item Si creano $k_i$ "stub" (mezzo arco) per ogni nodo $i$. Il numero totale di stub è $2m = \sum k_i$.
\item Si selezionano casualmente due stub dall'intero insieme e connetterli per formare un arco.
\item Si ipete finché tutti gli stub sono connessi.
\end{enumerate}

\subsection{Multi-archi e Auto-archi}
Questo processo non è "semplice" (nel senso della teoria dei grafi) perché può creare:
\begin{itemize}
\item \textbf{Auto-archi (loop):} Uno stub del nodo $i$ può essere connesso a un altro stub dello stesso nodo $i$.
\item \textbf{Multi-archi:} Due nodi $i$ e $j$ possono essere connessi da più di un arco.
\end{itemize}
La strategia più sicura è permettere che questi si formino e poi stimare la loro probabilità.

\begin{itemize}
\item \textbf{Prob. di un collegamento $P_{ij}$:}
\begin{equation}
 P_{ij} \approx \frac{k_i k_j}{2m-1} \approx \frac{k_i k_j}{2m}
 \end{equation}

\item \textbf{Numero atteso di Auto-Archi:}
\begin{equation}
 P_{\text{self}} = \sum_i \binom{k_i}{2} \frac{1}{2m-1} \approx \frac{1}{2m} \sum_i \frac{k_i(k_i-1)}{2} = \frac{\langle k^2 \rangle - \langle k \rangle}{2 \langle k \rangle}
 \end{equation}

\item \textbf{Numero atteso di Multi-Archi:} \color{red}{VEDERE PASSSAGGI!!}
\begin{equation}
 P_{\text{double\_edges}} \approx \frac{1}{2} \sum_{i \neq j} \binom{k_i}{2} \binom{k_j}{2} \left( \frac{1}{2m-1} \right)^2 \dots \approx \frac{1}{2} \left( \frac{\langle k^2 \rangle - \langle k \rangle}{\langle k \rangle} \right)^2
 \end{equation}
\end{itemize}
Il numero di questi archi "problematici" rimane finito (e quindi la loro densità svanisce per $n \to \infty$) \textbf{se} i momenti $\langle k \rangle$ e $\langle k^2 \rangle$ sono finiti. Se $p_k$ è una legge di potenza $p_k \sim k^{-\alpha}$ con $\alpha < 3$, allora $\langle k^2 \rangle$ diverge, e questo modello risulta "patologico" con un numero divergente di auto-loop.

\subsection{Assortatività e Paradosso dell'Amicizia}
Possiamo usare il modello di configurazione per guardare all'assortatività da una prospettiva diversa.


La distribuzione di grado dei miei vicini, $P_{nn}(k')$, è la stessa della distribuzione di grado generale $P(k')$?
No. Un nodo isolato ($k=0$) può esistere in $P(k)$, ma $P_{nn}(k=0) = 0$ per definizione (un vicino non può essere isolato).

\subsubsection{Derivazione: Distribuzione di Grado dei Vicini}
Vogliamo la probabilità $P_{nn}(k')$ che un nodo "all'altra estremità" di un arco scelto a caso abbia grado $k'$.
\begin{itemize}
\item Il numero totale di stub nella rete è $2m = N \langle k \rangle$.
\item Il numero di stub attaccati ai nodi di grado $k'$ è:
$$k' \cdot(\text{\# di nodi con grado } k')~=~k' N p_{k'}$$
\item La probabilità di scegliere uno stub che appartiene a un nodo di grado $k'$ è:
$$ P_{nn}(k') = \frac{\text{\# stub su nodi } k'}{\text{stub totali}} = \frac{k' N p_{k'}}{N \langle k \rangle} = \frac{k' p_{k'}}{\langle k \rangle} $$
\end{itemize}



\begin{equation} \label{eq:nn_dist}
P_{nn}(k') = \frac{k' p_{k'}}{\langle k \rangle}
\end{equation}


\subsubsection{Derivazione: Paradosso dell'Amicizia}
\begin{align*}
\langle k \rangle_{nn} &= \sum_{k'} k' P_{nn}(k') = \sum_{k'} k' \left( \frac{k' p_{k'}}{\langle k \rangle} \right) \\
&= \frac{1}{\langle k \rangle} \sum_{k'} (k')^2 p_{k'} \\
&= \frac{\langle k^2 \rangle}{\langle k \rangle}
\end{align*}
Il "paradosso" è il fatto che $\langle k \rangle_{nn} \ge \langle k \rangle$:

$$ \langle k \rangle_{nn} - \langle k \rangle = \frac{\langle k^2 \rangle}{\langle k \rangle} - \langle k \rangle = \frac{\langle k^2 \rangle - \langle k \rangle^2}{\langle k \rangle} = \frac{\sigma_k^2}{\langle k \rangle} $$
Poiché la varianza $\sigma_k^2 \ge 0$ e $\langle k \rangle > 0$, la differenza è non negativa.


\begin{tcolorbox}[colback=colorA!15,colframe=colorB!75!colorA,title=Paradosso dell'Amicizia]
Il grado medio di un vicino è maggiore o uguale al grado medio di un nodo.
\begin{equation} \label{eq:knn}
\langle k \rangle_{nn} = \frac{\langle k^2 \rangle}{\langle k \rangle} \ge \langle k \rangle
\end{equation}
"I tuoi amici sono, in media, più popolari di te."
\end{tcolorbox}
Questo accade perché i nodi ad alto grado sono sovra-rappresentati nell'insieme dei vicini (sono, per definizione, vicini di più nodi).

\subsection{Distribuzione del Grado in Eccesso}
Per processi dinamici come le epidemie, abbiamo bisogno della \textbf{distribuzione del grado in eccesso} $q_k$. Questa è la probabilità che un vicino, raggiunto tramite un arco, abbia $k$ \textit{altri} archi (in eccesso) su cui diffondere.
Ciò corrisponde al fatto che il vicino abbia un grado \textit{totale} di $k+1$.
\begin{equation} \label{eq:excess_deg}
q_k = P_{nn}(k+1) = \frac{(k+1) p_{k+1}}{\langle k \rangle}
\end{equation}
Questa distribuzione è propriamente normalizzata: $$\sum_{k=0}^\infty q_k = \frac{1}{\langle k \rangle} \sum_{k=0}^\infty (k+1) p_{k+1} = \frac{1}{\langle k \rangle} \sum_{k'=1}^\infty k' p_{k'} = \frac{\langle k \rangle}{\langle k \rangle} = 1$$

\subsection{Clustering nel Modello di Configurazione}
Il coefficiente di clustering globale $C$ può essere stimato come la probabilità che due vicini qualsiasi di un nodo scelto a caso siano essi stessi connessi.
Questa è la probabilità che due stub in eccesso scelti a caso da due vicini si connettano tra loro.
\begin{equation}
 C = \sum_{k_i, k_j=0}^\infty q_{k_i} q_{k_j} P(\text{link tra } i \text{ e } j)
\end{equation}
Usando $P_{ij} \approx k_i k_j / (2m)$:
\begin{equation}
 C \approx \sum_{k_i=0}^\infty \sum_{k_j=0}^\infty q_{k_i} q_{k_j} \frac{k_i k_j}{2m} = \frac{1}{2m} \left( \sum_k k q_k \right)^2
\end{equation}
Usando $\sum k q_k = (\langle k^2 \rangle - \langle k \rangle) / \langle k \rangle$ e $2m = N \langle k \rangle$:
\begin{equation}
C_{\text{config}} = \frac{1}{N \langle k \rangle} \left( \frac{\langle k^2 \rangle - \langle k \rangle}{\langle k \rangle} \right)^2
\end{equation}
Come in $G(n,p)$, il coefficiente di clustering per il modello di configurazione scala come $1/N$ e svanisce per reti grandi (assumendo momenti finiti). Non è un buon modello per le reti reali, altamente clusterizzate.

\section{Modelli di Rete III: Small-World (Watts-Strogatz)}
Questo modello, pubblicato nel 1998, interpola tra un reticolo regolare e un grafo casuale.
\begin{enumerate}
\item Inizio ($p=0$): Un reticolo regolare 1D dove ogni nodo è connesso ai suoi primi vicini. Questo grafo ha un \textbf{alto $C$} e un \textbf{alto $L$} (diametro, $L \sim N$).
\item Con probabilità $p$, ogni arco viene ricablato (rewire)  staccando un'estremità e riattaccandola a un nodo casuale nella rete.
\item Fine ($p=1$): Un grafo casuale $G(n, \langle k \rangle)$. Questo grafo ha un \textbf{basso $C$} e un \textbf{basso $L$} ($L \sim \log N$).
\end{enumerate}

\begin{figure}[htbp]
    \centering
\begin{tikzpicture}[scale=0.6]
    % --- CONFIGURAZIONE ---
    \tikzset{
        node_style/.style={circle, fill=black, inner sep=2pt},
        edge_style/.style={thick, line cap=round},
        curved_edge/.style={thick, bend right=75} 
    }

    % Numero di nodi
    \def\N{20}
    % Raggio
    \def\R{2.5}

    % ===================================================
    % 1. REGULAR (Reticolo Regolare)
    % ===================================================
    \begin{scope}[xshift=0cm]
        \node[] at (0, 3.8) {Regular};
        
        \foreach \i in {1,...,\N} {
            \node[node_style] (n\i) at ({90 - 360/\N * (\i - 1)}:\R) {};
        }
        
        \foreach \i in {1,...,\N} {
            \pgfmathsetmacro{\next}{int(mod(\i, \N) + 1)}
            \pgfmathsetmacro{\nextnext}{int(mod(\i + 1, \N) + 1)}
            
            \draw[edge_style] (n\i) -- (n\next);
            \draw[curved_edge] (n\i) to (n\nextnext);
        }
    \end{scope}

    % ===================================================
    % 2. SMALL-WORLD (Piccolo Mondo)
    % ===================================================
    \begin{scope}[xshift=7.5cm]
        \node[] at (0, 3.8) {Small-world};
        
        \foreach \i in {1,...,\N} {
            \node[node_style] (m\i) at ({90 - 360/\N * (\i - 1)}:\R) {};
        }
        
        \foreach \i in {1,...,\N} {
            \pgfmathsetmacro{\next}{int(mod(\i, \N) + 1)}
            \pgfmathsetmacro{\nextnext}{int(mod(\i + 1, \N) + 1)}
            
            \draw[edge_style] (m\i) -- (m\next);
            
            % Rimuoviamo il collegamento uscente dai nodi 3, 8 e 15
            \ifnum \i=3 \else \ifnum \i=8 \else \ifnum \i=15 \else
                \draw[curved_edge] (m\i) to (m\nextnext);
            \fi\fi\fi
        }
        
        % Shortcuts
        \draw[edge_style] (m3) to[bend right=10] (m12); 
        \draw[edge_style] (m8) to[bend left=10] (m18);  
        \draw[edge_style] (m15) to[bend right=20] (m5); 
    \end{scope}

    % ===================================================
    % 3. RANDOM (Casuale)
    % ===================================================
    \begin{scope}[xshift=15cm]
        \node[] at (0, 3.8) {Random};
        
        % Creazione nodi r1...r20
        \foreach \i in {1,...,\N} {
            \node[node_style] (r\i) at ({90 - 360/\N * (\i - 1)}:\R) {};
        }
        
        % Lista degli archi compattata per evitare problemi di spazi
        \foreach \u/\v in {
            1/2, 2/3, 3/4, 4/5, 5/6, 6/7, 7/8, 8/9, 9/10,
            10/11, 11/12, 12/13, 13/14, 14/15, 15/16, 16/17, 17/18, 18/19, 19/20, 20/1,
            1/6, 2/14, 3/10, 4/18, 5/12, 7/15, 8/19, 9/4, 11/17, 13/20, 16/2, 6/13, 10/18%
        } {
            % FIX: eval(u) e eval(v) rimuovono eventuali spazi fantasma
            \pgfmathsetmacro{\U}{int(\u)}
            \pgfmathsetmacro{\V}{int(\v)}
            \draw[edge_style] (r\U) to[bend right=15] (r\V);
        }
        
        \draw[edge_style] (r1) to (r11);
        \draw[edge_style] (r5) to (r15);
    \end{scope}

    % ===================================================
    % BARRA INFERIORE
    % ===================================================
    \begin{scope}[yshift=-4.5cm]
        \draw[->, >=stealth, very thick] (0, 0) -- (15, 0);
        \node[font=\large, anchor=east] at (-0.2, 0) {$p=0$};
        \node[font=\large, anchor=west] at (15.2, 0) {$p=1$};
        \node[] at (7.5, -0.5) {Increasing randomness};
    \end{scope}

\end{tikzpicture}
\end{figure}

\hfill \\


La scoperta chiave è ciò che accade per valori intermedi di $p$:
\begin{itemize}
\item La lunghezza del cammino $L(p)$ cala \textit{molto velocemente} (logaritmicamente) non appena $p > 0$. Le prime poche "scorciatoie" riducono drasticamente il diametro.
\item Il clustering $C(p)$ cala \textit{lentamente} (linearmente) con $p$.
\end{itemize}
Questo modello mostra un ampio regime "small-world" ($p$ piccolo) dove $L(p)$ è piccolo (come un grafo casuale) ma $C(p)$ è alto (come un reticolo regolare), corrispondendo ai dati empirici di molte reti reali.


\hfill \\

\begin{figure}[htbp]
    \centering
    \begin{tikzpicture}[scale=0.95]
        \begin{semilogxaxis}[
            width=10cm, height=8cm,
            xmin=0.0001, xmax=1,
            ymin=0, ymax=1.05,
            xlabel={$p$},
            ylabel={Characteristic path length},
            % Rimuove i numeri sull'asse Y per metterli come nell'immagine o lascia standard
            ytick={0, 0.2, 0.4, 0.6, 0.8, 1.0},
            % Gestione dei tick dell'asse X (potenze di 10)
            xtick={0.0001, 0.001, 0.01, 0.1, 1},
            xticklabels={0.0001, 0.001, 0.01, 0.1, 1},
            % Permette di disegnare fuori dal grafico (per le annotazioni)
            clip=false,
            % Stile dei tick
            tick pos=left,
            tick align=inside,
            major tick length=0.15cm,
            minor tick length=0.08cm,
        ]

            % ==========================
            % SERIE DATI: L(p)/L(0) - Cerchi neri
            % ==========================
            \addplot[only marks, mark=*, mark size=2pt, color=black!90] coordinates {
                (0.0001, 0.95)
                (0.0002, 0.85)
                (0.0004, 0.75)
                (0.0008, 0.55)
                (0.0015, 0.38)
                (0.003, 0.27)
                (0.006, 0.20)
                (0.01, 0.16)
                (0.02, 0.13)
                (0.05, 0.11)
                (0.1, 0.09)
                (0.2, 0.08)
                (0.5, 0.075)
                (1.0, 0.07)
            };
            
            % Etichetta interna L(p)/L(0)
            \node[anchor=west] at (axis cs: 0.0005, 0.35) {\large $\frac{L(p)}{ L(0)}$};


            % ==========================
            % SERIE DATI: C(p)/C(0) - Quadrati vuoti
            % ==========================
            \addplot[only marks, mark=square, mark size=2pt, color=black!90] coordinates {
                (0.0001, 1.0)
                (0.0003, 1.0)
                (0.0007, 1.0)
                (0.0015, 1.0)
                (0.003, 0.99)
                (0.006, 0.98)
                (0.01, 0.95)
                (0.02, 0.91)
                (0.05, 0.83)
                (0.1, 0.68)
                (0.2, 0.45)
                (0.4, 0.15)
                (1.0, 0.02)
            };

            % Etichetta interna C(p)/C(0)
            \node[anchor=west] at (axis cs: 0.004, 0.85) {\large $\frac{C(p)}{ C(0})$};


            % ==========================
            % ANNOTAZIONE BLU (In alto)
            % ==========================
            % Nodo per il testo\node[align=left, font=\small, color=black] (blueText) at (axis cs: 0.05, 1.25) {C(p) remains almost constant at its value for the regular lattice, indicating\\ that the transition to a small world is almost undetectable at the local level. };
            % Linea blu sotto il testo \draw[blue, thick] (blueText.south west) -- (blueText.south east);
            % Freccia blu \draw[->, >=stealth, blue, thick] (blueText.south) -- (axis cs: 0.05, 0.85);


            % ==========================
            % ANNOTAZIONE ROSSA (In basso)
            % ==========================
            % Nodo per il testo \node[align=left, font=\small, color=black, anchor=north] (redText) at (axis cs: 0.004, -0.25) {Rapid drop in L(p), corresponding to the onset of the small-world phenomenon.};
            % Linea rossa sopra il testo\draw[red!80!orange, thick] (redText.north west) -- (redText.north east);
            % Freccia rossa \draw[->, >=stealth, red!80!orange, thick] (redText.north) -- (axis cs: 0.002, 0.25);


            % ==========================
            % ETICHETTA ASSE DESTRO
            % ==========================
            \node[rotate=-90, anchor=south] at (rel axis cs: 1.05, 0.5) {Clustering coeffiecinet};

        \end{semilogxaxis}
    \end{tikzpicture}
    \caption{Il grafico mostra la lunghezza del cammino caratteristica normalizzata $L(p)/L(0)$ e il coefficiente di clustering normalizzato $C(p)/C(0)$ in funzione della probabilità di ri-cablaggio $p$. Si nota un ampio intervallo di $p$ dove la rete mantiene un alto clustering (tipico dei reticoli regolari) ma acquisisce una bassa lunghezza di cammino (tipica dei grafi casuali).}
    \label{fig:watts_strogatz_chart}
\end{figure}

\newpage
\section{Modelli di Rete IV: Modello Barabási-Albert}
Questo modello (1999) introduce un concetto diverso: la rete non è statica ma \textbf{cresce} nel tempo. Ha due ingredienti chiave:

\begin{enumerate}
\item \textbf{Crescita:} Si inizia con un piccolo insieme di $n_0$ nodi. Ad ogni passo temporale, si aggiunge un nuovo nodo al sistema.
\item \textbf{Attaccamento Preferenziale:} Il nuovo nodo si connette a $m_0$ nodi esistenti. La probabilità $\mathbb{P}(k_i)$ di connettersi a un nodo esistente $i$ non è uniforme; è proporzionale al grado attuale di quel nodo $k_i$.
\end{enumerate}
\begin{tcolorbox}[colback=yellow!25, colframe=yellow!75!orange, coltitle=black, title=\textbf{Regola dell'Attaccamento Preferenziale}]
\begin{equation} \label{eq:ba}
\mathbb{P}(k_i) = \frac{k_i}{\sum_j k_j}
\end{equation}
\end{tcolorbox}

Questo meccanismo rich get richer implica che i nodi che sono già altamente connessi hanno maggiori probabilità di ottenere nuove connessioni. Questo modello genera naturalmente una distribuzione di grado a legge di potenza (che risolveremo nella prossima lezione) e, per definizione, crea una singola componente gigante ($S=1$).


\chapter{Lezione 10}
\label{chap:lezione_10} 

\begin{flushright}
\textit{Data: 30/10/2025}
\end{flushright}

\section{Modelli di Crescita delle Reti: Il Modello Barabási-Albert}

Si analizzano i modelli dinamici per la generazione di reti, con particolare riferimento al modello Barabási-Albert (BA), introdotto nel 1999 per spiegare l'emergere delle proprietà \textit{Scale-Free} nelle reti reali.

\subsection{Meccanismo di Attaccamento Preferenziale}

Il modello si fonda sul meccanismo dell'\textit{Attaccamento Preferenziale} (Preferential Attachment), riassumibile nel principio "the rich get richer". 
La dinamica di evoluzione della rete è definita come segue:
\begin{enumerate}
    \item \textbf{Condizione Iniziale:} Si parte da un nucleo di $n_0$ nodi esistenti.
    \item \textbf{Crescita:} Ad ogni passo temporale $t$, viene aggiunto un nuovo nodo alla rete.
    \item \textbf{Connessione:} Il nuovo nodo crea $m$ collegamenti ($m \le n_0$) verso nodi già presenti nel sistema.
    \item \textbf{Regola di Attaccamento:} La probabilità $\mathbb{P}(k_i)$ che un nuovo collegamento si attacchi al nodo $i$-esimo dipende linearmente dal grado $k_i$ di tale nodo.
\end{enumerate}

Formalmente, la probabilità di connessione è data da:
\begin{equation}
    \mathbb{P}(k_i) = \frac{k_i}{\sum_j k_j}
\end{equation}
Poiché ogni nuovo nodo porta $m$ nuovi archi, il denominatore rappresenta la somma totale dei gradi nella rete al tempo $t$. 
Notiamo che:
\begin{itemize}
    \item Il numero di nodi cresce linearmente: $n(t) = n_0 + t$.
    \item Il numero totale di archi al tempo $t$ è $m_{tot} = m \cdot t$ (assumendo $t$ grande e trascurando i collegamenti iniziali).
    \item La somma dei gradi è il doppio del numero di archi: $\sum_j k_j = 2 m t$.
\end{itemize}
Per definizione, questo modello genera una rete con una Componente Gigante (GC) unitaria.

\begin{figure}[h]
    \centering
    \begin{tikzpicture}
        % Nodo centrale Hub
        \node[circle, draw=black, fill=colorC!40, minimum size=1.2cm] (hub) at (0,0) {\textbf{Hub}};
        
        % Nodi periferici esistenti
        \foreach \angle in {0, 60, 120, 180, 240, 300} {
            \node[circle, draw=black, fill=colorA!40, minimum size=0.5cm] (n\angle) at (\angle:2cm) {};
            \draw (hub) -- (n\angle);
        }
        
        % Nuovo nodo
        \node[circle, draw=black, fill=colorE!40, minimum size=0.8cm, label=above:New Node] (new) at (4,1) {};
        
        % Frecce di probabilità
        \draw[->, dashed, thick, colorE] (new) -- (hub) node[midway, above, sloped] {High $P$};
        \draw[->, dashed, thin, gray] (new) -- (n0) node[midway, below, sloped] {Low $P$};
    \end{tikzpicture}
    \caption{Rappresentazione schematica dell'attaccamento preferenziale. Il nuovo nodo ha una probabilità maggiore di connettersi all'Hub (grado elevato) rispetto ai nodi periferici.}
\end{figure}

\subsection{Approssimazione nel Continuo}

Per derivare la distribuzione dei gradi, trattiamo $k_i$ come una variabile continua nel tempo. L'evoluzione temporale del grado del nodo $i$ è descritta dall'equazione differenziale:
\begin{equation}
    \frac{\partial k_i}{\partial t} = m \cdot \mathbb{P}(k_i) = m \frac{k_i}{\sum_j k_j}
\end{equation}
Sostituendo la somma dei gradi $\sum_j k_j = 2mt$:
\begin{equation}
    \frac{\partial k_i}{\partial t} = m \frac{k_i}{2mt} = \frac{k_i}{2t}
\end{equation}
Separando le variabili e integrando tra il tempo di nascita del nodo $t_i$ e il tempo corrente $t$:
\begin{equation}
    \int_{k_i(t_i)}^{k_i(t)} \frac{dk_i}{k_i} = \int_{t_i}^{t} \frac{dt'}{2t'}
\end{equation}
La condizione iniziale è che al momento della nascita $t_i$, il nodo ha grado $m$, ovvero $k_i(t_i) = m$.
\begin{equation}
    \ln\left(\frac{k_i(t)}{m}\right) = \frac{1}{2} \ln\left(\frac{t}{t_i}\right)
\end{equation}
Esponenziando entrambi i membri otteniamo la legge di potenza temporale:
\begin{equation}
    k_i(t) = m \left( \frac{t}{t_i} \right)^{\beta} \quad \text{con } \beta = \frac{1}{2}
\end{equation}

\subsection{Derivazione della Distribuzione del Grado $P(k)$}

La probabilità che un nodo abbia un grado $k_i(t)$ inferiore a un valore soglia $k$ è equivalente alla probabilità che il suo tempo di nascita $t_i$ sia successivo a un certo istante (poiché i nodi più vecchi hanno grado maggiore).
Dalla relazione precedente ricaviamo $t_i$:
\begin{equation}
    k = m \left( \frac{t}{t_i} \right)^{1/2} \implies t_i = t \left( \frac{m}{k} \right)^2
\end{equation}
Assumiamo che i tempi di inserimento $t_i$ siano uniformemente distribuiti nell'intervallo $[0, t]$. La funzione di distribuzione cumulativa per $t_i$ è $F(t_i) = t_i / t$.
La probabilità cumulativa per il grado è:
\begin{equation}
    P(k_i(t) < k) = P\left( t_i > \frac{m^2 t}{k^2} \right) = 1 - P\left( t_i \le \frac{m^2 t}{k^2} \right)
\end{equation}
Sostituendo la distribuzione uniforme (assumendo $n_0 \ll t$):
\begin{equation}
    P(k_i(t) < k) = 1 - \frac{1}{t} \left( \frac{m^2 t}{k^2} \right) = 1 - \frac{m^2}{k^2}
\end{equation}
La densità di probabilità $P(k)$ si ottiene derivando la cumulativa rispetto a $k$:
\begin{equation}
    P(k) = \frac{\partial}{\partial k} \left( 1 - \frac{m^2}{k^2} \right) = -(-2) m^2 k^{-3} = 2m^2 k^{-3}
\end{equation}
Si ottiene quindi il celebre risultato che la distribuzione dei gradi decade come una legge di potenza con esponente $\gamma = 3$:
\begin{tcolorbox}[colback=colorC!15, colframe=colorC!70!colorB,  title=\textbf{Decadimento della Distribuzione dei Gradi a Legge di Potenza}]
\begin{equation}
    P(k) \sim k^{-3}
    \label{eq: 10.1}
\end{equation}
\end{tcolorbox}

\subsection{In-Degree Distribution}

Per analizzare più in dettaglio la struttura della rete, è utile distinguere tra i collegamenti uscenti ed entranti. Nel modello Barabási-Albert, ogni nuovo nodo introdotto nella rete genera $m$ collegamenti verso nodi preesistenti. Pertanto, il grado uscente (out-degree) è fisso per ogni nodo e pari a $m$.
La variabilità del grado totale $k$ è determinata interamente dal grado entrante (in-degree), che denotiamo con $q$.

La relazione tra il grado totale $k$ e il grado entrante $q$ è data da:
\begin{equation}
    k = q + m
\end{equation}
Dove $m$ rappresenta il numero di link che il nodo stabilisce al momento della sua nascita (out-degree costante).

Consideriamo la probabilità di attaccamento preferenziale. La probabilità $\mathbb{P}(k_i)$ che un nuovo nodo si connetta al nodo $i$ è proporzionale al grado totale di $i$. Possiamo riscrivere questa probabilità in termini di in-degree $q_i$:
\begin{equation}
    \mathbb{P}(q_i) = \frac{k_i}{\sum_j k_j} = \frac{q_i + m}{\sum_j (q_j + m)}
\end{equation}
Al denominatore, la somma è estesa a tutti i nodi presenti al tempo $t$. Se al tempo $t$ ci sono $t$ nodi (assumendo $n_0 \ll t$), la somma dei gradi totali è pari a due volte il numero totale di archi ($m \cdot t$):
\begin{equation}
    \sum_j k_j = 2 m t
\end{equation}
Pertanto, la probabilità che un singolo nuovo link si connetta a un nodo con in-degree $q$ è:
\begin{equation}
    \mathbb{P}(q) = \frac{q + m}{2 m t}
\end{equation}
Poiché ogni nuovo nodo introduce $m$ nuovi collegamenti, la probabilità complessiva che un nodo con in-degree $q$ riceva un nuovo collegamento durante un passo temporale è:
\begin{equation}
    P(\text{link} \to q) = m \cdot \mathbb{P}(q) = m \frac{q + m}{2 m t} = \frac{q + m}{2 t}
\end{equation}

\subsection{Approccio della Master Equation}

L'approccio della Master Equation ci permette di descrivere l'evoluzione temporale esatta della distribuzione dei gradi.
Definiamo $N(q, t)$ come il numero medio di nodi che possiedono in-degree pari a $q$ al tempo $t$.
La probabilità che un nodo scelto a caso al tempo $t$ abbia in-degree $q$ è data da:
\begin{equation}
    P(q, t) = \frac{N(q, t)}{t}
\end{equation}

L'evoluzione del numero di nodi $N(q, t)$ dal tempo $t$ al tempo $t+1$ è governata dai flussi di probabilità tra gli stati (i diversi valori di $q$).
Dobbiamo considerare tre contributi:
\begin{enumerate}
    \item Guadagno da $q-1$: Un nodo che ha in-degree $q-1$ riceve un nuovo collegamento e il suo grado diventa $q$.
    \item Perdita da $q$: Un nodo che ha già in-degree $q$ riceve un nuovo collegamento e il suo grado diventa $q+1$.
    \item Nuovo nodo: Ad ogni passo temporale viene introdotto un nuovo nodo. Questo nodo nasce con in-degree $q=0$ (poiché i suoi link sono solo uscenti).
\end{enumerate}

\noindent Scriviamo l'equazione di bilancio per $N(q, t+1)$. Per $q > 0$:
\begin{equation}
    N(q, t+1) = N(q, t) + \underbrace{\frac{m[(q-1) + m]}{2mt} N(q-1, t)}_{\text{Guadagno da } q-1} - \underbrace{\frac{m[q + m]}{2mt} N(q, t)}_{\text{Perdita verso } q+1}
\end{equation}
Semplificando i termini $m$ e $t$:
\begin{equation}
    N(q, t+1) = N(q, t) + \frac{q - 1 + m}{2t} N(q-1, t) - \frac{q + m}{2t} N(q, t)
\end{equation}

\noindent Per $q = 0$: il nuovo nodo entra sempre con in-degree 0. Inoltre, nodi con $q=0$ possono ricevere un link e passare a $q=1$.
\begin{equation}
    N(0, t+1) = N(0, t) + 1 - \frac{m}{2t} N(0, t)
\end{equation}
Il termine $+1$ rappresenta il nuovo nodo aggiunto al sistema al tempo $t+1$.

\subsubsection{Soluzione nello Stato Stazionario}

Cerchiamo una soluzione stazionaria in cui la distribuzione di probabilità $P(q)$ non dipenda dal tempo per $t \to \infty$.
Assumiamo $N(q, t) = t P(q)$ e sostituendo nella Master Equation per $q > 0$:
\begin{equation}
    P_q = \frac{1}{2}(q - 1 + m) P_{q-1} - \frac{1}{2}(q + m) P_q
\end{equation}
Raggruppiamo i termini con $P(q)$ a sinistra e otteniamo la relazione ricorsiva fondamentale:
\begin{equation}
    P_q^\infty  = P_{q-1}^\infty  \frac{q + m - 1}{q + m + 2}
\end{equation}

\noindent Analizziamo ora l'equazione per $q=0$ nel limite stazionario. Si trova la seguente relazione:
\begin{equation}
 P_0^\infty = \frac{2}{m + 2}
\end{equation}

\subsection{Soluzione Esatta e Comportamento Asintotico}

Utilizzando la relazione ricorsiva, possiamo scrivere $P(q)$ come prodotto:
\begin{equation}
    P_q^\infty  = P_0^\infty  \prod_{j=1}^{q} \frac{j + m - 1}{j + m + 2}
\end{equation}
Questa espressione può essere riscritta utilizzando la funzione Gamma di Eulero, $\Gamma(x)$, sfruttando la proprietà $\Gamma(x+1) = x \Gamma(x)$.
La relazione ricorsiva $P(q) / P(q-1)$ suggerisce la forma:
\begin{equation}
    P(q) = \frac{2}{m+2} \frac{\Gamma(q+m) \Gamma(m+3)}{\Gamma(m) \Gamma(q+m+3)}
\end{equation}
Verifichiamo i termini costanti. Sapendo che $P(0) = \frac{2}{m+2}$, l'espressione generale diventa:
\begin{equation}
    P(q) = \frac{\Gamma(m+3)}{\Gamma(m)} \frac{\Gamma(q+m)}{\Gamma(q+m+3)}
\end{equation}
Il pre-fattore $\frac{\Gamma(m+3)}{\Gamma(m)} = (m+2)(m+1)m! / (m-1)! \dots$ è una costante di normalizzazione.

\noindent  Per studiare il comportamento asintotico per grandi $q$ ($q \gg 1$), utilizziamo l'approssimazione di Stirling o la proprietà asintotica del rapporto di funzioni Gamma:
\begin{equation}
    \frac{\Gamma(z+a)}{\Gamma(z+b)} \approx z^{a-b} \quad \text{per } z \to \infty
\end{equation}
Nel nostro caso, $z = q$, $a = m$, $b = m+3$.
\begin{equation}
    \frac{\Gamma(q+m)}{\Gamma(q+m+3)} \approx q^{m - (m+3)} = q^{-3}
\end{equation}
Quindi, la distribuzione del grado entrante segue una legge di potenza:
\begin{equation}
    P(q) \sim q^{-3}
\end{equation}

Poiché per grandi $q$ si ha $k \approx q$ (dato che $k = q+m$ e $m$ è costante), la distribuzione del grado totale $P(k)$ avrà lo stesso esponente asintotico:
\begin{equation}
    P(k) \sim k^{-3}
\end{equation}
Questo conferma rigorosamente il risultato riportato nell'Equazione \ref{eq: 10.1}, dimostrando che l'esponente $\gamma = 3$ è una proprietà robusta del modello Barabási-Albert.


\newpage
\section{Dinamiche di Diffusione Epidemica}

La modellizzazione matematica in epidemiologia vanta una storia lunga e ricca, che risale agli anni '20 con la teoria di Kermack-McKendrick. L'idea fondamentale, seppur semplice, si rivela estremamente potente: è possibile suddividere la popolazione in diversi \textit{compartimenti} che rappresentano gli stadi della malattia e utilizzare le dimensioni relative di tali compartimenti per modellare l'evoluzione temporale del contagio.

Sebbene nasca in ambito biologico, questa teoria ha trovato una forte sinergia con la \textit{Network Science}. Oggi, i modelli epidemici sono utilizzati per comprendere non solo la diffusione di patogeni biologici, ma anche la propagazione di idee, virus informatici o "meme" sui social network.
Le domande centrali a cui questi modelli cercano di rispondere sono:
\begin{itemize}
    \item Come si diffonde un'infezione all'interno di una popolazione?
    \item La malattia evolverà in una pandemia globale o si estinguerà localmente?
    \item Qual è la strategia di vaccinazione ottimale?
\end{itemize}

\subsubsection{Dinamica Co-evolutiva}
È fondamentale considerare che esistono due popolazioni interagenti:
\begin{enumerate}
    \item \textbf{Ospiti (Hosts):} Gli individui che contraggono il virus.
    \item \textbf{Patogeni (Virus):} Gli agenti infettivi che evolvono per eludere il sistema immunitario dell'ospite.
\end{enumerate}
Questo sistema è intrinsecamente co-evolutivo. Ad esempio, nel caso dell'influenza stagionale, i vaccini devono essere aggiornati regolarmente per adattarsi ai ceppi dominanti che emergono dalla pressione selettiva.

\hfill

\noindent Iniziamo analizzando i modelli in approssimazione di \textbf{well-mixed population} (popolazione ben mescolata), dove si assume che ogni individuo possa interagire con chiunque altro con la stessa probabilità.

\subsection{Il Modello SI (Susceptible-Infected)}
Il modello più semplice divide la popolazione $N$ in due classi:
\begin{itemize}
    \item \textbf{S (Susceptible):} Individui sani ma suscettibili all'infezione.
    \item \textbf{I (Infected):} Individui infetti in grado di trasmettere la malattia.
\end{itemize}
Sia $\beta$ il tasso di infezione, che dipende dalla contagiosità del patogeno e dalla frequenza dei contatti sociali. La dinamica segue la \textbf{Legge di Azione di Massa}: la variazione degli infetti è proporzionale al prodotto tra il numero di suscettibili e la frazione di infetti presenti.
Le equazioni differenziali che governano il sistema sono:
\begin{align}
    \frac{dI(t)}{dt} &= \beta S(t) \frac{I(t)}{N} \\
    \frac{dS(t)}{dt} &= -\beta S(t) \frac{I(t)}{N}
\end{align}
Con la condizione di normalizzazione $S(t) + I(t) = N$.
Passando alle variabili normalizzate $s(t) = S(t)/N$ e $i(t) = I(t)/N$, otteniamo:
\begin{equation}
    \frac{di}{dt} = \beta s i = \beta (1-i) i
\end{equation}
La soluzione di questa equazione differenziale è la \textbf{funzione logistica}:
\begin{equation}
    i(t) = \frac{i(0) e^{\beta t}}{1 - i(0) + i(0)e^{\beta t}}
\end{equation}
Questo modello prevede che, per $t \to \infty$, l'intera popolazione venga infettata ($i \to 1$). È un modello adatto per malattie incurabili senza recupero, ma limitato per la maggior parte dei casi reali.

\begin{figure}[h]
\centering
\begin{tikzpicture}
    \begin{axis}[
        width=8cm, height=6cm,
        xlabel={$t$}, 
        ymin=0, ymax=1.1, xmin=0, xmax=10,
        axis lines=left,
        grid=major,
        legend pos=outer north east
       ]
    % Curva degli Infetti i(t) - Crescita Logistica
    \addplot[domain=0:10, samples=100, thick, colorC] {1/(1 + 99*exp(-1*x))};
    \addlegendentry{$i(t)$}
    
    % Curva dei Suscettibili s(t) - Decadimento (s = 1 - i)
    \addplot[domain=0:10, samples=100, thick, colorE, dashed] {1 - (1/(1 + 99*exp(-1*x)))};
    \addlegendentry{$s(t)$}
    \end{axis}
\end{tikzpicture}
\caption{Andamento complementare di suscettibili $s(t)$ e infetti $i(t)$ nel modello SI.}
\end{figure}

\subsection{Il Modello SIR (Susceptible-Infected-Recovered)}
Per descrivere malattie da cui si può guarire (o morire), acquisendo immunità, si introduce il compartimento \textbf{R (Recovered)}.
\begin{itemize}
    \item $S \to I$: transizione con tasso $\beta$ (infezione).
    \item $I \to R$: transizione con tasso $\mu$ (guarigione/rimozione).
\end{itemize}
Il sistema di equazioni differenziali è:
\begin{align}
    \frac{ds}{dt} &= -\beta s i \\
    \frac{di}{dt} &= \beta s i - \mu i \\
    \frac{dr}{dt} &= \mu i
\end{align}
Dove $\mu$ è il tasso di recupero, inverso del tempo medio di infezione $\tau \sim 1/\mu$. La probabilità di rimanere infetti per un tempo $\tau$ decade esponenzialmente come $P(\tau) \sim e^{-\mu \tau}$.

\subsubsection{Derivazione della Dimensione Finale dell'Epidemia}
Dividendo la prima equazione per la terza, eliminiamo la dipendenza temporale:
\begin{equation}
    \frac{ds}{dr} = \frac{-\beta s i}{\mu i} = -\frac{\beta}{\mu} s
\end{equation}
Definiamo il \textbf{Basic Reproduction Number} come $R_0 = \frac{\beta}{\mu}$. Integrando l'equazione differenziale tra lo stato iniziale (tempo 0) e un tempo generico $t$:
\begin{equation}
    \int_{s(0)}^{s(t)} \frac{ds'}{s'} = -R_0 \int_{r(0)}^{r(t)} dr' \implies \ln\left(\frac{s(t)}{s(0)}\right) = -R_0 (r(t) - r(0))
\end{equation}
Assumendo che all'inizio nessuno sia guarito ($r(0)=0$) e quasi tutti siano suscettibili ($s(0) \approx 1$), otteniamo:
\begin{equation}
    s(t) = e^{-R_0 r(t)}
\end{equation}
Nel limite asintotico $t \to \infty$, il numero di infetti si annulla ($i_\infty = 0$), quindi vale la conservazione $s_\infty + r_\infty = 1$. Sostituendo $s_\infty = 1 - r_\infty$:
\begin{equation}
    1 - r_\infty = e^{-R_0 r_\infty}
\end{equation}
Otteniamo così un'equazione trascendente per la dimensione totale dell'epidemia $r_\infty$:
\begin{equation}
    r_\infty = 1 - e^{-R_0 r_\infty}
\end{equation}

\begin{figure}[h]
\centering
\begin{tikzpicture}
    \begin{axis}[
        width=8cm, height=6cm,
        xlabel={$r_\infty$},
        xmin=0, xmax=1, ymin=0, ymax=1,
        legend pos=south east,
        axis lines=middle
    ]
    \addplot[domain=0:1, dashed, black] {x};
    \addplot[domain=0:1, thick, colorE] {1 - exp(-2.5*x)}; % R0 > 1
    \addplot[domain=0:1, thick, colorB] {1 - exp(-0.8*x)}; % R0 < 1
    
    \node[colorE] at (axis cs: 0.7, 0.9) {$R_0 > 1$};
    \node[colorB] at (axis cs: 0.3, 0.1) {$R_0 < 1$};
    \end{axis}
\end{tikzpicture}
\caption{Soluzione grafica per $r_\infty$. Per $R_0 > 1$  esiste una soluzione non banale $r_\infty \neq 0$. Per $R_0 < 1$ l'unica soluzione è l'assenza di epidemia ($r_\infty = 0$).}
\end{figure}

Il parametro $R_0$ determina la soglia critica:
\begin{itemize}
    \item Se $R_0 < 1$: L'epidemia non decolla (soluzione $r_\infty = 0$).
    \item Se $R_0 > 1$: Esiste una soluzione positiva $r_\infty > 0$, si verifica un'epidemia macroscopica.
\end{itemize}

\newpage
\subsection{Epidemie su Reti e Percolazione}

Le assunzioni di "well-mixed population" sono spesso eccessivamente semplicistiche. Nella realtà, la struttura dei contatti (la rete sociale, i trasporti aerei, ecc.) gioca un ruolo cruciale.
Passiamo quindi all'analisi delle epidemie su \textbf{Reti Complesse}.

\subsection{Probabilità di Trasmissione e Percolazione di Legame}
Consideriamo un modello in cui ogni individuo infetto rimane tale per un tempo fisso $\tau$. La probabilità di trasmettere l'infezione a un vicino connesso è data da $\beta$ per unità di tempo.
La probabilità totale che l'infezione passi attraverso un arco (link) durante il periodo di infezione $\tau$ è definita come \textbf{trasmissibilità} $\phi$ (o $T$):
\begin{equation}
    \phi = 1 - e^{-\beta \tau}
\end{equation}
Questo ci permette di mappare il problema dinamico dell'epidemia su un problema statico di \textbf{Percolazione di Legame} (\textit{Bond Percolation}).
Ogni arco della rete viene dichiarato "occupato" (ovvero capace di trasmettere l'infezione) con probabilità $\phi$.
\begin{itemize}
    \item Se esiste una \textbf{Componente Gigante} (Giant Component) formata dagli archi occupati, l'epidemia può diffondersi a una frazione macroscopica della popolazione.
    \item La soglia epidemica corrisponde esattamente alla soglia di percolazione $\phi_c$.
\end{itemize}


\subsection{Derivazione della Soglia Critica $\phi_c$}
Vogliamo calcolare la probabilità che un nodo \textbf{non} appartenga alla Componente Gigante (GC) dell'infezione.
Definiamo $u$ come la probabilità che seguendo un arco scelto a caso in una direzione, \textbf{non}  si raggiunga la Componente Gigante. 

Consideriamo un nodo $i$ raggiunto tramite un arco. Affinché questo percorso non conduca alla GC, deve verificarsi una delle seguenti condizioni esclusive:
\begin{enumerate}
    \item L'arco che stiamo percorrendo non è attivo (l'infezione non passa). Questo avviene con probabilità $1 - \phi$.
    \item L'arco è attivo (probabilità $\phi$), ma il nodo $i$ a cui arriviamo non è connesso alla GC tramite nessuno dei suoi \textit{altri} archi uscenti.
\end{enumerate}

\noindent Se il nodo $i$ ha grado totale $k+1$, significa che ha altri $k$ archi oltre a quello di arrivo. La probabilità che nessuno di questi $k$ rami porti alla GC è $u^k$ (assumendo indipendenza locale, valida per reti senza troppi cicli corti).

Dobbiamo mediare questa probabilità sulla distribuzione dei gradi dei nodi che incontriamo navigando la rete. Questa non è la distribuzione di grado standard $P(k)$, ma la \textbf{Excess Degree Distribution} (distribuzione del grado in eccesso), denotata con $q_k$.
La probabilità di arrivare a un nodo con grado $k+1$ è proporzionale a $(k+1)P(k+1)$. Normalizzando:
\begin{equation}
    q_k = \frac{(k+1)P(k+1)}{\langle k \rangle}
\end{equation}
Dove $q_k$ rappresenta la probabilità che un nodo raggiunto abbia $k$ collegamenti ulteriori.

Possiamo ora scrivere l'equazione ricorsiva per $u$:
\begin{equation}
    u = \underbrace{(1 - \phi)}_{\text{Arco inattivo}} + \underbrace{\phi \sum_{k=0}^{\infty} q_k u^k}_{\text{Arco attivo } \cdot \text{ rami morti}}
\end{equation}
Introduciamo la \textbf{Funzione Generatrice} $g_1(u)$ per la distribuzione di eccesso:
\begin{equation}
    g_1(u) = \sum_{k=0}^{\infty} q_k u^k
\end{equation}
L'equazione di auto-consistenza diventa quindi:
\begin{equation} \label{eq:self_consistency}
    u = 1 - \phi + \phi g_1(u)
\end{equation}

\subsubsection{Soluzione Grafica e Transizione di Fase}

Cerchiamo le soluzioni dell'equazione (\ref{eq:self_consistency}) nell'intervallo $u \in [0, 1]$.




\begin{itemize}
    \item Una soluzione banale è sempre $u=1$. Infatti, poiché $\sum q_k = 1$, si ha $g_1(1)~=~1$, quindi $1 = 1 - \phi + \phi(1)$ è verificata. Questo corrisponde all'assenza di una Componente Gigante (l'epidemia si estingue).
    \item Una soluzione fisica $u < 1$ (che implica l'esistenza della GC) emerge solo se la curva interseca la bisettrice $y=u$ con una pendenza sufficiente.
\end{itemize}

La transizione di fase (il punto critico) avviene quando la derivata del membro di destra rispetto a $u$, calcolata in $u=1$, è esattamente uguale a 1 (tangenza):
\begin{equation}
    \frac{d}{du} \Big( 1 - \phi + \phi g_1(u) \Big) \Big|_{u=1} = 1
\end{equation}
Calcolando la derivata:
\begin{equation}
    0 - 0 + \phi \cdot g_1'(u) \Big|_{u=1} = 1 \implies \phi_c g_1'(1) = 1
\end{equation}
La soglia critica è dunque:
\begin{equation}
    \phi_c = \frac{1}{g_1'(1)}
\end{equation}

\subsubsection{Calcolo Esplicito dei Momenti}

Dobbiamo calcolare il valore di $g_1'(1)$. Dalla definizione di funzione generatrice:
\begin{equation}
    g_1'(u) = \frac{d}{du} \sum_{k=0}^{\infty} q_k u^k = \sum_{k=1}^{\infty} k q_k u^{k-1}
\end{equation}
Valutando in $u=1$:
\begin{equation}
    g_1'(1) = \sum_{k=0}^{\infty} k q_k
\end{equation}
Questo rappresenta il valor medio della distribuzione di eccesso. Sostituiamo ora la definizione di $q_k = \frac{(k+1)P(k+1)}{\langle k \rangle}$:
\begin{equation}
    g_1'(1) = \sum_{k=0}^{\infty} k \frac{(k+1)P(k+1)}{\langle k \rangle}
\end{equation}
Per risolvere questa sommatoria, effettuiamo un cambio di variabile. Poniamo $j = k+1$, da cui $k = j-1$. Quando $k=0$, $j=1$.
\begin{equation}
    g_1'(1) = \frac{1}{\langle k \rangle} \sum_{j=1}^{\infty} (j-1) j P(j)
\end{equation}
Espandiamo il prodotto $(j-1)j = j^2 - j$:
\begin{equation}
    g_1'(1) = \frac{1}{\langle k \rangle} \left( \sum_{j} j^2 P(j) - \sum_{j} j P(j) \right)
\end{equation}
Riconosciamo i momenti della distribuzione originale $P(k)$: $\sum j^2 P(j) = \langle k^2 \rangle$ e $\sum j P(j) = \langle k \rangle$.
Quindi:
\begin{equation}
    g_1'(1) = \frac{\langle k^2 \rangle - \langle k \rangle}{\langle k \rangle}
\end{equation}

\subsubsection{Risultato Finale per la Soglia Epidemica}

Sostituendo l'espressione di $g_1'(1)$ nella formula della soglia critica $\phi_c = 1/g_1'(1)$, otteniamo il risultato fondamentale:

\begin{tcolorbox}[colback=colorC!15, colframe=colorC!70!colorB,  title=\textbf{Soglia Epidemica}]
\begin{equation} \label{eq:threshold}
    \phi_c = \frac{\langle k \rangle}{\langle k^2 \rangle - \langle k \rangle}
\end{equation}
\end{tcolorbox}


Questa equazione lega la dinamica epidemica (rappresentata da $\phi_c$) esclusivamente alle proprietà topologiche della rete (i momenti $\langle k \rangle$ e $\langle k^2 \rangle$).

\subsection{Implicazioni per la Robustezza delle Reti}

Analizziamo il comportamento della soglia (\ref{eq:threshold}) per diverse topologie di rete.

\subsubsection{Caso 1: Reti Omogenee (Erdős-Rényi)}
In una rete random con distribuzione di Poisson, la varianza $\sigma^2$ è uguale alla media $\langle k \rangle$.
Ricordando che $\sigma^2 = \langle k^2 \rangle - \langle k \rangle^2$, possiamo scrivere:
\begin{equation}
    \langle k^2 \rangle = \langle k \rangle^2 + \sigma^2 = \langle k \rangle^2 + \langle k \rangle
\end{equation}
Sostituendo nel denominatore della soglia critica:
\begin{equation}
    \phi_c = \frac{\langle k \rangle}{(\langle k \rangle^2 + \langle k \rangle) - \langle k \rangle} = \frac{\langle k \rangle}{\langle k \rangle^2} = \frac{1}{\langle k \rangle}
\end{equation}
Se la rete ha grado medio $\langle k \rangle = 4$, allora $\phi_c = 1/4 = 0.25$.
Ciò implica una notevole \textbf{robustezza}: affinché l'epidemia si propaghi (o affinché la rete rimanga connessa sotto attacco casuale), la trasmissibilità deve essere superiore al 25\%. Inversamente, possiamo rimuovere fino al 75\% dei nodi prima che la componente gigante venga distrutta.

\subsubsection{Caso 2: Reti Scale-Free (Barabási-Albert)}
Le reti reali (Internet, WWW, reti sociali) seguono una legge di potenza $P(k) \sim k^{-\gamma}$ con $2 < \gamma \le 3$.
Una proprietà peculiare di queste distribuzioni è che, nel limite di rete infinita ($N \to \infty$), il momento secondo diverge:
\begin{equation}
    \langle k^2 \rangle \to \infty
\end{equation}
Se il denominatore nella (\ref{eq:threshold}) tende a infinito, la soglia critica tende a zero:
\begin{equation}
    \phi_c = \frac{\langle k \rangle}{\infty} \longrightarrow 0
\end{equation}

Questo risultato dimostra l'estrema \textbf{fragilità} delle reti scale-free rispetto alla diffusione virale.
\begin{itemize}
    \item Poiché $\phi_c = 0$, qualsiasi virus anche con trasmissibilità infinitesimale è in grado di diffondersi e diventare pandemico.
    \item Non esiste una soglia di sicurezza naturale.
    \item Un numero infinitesimale di nodi iniziali è sufficiente a sostenere il contagio 
\end{itemize}


\section{Disclaimer}

Questi appunti non coprono la seconda parte del corso.

%\include{l_12}
%\include{l_13}
%\include{l_14}
%\include{l_15}
%\include{l_16}
%\include{l_17}

\end{document}
